%% GENERATED FILE. Do not edit, with one exception.
%% Written by .bibaudit/build_audit.py. Regenerate instead of correcting by hand:
%%     .bibaudit/build_audit.py && latexmk -pdf source_audit.tex
%%
%% The exception is Part 1's last column, which is the author's review tick. Change a
%% $\square$ to a $\blacksquare$ to record that an entry has been checked. The generator
%% reads those ticks back before it overwrites this file and reprints them, and mirrors
%% them to .bibaudit/reviewed.txt so they survive the file being deleted. Everything else
%% here is derived, so a hand edit to it is lost on the next run.
%%
%% Every reference below is printed by biblatex from bibliography.bib. Every printed
%% citation number is read from the corresponding .bbl. Every citing sentence is copied
%% from the thesis source. None of it is retyped here, so this document cannot disagree
%% with the thesis about what a reference says or which number it carries.
\documentclass[10pt,a4paper]{article}

\usepackage[T1]{fontenc}
\usepackage[utf8]{inputenc}
\usepackage{lmodern}
\usepackage[a4paper,margin=22mm,bottom=24mm]{geometry}
\usepackage{microtype}
\usepackage{longtable}
\usepackage{booktabs}
\usepackage{array}
\usepackage{amsmath}
\usepackage{amssymb}
\usepackage{enumitem}
\usepackage{csquotes}
\usepackage{xcolor}
\usepackage{fancyhdr}

%% The thesis's own biblatex configuration, so \fullcite formats a reference with the same
%% code that formats the thesis bibliography.
\usepackage[style=numeric,bibstyle=numeric,backend=biber,natbib=true,
            sorting=none,sortcites=true,maxbibnames=99,giveninits=true,uniquename=init]{biblatex}
\addbibresource{bibliography.bib}
\DeclareNameAlias{sortname}{family-given}

\usepackage[colorlinks=true,linkcolor=black,citecolor=black,urlcolor=blue!55!black]{hyperref}

%% The generated macro files the thesis reads. A citing sentence that quotes a measured
%% value therefore prints the value, not the macro name.
% Generated by bench/plot_figures.py -- do not edit by hand.
\newcommand{\ReprLoad}{28}
\newcommand{\ReprLoadAch}{28.0}
\newcommand{\ReprLoadShortfall}{0.1}
\newcommand{\SatKnee}{334}
\newcommand{\TrackingLoad}{336}
\newcommand{\NRepeats}{20}
\newcommand{\NLoadLevels}{14}
\newcommand{\NVariants}{4}
\newcommand{\NRuns}{1120}
\newcommand{\PNineNineSpread}{14}
\newcommand{\PFiftySpread}{2.9}
\newcommand{\PFiftyRange}{3}
\newcommand{\PNineNineRange}{460}
\newcommand{\PNineNineMaxCI}{2540}
\newcommand{\MaxOfferedLoad}{784}
\newcommand{\NTests}{10}
\newcommand{\MinAttainableP}{0.0000019}
\newcommand{\CorrectedFloor}{0.000019}
\newcommand{\NSignificant}{2}
\newcommand{\SmallestHolm}{0.00010}
\newcommand{\NExclZero}{4}
\newcommand{\NTailExclZero}{1}
\newcommand{\TailExclZeroRows}{p99 latency under system}
\newcommand{\RobustExclZeroRows}{p50 latency, Peak RSS and CPU utilisation under system}
\newcommand{\PFiftyDiffAbsMax}{1.6}
\newcommand{\PFiftyDiffCIMax}{3}
\newcommand{\RssLangDiffLo}{1.3}
\newcommand{\RssLangDiffHi}{1.7}
\newcommand{\RssLangCILo}{-0.5}
\newcommand{\RssLangCIHi}{1.9}
\newcommand{\CpuDiffAbsMax}{3.4}
\newcommand{\CrossAllocConsistent}{p50 latency, p99 latency, p99.9 latency, Peak RSS and CPU utilisation}
\newcommand{\CrossAllocReversed}{none}
\newcommand{\NCrossAllocReversed}{0}
\newcommand{\RssAllocDiffLo}{55.9}
\newcommand{\RssAllocDiffHi}{56.5}
\newcommand{\BudgetRepeats}{15}
\newcommand{\BudgetCoreCpp}{34.4}
\newcommand{\BudgetCoreRust}{21.3}
\newcommand{\BudgetSvcCpp}{13.52}
\newcommand{\BudgetSvcRust}{12.63}
\newcommand{\BudgetHandoffCpp}{13.49}
\newcommand{\BudgetHandoffRust}{12.61}
\newcommand{\BudgetCoreShareMax}{0.250}
\newcommand{\BudgetCoreFaster}{Rust}
\newcommand{\BudgetCoreRatio}{1.6}
\newcommand{\BudgetHandoffSlower}{C++}
\newcommand{\BudgetBoundVsCore}{101}
\newcommand{\FfiCrossing}{2.54}
\newcommand{\FfiOpRepeats}{15}
\newcommand{\FfiOpNative}{12.3}
\newcommand{\FfiOpSeam}{12.7}
\newcommand{\FfiOpDelta}{+192}
\newcommand{\FfiOpDeltaLo}{-954}
\newcommand{\FfiOpDeltaHi}{+1341}
\newcommand{\FfiOpP}{0.239}
\newcommand{\FfiOpBoundPct}{10.91}

% Generated by bench/complexity.py -- do not edit by hand.
\newcommand{\CppFuncs}{37}
\newcommand{\RustFuncs}{40}
\newcommand{\CppNloc}{511}
\newcommand{\RustNloc}{608}
\newcommand{\CppCcnTotal}{93}
\newcommand{\RustCcnTotal}{79}
\newcommand{\CppCcnMean}{2.5}
\newcommand{\RustCcnMean}{2.0}
\newcommand{\CppCcnMax}{13}
\newcommand{\RustCcnMax}{6}
\newcommand{\CppCcnMaxFn}{\texttt{worker\_loop}}
\newcommand{\RustCcnMaxFn}{\texttt{add\_drb}}
\newcommand{\CcnTotalDrop}{15}
\newcommand{\NlocRise}{19}

% Generated by bench/artefact_facts.py -- do not edit by hand.
\newcommand{\CppUeCtxBytes}{96}
\newcommand{\RustUeCtxBytes}{88}
\newcommand{\UeCtxGapBytes}{8}
\newcommand{\CppUeCtxStatePad}{3}
\newcommand{\CppUeCtxKeyEnd}{50}
\newcommand{\CppUeCtxKeyPad}{6}
\newcommand{\CppRequestBytes}{160}
\newcommand{\RustRequestBytes}{96}
\newcommand{\SignalingMsgBytes}{68}
\newcommand{\SweepUes}{5\,000}
\newcommand{\UeCtxGapTotalKB}{40}
\newcommand{\CppQueueResidentMB}{1.3}
\newcommand{\RustQueueResidentMB}{0.8}
\newcommand{\CppLineCov}{94.4}
\newcommand{\RustLineCov}{97.9}
\newcommand{\RustRegionCov}{99.1}
\newcommand{\ConfHandOps}{223}
\newcommand{\ConfGenOps}{9\,000}
\newcommand{\FfiSeamLines}{231}
\newcommand{\FfiHeaderLines}{67}
\newcommand{\FfiEntryPoints}{11}
\newcommand{\FfiUnsafeFns}{9}
\newcommand{\DeadlockTsanHits}{3}
\newcommand{\DeadlockTsanRuns}{12}
\newcommand{\FaultClasses}{14}
\newcommand{\FaultClassesPrevented}{7}
\newcommand{\FaultClassesBoundary}{7}
\newcommand{\FaultClassesUndetected}{4}
\newcommand{\LostUpdateExpected}{40\,000}
\newcommand{\LostUpdateRustHandled}{28\,714}
\newcommand{\LostUpdateRustLost}{11\,286}
\newcommand{\LostUpdateCppHandled}{27\,320}
\newcommand{\LostUpdateCppLost}{12\,680}
\newcommand{\CppTestBinaries}{6}
\newcommand{\RustTests}{27}

% Generated by bench/adoption_cost.py -- do not edit by hand.
\newcommand{\AdoptRepeats}{15}
\newcommand{\AdoptJobs}{14}
\newcommand{\BuildLibCppWall}{1.9}
\newcommand{\BuildLibRustWall}{1.2}
\newcommand{\BuildLibCppCpu}{2.6}
\newcommand{\BuildLibRustCpu}{3.9}
\newcommand{\BuildLibWallRatio}{1.7}
\newcommand{\BuildLibCpuRatio}{1.5}
\newcommand{\BuildLibWallFaster}{Rust}
\newcommand{\BuildLibCpuCheaper}{C++}
\newcommand{\BuildAllCppWall}{5.7}
\newcommand{\BuildAllRustWall}{8.1}
\newcommand{\BuildAllCppCpu}{35.9}
\newcommand{\BuildAllRustCpu}{64.8}
\newcommand{\BuildAllCpuRatio}{1.8}
\newcommand{\BuildAllFirstPartyCpu}{29.6}
\newcommand{\BuildAllFirstPartyWall}{3.8}
\newcommand{\BuildDepShareCpu}{54}
\newcommand{\BuildAllFirstPartyCpuRatio}{0.8}
\newcommand{\BuildIncCppTu}{1.2}
\newcommand{\BuildIncCppHeader}{5.8}
\newcommand{\BuildIncRust}{3.6}
\newcommand{\BuildIncTuRatio}{3.1}
\newcommand{\BuildIncHeaderRatio}{1.6}
\newcommand{\BuildConfigure}{0.91}
\newcommand{\FootTarget}{conformance\_runner}
\newcommand{\FootCppShipped}{106}
\newcommand{\FootRustShipped}{723}
\newcommand{\FootShippedRatio}{6.8}
\newcommand{\FootCppStatic}{1577}
\newcommand{\FootRustStatic}{723}
\newcommand{\FootStaticRatio}{0.5}
\newcommand{\FootCppDynamic}{106}
\newcommand{\FootRustDynamic}{269}
\newcommand{\FootDynamicRatio}{2.5}
\newcommand{\FootRustStdShare}{63}
\newcommand{\FootCppStdShare}{93}
\newcommand{\FootSharedCpp}{2531}
\newcommand{\FootSharedRust}{5295}
\newcommand{\FootStaticSmaller}{Rust}
\newcommand{\FootDynamicSmaller}{C++}
\newcommand{\FootStaticRatioAbs}{2.2}
\newcommand{\FootCppLocaleShare}{47}
\newcommand{\DepsCore}{0}
\newcommand{\DepsJemalloc}{8}
\newcommand{\DepsHarness}{17}
\newcommand{\DepsUnsafeJemalloc}{623}
\newcommand{\DepsUnsafeHarness}{1064}
\newcommand{\CoreUnsafe}{0}
\newcommand{\DepsJemallocShipping}{5}
\newcommand{\DepsJemallocBuildOnly}{3}
\newcommand{\DepsHarnessShipping}{16}
\newcommand{\DepsHarnessBuildOnly}{1}

% Generated by bench/baremetal_facts.py -- do not edit by hand.
\newcommand{\BmCores}{24}
\newcommand{\BmCounterRepeats}{10}
\newcommand{\BmCpu}{AMD EPYC 7443P}
\newcommand{\BmCpuDiffHi}{7.6}
\newcommand{\BmCpuDiffLo}{0.3}
\newcommand{\BmCpuLower}{Rust}
\newcommand{\BmCpuOneSign}{yes}
\newcommand{\BmCpuSig}{16}
\newcommand{\BmCpuTested}{16}
\newcommand{\BmCycleDelta}{15}
\newcommand{\BmCycleDeltaDir}{fewer}
\newcommand{\BmCycleDeltaP}{0.0117}
\newcommand{\BmDistro}{Ubuntu 24.04.4 LTS}
\newcommand{\BmFarTailCppHi}{228}
\newcommand{\BmFarTailCppLo}{147}
\newcommand{\BmFarTailRustHi}{115}
\newcommand{\BmFarTailRustLo}{88}
\newcommand{\BmFarTailSepHi}{96}
\newcommand{\BmFarTailSepLoad}{32}
\newcommand{\BmFarTailTopLoad}{128}
\newcommand{\BmGovernor}{performance}
\newcommand{\BmGpp}{13.3.0}
\newcommand{\BmInstrDelta}{11}
\newcommand{\BmInstrDeltaDir}{fewer}
\newcommand{\BmInstrDeltaP}{0.002}
\newcommand{\BmIsolcpus}{none}
\newcommand{\BmJeActiveCpp}{1.43}
\newcommand{\BmJeActiveDelta}{-0.64}
\newcommand{\BmJeActiveRust}{0.79}
\newcommand{\BmJeAllocatedCpp}{0.78}
\newcommand{\BmJeAllocatedDelta}{-0.16}
\newcommand{\BmJeAllocatedRust}{0.62}
\newcommand{\BmJeLiveLower}{Rust}
\newcommand{\BmJeMappedCpp}{42.82}
\newcommand{\BmJeMappedDelta}{+2.01}
\newcommand{\BmJeMappedRust}{44.83}
\newcommand{\BmJeMetaDrift}{0.5}
\newcommand{\BmJeMetadataCpp}{4.19}
\newcommand{\BmJeMetadataDelta}{+2.07}
\newcommand{\BmJeMetadataRust}{6.26}
\newcommand{\BmJeOutsideCpp}{4.0}
\newcommand{\BmJeOutsideDelta}{+3.8}
\newcommand{\BmJeOutsideRust}{7.9}
\newcommand{\BmJeResidentCpp}{11.02}
\newcommand{\BmJeResidentDelta}{-1.92}
\newcommand{\BmJeResidentRust}{9.10}
\newcommand{\BmJeRetainedCpp}{27.72}
\newcommand{\BmJeRetainedDelta}{+3.95}
\newcommand{\BmJeRetainedRust}{31.67}
\newcommand{\BmJeRssCpp}{15.08}
\newcommand{\BmJeRssDelta}{+1.74}
\newcommand{\BmJeRssRust}{16.82}
\newcommand{\BmJeRuns}{240}
\newcommand{\BmKernel}{6.17.0-23-generic}
\newcommand{\BmKneeHi}{512}
\newcommand{\BmKneeLo}{448}
\newcommand{\BmLoadHi}{128}
\newcommand{\BmLoadLevels}{8}
\newcommand{\BmLoadLo}{8}
\newcommand{\BmMetastableCv}{37}
\newcommand{\BmMetastableLoad}{512}
\newcommand{\BmMissDelta}{29}
\newcommand{\BmMissDeltaDir}{more}
\newcommand{\BmMissDeltaP}{0.002}
\newcommand{\BmNohzFull}{none}
\newcommand{\BmNumaNodes}{1}
\newcommand{\BmOverloadCv}{0.4}
\newcommand{\BmOverloadFaster}{C++}
\newcommand{\BmOverloadLoad}{640}
\newcommand{\BmOverloadThrCpp}{158}
\newcommand{\BmOverloadThrRatio}{1.64}
\newcommand{\BmOverloadThrRust}{96}
\newcommand{\BmPFiftyDiffHi}{1.1}
\newcommand{\BmPFiftyDiffLo}{0.5}
\newcommand{\BmPFiftyFaster}{Rust}
\newcommand{\BmPFiftyMedian}{74}
\newcommand{\BmPFiftySig}{14}
\newcommand{\BmPFiftyTested}{16}
\newcommand{\BmPNineNineCvMax}{2.3}
\newcommand{\BmPNineNineCvMed}{1.4}
\newcommand{\BmPNineNineMaxCI}{2.8}
\newcommand{\BmPNineNineRange}{3.6}
\newcommand{\BmPinned}{2-17}
\newcommand{\BmPoolSamplesMax}{5\,761\,198}
\newcommand{\BmPoolSamplesMin}{360\,000}
\newcommand{\BmProducers}{8}
\newcommand{\BmRepeats}{15}
\newcommand{\BmRssAllocHi}{3.2}
\newcommand{\BmRssAllocLo}{1.9}
\newcommand{\BmRssLangHi}{1.8}
\newcommand{\BmRssLangLo}{0.7}
\newcommand{\BmRssLargerEffect}{allocator}
\newcommand{\BmRssSig}{16}
\newcommand{\BmRssSpreadMax}{1.7}
\newcommand{\BmRssTested}{16}
\newcommand{\BmRuns}{480}
\newcommand{\BmRustc}{1.98}
\newcommand{\BmShards}{8}
\newcommand{\BmSmt}{off}
\newcommand{\BmSockets}{1}
\newcommand{\BmStealCores}{16}
\newcommand{\BmStealIntervals}{60}
\newcommand{\BmStealMs}{0.0}
\newcommand{\BmStealWindow}{61}
\newcommand{\BmStealWithSteal}{0}
\newcommand{\BmTrackingMax}{448}
\newcommand{\BmVirt}{none}

% Generated by bench/guest_facts.py -- do not edit by hand.
\newcommand{\GuestCores}{16}
\newcommand{\GuestCpu}{Intel(R) Xeon(R) CPU E5-2680 v4}
\newcommand{\GuestDeadlineLoad}{56}
\newcommand{\GuestDistro}{Rocky Linux 8.10 (Green Obsidian)}
\newcommand{\GuestIsolcpus}{none}
\newcommand{\GuestKernel}{7.0.2-1.el8.elrepo.x86_64}
\newcommand{\GuestLoadLevels}{14}
\newcommand{\GuestOverloadThrRatio}{1.8}
\newcommand{\GuestPFifty}{80}
\newcommand{\GuestPNineNine}{4518}
\newcommand{\GuestPNineNineMaxCI}{1480}
\newcommand{\GuestPNineNineRange}{293}
\newcommand{\GuestProducers}{7}
\newcommand{\GuestRepeats}{20}
\newcommand{\GuestRssAllocHi}{34.3}
\newcommand{\GuestRssAllocLo}{32.5}
\newcommand{\GuestRssAllocOverLang}{13}
\newcommand{\GuestRssLangHi}{3.5}
\newcommand{\GuestRssLangLo}{1.8}
\newcommand{\GuestRssLowRunOffset}{31}
\newcommand{\GuestRuns}{1120}
\newcommand{\GuestSatKnee}{336}
\newcommand{\GuestShards}{7}
\newcommand{\GuestSteal}{0}
\newcommand{\GuestThp}{always}
\newcommand{\GuestTrackingLoad}{336}
\newcommand{\GuestVirt}{kvm}


%% A citation inside a quoted sentence prints as the number the thesis prints, and a
%% cross-reference prints as its label, since this document has no such target of its own.
\newcommand{\citemark}[1]{\textbf{[#1]}}
\newcommand{\refmark}[1]{\textsl{#1}}

\newcommand{\entryrule}{\par\medskip\noindent\textcolor{black!35}{\rule{\linewidth}{0.6pt}}\par\nobreak\smallskip}
\newcommand{\auditentry}[3]{%
  \entryrule
  \phantomsection\addcontentsline{toc}{subsection}{[#1] #2}%
  \noindent{\large\textbf{[#1]}}\quad\texttt{\textbf{#2}}\hfill{\small #3}\par\nobreak\smallskip}
\newcommand{\afield}[2]{\par\noindent\textbf{#1.}\enspace #2\par}
\newenvironment{saidquote}%
  {\begin{quote}\small\itshape}%
  {\end{quote}}
\newenvironment{citedsentence}%
  {\begin{quote}\small}%
  {\end{quote}}

\setlength{\parindent}{0pt}
\setlength{\parskip}{0.45em}

%% Quoted passages carry long identifiers and URLs that cannot be hyphenated, so give the
%% line breaker room rather than leaving overfull boxes in a document meant to be read.
\sloppy
\emergencystretch=3em

\pagestyle{fancy}
\fancyhf{}
\fancyhead[L]{\small Citation audit}
\fancyhead[R]{\small\thepage}
\renewcommand{\headrulewidth}{0.4pt}

\title{Citation audit\\[2pt]\large Every cited work, the sentence citing it, and where the
claim sits in the source}
\author{Ha Vien}

\date{Generated 2026-09-08}

\begin{document}
\maketitle


\section*{How to read this document}
\addcontentsline{toc}{section}{How to read this document}

This document lists every work in \texttt{bibliography.bib}. It is ordered by the number the
work carries in the thesis's reference list, so entry \textbf{[17]} here is reference
\textbf{[17]} there and the two can be read side by side. It covers
\texttt{Ha\_Vien\_Mastersthesis\_2026} alone, which is the submission document and the only one
that cites every entry, so a number here is always that document's number and never needs
translating.

Each entry answers three questions in order. What is the source. Which sentences cite it, printed
as they stand in the thesis with the citation shown as \citemark{n}. And what the source itself
says, quoted with the page or clause it was read from. When the three agree, the citation is
sound. When the quoted passage does not carry the sentence above it, that is the defect this
document is built to expose.

\textbf{Checking one entry.} Read the citing sentence, then read the quotation under
\emph{Source says}. Ask only whether the second supports the first. The evidence level tells you
how far to trust the quotation without opening the paper yourself.

\begin{description}[leftmargin=2.4em,style=nextline,itemsep=1pt,topsep=2pt]
  \item[Level A] The source was opened and the supporting passage located. A page, clause or
        section is given. Nothing further is needed.
  \item[Level B] A publisher-deposited record, such as a Crossref abstract or an official
        landing page, carries the claim. The full text was not opened.
  \item[Level C] Bibliographic metadata is confirmed and the passage sits behind a paywall.
        Worth opening through the university proxy if the claim is load-bearing.
  \item[No level] The record is a literature-sweep note rather than a located passage, so it
        gives no page or clause. Part~3 lists these together and they are the ones to close
        first.
\end{description}

\textbf{Provenance.} Nothing in this document is typed by hand. References are printed by
biblatex from \texttt{bibliography.bib}. Printed citation numbers are read from the thesis's own
\texttt{.bbl}, which is what biber wrote, and the generator refuses to run when a
\texttt{.bbl} disagrees with its source. Citing sentences are copied from the \texttt{.tex}
sources with the generated measurement macros resolved, so a sentence quoting a measured value
shows the value. Evidence comes from \texttt{docs/Sources\_evidence.md}, with
\texttt{docs/LITERATURE\_NOTES.md} as the fallback, and the origin is named under each entry.
Regenerate with:

\begin{quote}\ttfamily\small
.bibaudit/build\_audit.py \&\& latexmk -pdf source\_audit.tex
\end{quote}

\textbf{One note on punctuation.} The thesis style forbids em dashes and prose semicolons
anywhere a reader sees them, and this document follows that rule in its own prose. Quotations
from cited sources keep their original punctuation, because altering a quotation misquotes it.
Any dash or semicolon below therefore sits inside quoted source material and is deliberate.

\vspace{0.6em}
\textbf{Coverage.} 111 entries, every one of them cited in the thesis.
73 at level~A, 32 at level~B, 6 at level~C, and
0 with no located passage.


\clearpage
\section*{Part 1. Master checklist}
\addcontentsline{toc}{section}{Part 1. Master checklist}

One row per entry, ordered by the thesis's reference number. \textbf{cit} counts the places the thesis cites it.
\textbf{ev} is the evidence level, blank where no passage was located, and
\textbf{status} reads \emph{ok} where the verdict is clean, \emph{note} where the record raises
something, and \emph{sweep} where the only record is a literature-sweep note. Tick the last
column as you check.

\begin{small}
\setlength{\tabcolsep}{3pt}
\begin{longtable}{@{}>{\raggedleft\arraybackslash}p{7mm}p{50mm}p{31mm}p{20mm}>{\raggedleft\arraybackslash}p{7mm}>{\centering\arraybackslash}p{6mm}p{14mm}>{\centering\arraybackslash}p{5mm}@{}}
\toprule
\textbf{\#} & \textbf{key} & \textbf{source} & \textbf{kind} & \textbf{cit} & \textbf{ev} & \textbf{status} &  \\
\midrule
\endfirsthead
\toprule
\textbf{\#} & \textbf{key} & \textbf{source} & \textbf{kind} & \textbf{cit} & \textbf{ev} & \textbf{status} &  \\
\midrule
\endhead
\bottomrule
\endfoot
1 & \texttt{\small 3gpp38401} & {\small 3GPP 2023} & {\small standard} & 6 & A & {\small ok} & $\blacksquare$ \\
2 & \texttt{\small 3gpp38300} & {\small 3GPP 2023} & {\small standard} & 4 & A & {\small ok} & $\blacksquare$ \\
3 & \texttt{\small polese2023oran} & {\small Polese et al. 2023} & {\small peer reviewed} & 4 & A & {\small ok} & $\blacksquare$ \\
4 & \texttt{\small 3gpp38331} & {\small 3GPP 2023} & {\small standard} & 5 & C & {\small note} & $\blacksquare$ \\
5 & \texttt{\small 3gpp38913} & {\small 3GPP 2024} & {\small standard} & 4 & A & {\small note} & $\blacksquare$ \\
6 & \texttt{\small dean2013tail} & {\small Dean 2013} & {\small peer reviewed} & 7 & B & {\small ok} & $\blacksquare$ \\
7 & \texttt{\small barroso2017microseconds} & {\small Barroso et al. 2017} & {\small peer reviewed} & 4 & B & {\small ok} & $\blacksquare$ \\
8 & \texttt{\small msrc2019memory} & {\small Miller 2019} & {\small grey} & 5 & A & {\small ok} & $\blacksquare$ \\
9 & \texttt{\small google2022android} & {\small Vander Stoep 2022} & {\small grey} & 5 & A & {\small ok} & $\blacksquare$ \\
10 & \texttt{\small klabnik2023rust} & {\small Klabnik 2023} & {\small book} & 2 & C & {\small ok} & $\blacksquare$ \\
11 & \texttt{\small jung2018rustbelt} & {\small Jung et al. 2018} & {\small peer reviewed} & 3 & A & {\small ok} & $\blacksquare$ \\
12 & \texttt{\small balasubramanian2017rust} & {\small Balasubramanian et al. 2017} & {\small peer reviewed} & 7 & B & {\small ok} & $\blacksquare$ \\
13 & \texttt{\small itu2023imt2030} & {\small ITU-R 2023} & {\small standard} & 2 & A & {\small ok} & $\blacksquare$ \\
14 & \texttt{\small polese2026airan} & {\small Polese et al. 2026} & {\small peer reviewed} & 1 & B & {\small ok} & $\blacksquare$ \\
15 & \texttt{\small airan2026orchestrator} & {\small AI-RAN Alliance 2026} & {\small grey} & 1 & A & {\small ok} & $\blacksquare$ \\
16 & \texttt{\small vanderstoep2025android} & {\small Vander Stoep 2025} & {\small grey} & 2 & A & {\small ok} & $\blacksquare$ \\
17 & \texttt{\small google2024securebydesign} & {\small Rebert 2024} & {\small grey} & 2 & B & {\small ok} & $\blacksquare$ \\
18 & \texttt{\small watson2025memorysafety} & {\small Watson et al. 2025} & {\small peer reviewed} & 1 & A & {\small note} & $\blacksquare$ \\
19 & \texttt{\small cisa2023memorysafe} & {\small Cybersecurity and\ldots{} 2023} & {\small standard} & 3 & A & {\small note} & $\blacksquare$ \\
20 & \texttt{\small oncd2024building} & {\small Office of the National Cyber\ldots{} 2024} & {\small standard} & 3 & A & {\small ok} & $\blacksquare$ \\
21 & \texttt{\small eu2024cra} & {\small European Parliament and\ldots{} 2024} & {\small regulation} & 1 & A & {\small note} & $\blacksquare$ \\
22 & \texttt{\small agache2020firecracker} & {\small Agache et al. 2020} & {\small peer reviewed} & 1 & A & {\small note} & $\blacksquare$ \\
23 & \texttt{\small delloro2025ran3q} & {\small Dell'Oro Group 2025} & {\small grey} & 1 & A & {\small ok} & $\blacksquare$ \\
24 & \texttt{\small delloro2026ran} & {\small Dell'Oro Group 2026} & {\small grey} & 1 & A & {\small ok} & $\blacksquare$ \\
25 & \texttt{\small 3gpp2024road6g} & {\small Almodovar 2024} & {\small grey} & 1 & A & {\small ok} & $\blacksquare$ \\
26 & \texttt{\small cisa2025badpractices} & {\small Cybersecurity and\ldots{} 2025} & {\small standard} & 1 & A & {\small ok} & $\blacksquare$ \\
27 & \texttt{\small qin2020rust} & {\small Qin et al. 2020} & {\small peer reviewed} & 4 & A & {\small ok} & $\blacksquare$ \\
28 & \texttt{\small meneely2025justuserust} & {\small Meneely et al. 2025} & {\small peer reviewed} & 2 & A & {\small note} & $\square$ \\
29 & \texttt{\small 3gpp22261} & {\small 3GPP 2022} & {\small standard} & 2 & C & {\small note} & $\blacksquare$ \\
30 & \texttt{\small 3gpp38463} & {\small 3GPP 2022} & {\small standard} & 1 & A & {\small ok} & $\blacksquare$ \\
31 & \texttt{\small li2014tales} & {\small Li et al. 2014} & {\small peer reviewed} & 2 & A & {\small ok} & $\square$ \\
32 & \texttt{\small schroeder2006openclosed} & {\small Schroeder et al. 2006} & {\small peer reviewed} & 2 & A & {\small ok} & $\square$ \\
33 & \texttt{\small georges2007rigorous} & {\small Georges et al. 2007} & {\small peer reviewed} & 3 & A & {\small ok} & $\square$ \\
34 & \texttt{\small szekeres2013sok} & {\small Szekeres et al. 2013} & {\small peer reviewed} & 3 & A & {\small ok} & $\square$ \\
35 & \texttt{\small bensalem2005deadlock} & {\small Bensalem 2005} & {\small peer reviewed} & 1 & A & {\small ok} & $\square$ \\
36 & \texttt{\small tsandeadlock} & {\small The Sanitizers Project 2026} & {\small tool record} & 1 & A & {\small ok} & $\square$ \\
37 & \texttt{\small sutter2025profiles} & {\small Sutter 2025} & {\small committee draft} & 1 & A & {\small ok} & $\square$ \\
38 & \texttt{\small baxter2024safecpp} & {\small Baxter 2024} & {\small committee draft} & 1 & A & {\small ok} & $\square$ \\
39 & \texttt{\small vandevoorde2026profiles} & {\small Vandevoorde et al. 2026} & {\small committee draft} & 3 & A & {\small ok} & $\square$ \\
40 & \texttt{\small jung2021safe} & {\small Jung et al. 2021} & {\small peer reviewed} & 3 & A & {\small ok} & $\square$ \\
41 & \texttt{\small jung2020stacked} & {\small Jung et al. 2020} & {\small peer reviewed} & 2 & A & {\small ok} & $\square$ \\
42 & \texttt{\small astrauskas2020unsafe} & {\small Astrauskas et al. 2020} & {\small peer reviewed} & 2 & B & {\small ok} & $\square$ \\
43 & \texttt{\small xu2021rustcve} & {\small Xu et al. 2022} & {\small peer reviewed} & 3 & A & {\small note} & $\square$ \\
44 & \texttt{\small evans2020rust} & {\small Evans et al. 2020} & {\small peer reviewed} & 5 & A & {\small ok} & $\square$ \\
45 & \texttt{\small bugden2022rust} & {\small Bugden 2022} & {\small peer reviewed} & 2 & A & {\small ok} & $\square$ \\
46 & \texttt{\small enberg2019threadpercore} & {\small Enberg et al. 2019} & {\small peer reviewed} & 2 & A & {\small ok} & $\square$ \\
47 & \texttt{\small skoudlil2025asyncrust} & {\small Skoudlil et al. 2025} & {\small peer reviewed} & 2 & A & {\small ok} & $\square$ \\
48 & \texttt{\small rivera2021galeed} & {\small Rivera et al. 2021} & {\small peer reviewed} & 6 & B & {\small ok} & $\square$ \\
49 & \texttt{\small kirth2022pkrusafe} & {\small Kirth et al. 2022} & {\small peer reviewed} & 6 & B & {\small ok} & $\square$ \\
50 & \texttt{\small li2022crosslang} & {\small Li et al. 2022} & {\small peer reviewed} & 3 & C & {\small ok} & $\square$ \\
51 & \texttt{\small cai2026interop} & {\small Cai et al. 2026} & {\small peer reviewed} & 3 & B & {\small ok} & $\square$ \\
52 & \texttt{\small pearce2026rustmc} & {\small Pearce et al. 2026} & {\small peer reviewed} & 2 & A & {\small ok} & $\square$ \\
53 & \texttt{\small karlback2021ran} & {\small Karlback 2021} & {\small thesis} & 7 & A & {\small note} & $\square$ \\
54 & \texttt{\small olsson2023ranmgmt} & {\small Olsson 2023} & {\small thesis} & 6 & A & {\small ok} & $\square$ \\
55 & \texttt{\small olsson2023asyncpackets} & {\small Olsson 2023} & {\small thesis} & 4 & A & {\small ok} & $\square$ \\
56 & \texttt{\small arvidsson2025async} & {\small Arvidsson 2025} & {\small thesis} & 9 & A & {\small ok} & $\square$ \\
57 & \texttt{\small emberg2023asn1} & {\small Emberg 2023} & {\small thesis} & 6 & A & {\small ok} & $\square$ \\
58 & \texttt{\small gjomemo2016mesc} & {\small Gjomemo et al. 2016} & {\small peer reviewed} & 1 & B & {\small ok} & $\square$ \\
59 & \texttt{\small michaelides2026ranlatency} & {\small Michaelides et al. 2026} & {\small peer reviewed} & 1 & A & {\small ok} & $\square$ \\
60 & \texttt{\small natella2016fault} & {\small Natella et al. 2016} & {\small peer reviewed} & 2 & B & {\small ok} & $\square$ \\
61 & \texttt{\small efron1993bootstrap} & {\small Efron 1993} & {\small book} & 1 & B & {\small note} & $\square$ \\
62 & \texttt{\small rfc3393} & {\small Demichelis 2002} & {\small standard} & 2 & A & {\small ok} & $\square$ \\
63 & \texttt{\small holm1979simple} & {\small Holm 1979} & {\small peer reviewed} & 1 & B & {\small ok} & $\square$ \\
64 & \texttt{\small tene\_hdrhistogram} & {\small Tene} & {\small tool record} & 2 & A & {\small ok} & $\square$ \\
65 & \texttt{\small lizard} & {\small Yin 2026} & {\small tool record} & 2 & A & {\small ok} & $\square$ \\
66 & \texttt{\small zhou2025oltp} & {\small Zhou et al. 2025} & {\small peer reviewed} & 2 & A & {\small ok} & $\square$ \\
67 & \texttt{\small clangtidyuseaftermove} & {\small The LLVM Project 2026} & {\small tool record} & 1 & A & {\small ok} & $\square$ \\
68 & \texttt{\small hughes2025gcrust} & {\small Hughes 2025} & {\small peer reviewed} & 1 & B & {\small ok} & $\square$ \\
69 & \texttt{\small cui2024rcanary} & {\small Cui et al. 2024} & {\small peer reviewed} & 1 & A & {\small ok} & $\square$ \\
70 & \texttt{\small dietz2015overflow} & {\small Dietz et al. 2015} & {\small peer reviewed} & 1 & A & {\small ok} & $\square$ \\
71 & \texttt{\small bornholt2021shardstore} & {\small Bornholt et al. 2021} & {\small peer reviewed} & 1 & A & {\small ok} & $\square$ \\
72 & \texttt{\small mccormack2025ffi} & {\small McCormack et al. 2025} & {\small peer reviewed} & 1 & A & {\small ok} & $\square$ \\
73 & \texttt{\small emmerich2019drivers} & {\small Emmerich et al. 2019} & {\small peer reviewed} & 5 & A & {\small ok} & $\square$ \\
74 & \texttt{\small zhang2022runtime} & {\small Zhang et al. 2022} & {\small peer reviewed} & 4 & A & {\small ok} & $\square$ \\
75 & \texttt{\small harris2023morello} & {\small Harris et al. 2023} & {\small peer reviewed} & 4 & A & {\small ok} & $\square$ \\
76 & \texttt{\small zhang2026deadlock} & {\small Zhang et al. 2026} & {\small peer reviewed} & 1 & A & {\small ok} & $\square$ \\
77 & \texttt{\small tang2025converos} & {\small Tang et al. 2025} & {\small peer reviewed} & 1 & A & {\small ok} & $\square$ \\
78 & \texttt{\small lu2008mistakes} & {\small Lu et al. 2008} & {\small peer reviewed} & 1 & A & {\small ok} & $\square$ \\
79 & \texttt{\small sundvisson2026drivers} & {\small Sundvisson 2026} & {\small thesis} & 3 & A & {\small ok} & $\square$ \\
80 & \texttt{\small fulton2021adopting} & {\small Fulton et al. 2021} & {\small peer reviewed} & 4 & A & {\small ok} & $\square$ \\
81 & \texttt{\small crichton2024learning} & {\small Crichton 2024} & {\small peer reviewed} & 2 & A & {\small ok} & $\square$ \\
82 & \texttt{\small li2026rustlinux} & {\small Li et al. 2026} & {\small peer reviewed} & 3 & B & {\small ok} & $\square$ \\
83 & \texttt{\small fan2026android} & {\small Fan et al. 2026} & {\small peer reviewed} & 2 & B & {\small ok} & $\square$ \\
84 & \texttt{\small li2021ipc} & {\small Li et al. 2021} & {\small peer reviewed} & 3 & A & {\small ok} & $\square$ \\
85 & \texttt{\small narayanan2020redleaf} & {\small Narayanan et al. 2020} & {\small peer reviewed} & 2 & A & {\small ok} & $\square$ \\
86 & \texttt{\small suchert2023stamp} & {\small Suchert 2023} & {\small peer reviewed} & 3 & B & {\small ok} & $\square$ \\
87 & \texttt{\small gonzalez2023rustudp} & {\small Gonzalez et al. 2023} & {\small peer reviewed} & 3 & B & {\small ok} & $\square$ \\
88 & \texttt{\small zhuang2026flowprioritize} & {\small Zhuang 2026} & {\small peer reviewed} & 1 & B & {\small ok} & $\square$ \\
89 & \texttt{\small andersson2026stdmemsafety} & {\small Andersson 2026} & {\small thesis} & 1 & A & {\small ok} & $\square$ \\
90 & \texttt{\small okawa2024ruststudy} & {\small Okawa 2024} & {\small peer reviewed} & 1 & B & {\small ok} & $\square$ \\
91 & \texttt{\small pirelli2023safelowlevel} & {\small Pirelli 2023} & {\small peer reviewed} & 2 & B & {\small ok} & $\square$ \\
92 & \texttt{\small qin2024safety} & {\small Qin et al. 2024} & {\small peer reviewed} & 2 & B & {\small ok} & $\square$ \\
93 & \texttt{\small pereira2017energy} & {\small Pereira et al. 2017} & {\small peer reviewed} & 1 & A & {\small ok} & $\square$ \\
94 & \texttt{\small levy2017kernel} & {\small Levy et al. 2017} & {\small peer reviewed} & 1 & B & {\small ok} & $\square$ \\
95 & \texttt{\small mo2023rustshyper} & {\small Mo et al. 2023} & {\small peer reviewed} & 1 & C & {\small ok} & $\square$ \\
96 & \texttt{\small xiao2024rpkt} & {\small Xiao et al. 2025} & {\small peer reviewed} & 1 & C & {\small ok} & $\square$ \\
97 & \texttt{\small cui2025perfbugs} & {\small Cui 2025} & {\small peer reviewed} & 1 & B & {\small ok} & $\square$ \\
98 & \texttt{\small naylor2026cherisimt} & {\small Naylor et al. 2026} & {\small peer reviewed} & 1 & A & {\small ok} & $\square$ \\
99 & \texttt{\small bae2021rudra} & {\small Bae et al. 2021} & {\small peer reviewed} & 1 & B & {\small ok} & $\square$ \\
100 & \texttt{\small zheng2024security} & {\small Zheng et al. 2024} & {\small peer reviewed} & 1 & A & {\small ok} & $\square$ \\
101 & \texttt{\small schuermann2023encapsulated} & {\small Schuermann et al. 2023} & {\small peer reviewed} & 1 & B & {\small ok} & $\square$ \\
102 & \texttt{\small emre2023aliasing} & {\small Emre et al. 2023} & {\small peer reviewed} & 1 & B & {\small ok} & $\square$ \\
103 & \texttt{\small sharma2024embedded} & {\small Sharma et al. 2024} & {\small peer reviewed} & 1 & A & {\small ok} & $\square$ \\
104 & \texttt{\small mccormack2026unsafe} & {\small McCormack et al. 2026} & {\small peer reviewed} & 1 & A & {\small ok} & $\square$ \\
105 & \texttt{\small hong2025ctorust} & {\small Hong 2025} & {\small peer reviewed} & 1 & B & {\small ok} & $\square$ \\
106 & \texttt{\small huang2017predictability} & {\small Huang et al. 2017} & {\small peer reviewed} & 1 & B & {\small ok} & $\square$ \\
107 & \texttt{\small liu2023l25gc} & {\small Liu et al. 2023} & {\small peer reviewed} & 1 & B & {\small ok} & $\square$ \\
108 & \texttt{\small park2025ctrlbench} & {\small Park 2025} & {\small peer reviewed} & 1 & B & {\small ok} & $\square$ \\
109 & \texttt{\small ardito2021verbosity} & {\small Ardito et al. 2021} & {\small peer reviewed} & 1 & B & {\small ok} & $\square$ \\
110 & \texttt{\small qazi2017pepc} & {\small Qazi et al. 2017} & {\small peer reviewed} & 1 & A & {\small ok} & $\square$ \\
111 & \texttt{\small alsoran} & {\small Larkin} & {\small tool record} & 1 & A & {\small ok} & $\square$ \\
\end{longtable}
\end{small}


\clearpage
\section*{Part 2. Entry records}
\addcontentsline{toc}{section}{Part 2. Entry records}

Entry \textbf{[n]} is reference \textbf{[n]} of the thesis. Citing sentences are printed as
they stand, with the citation shown as \citemark{n} and a cross-reference as \refmark{label},
each located by the chapter or section it sits in.

\auditentry{1}{3gpp38401}{standard, level~A}
\afield{Reference}{\fullcite{3gpp38401}}
\afield{Entry}{\href{https://www.3gpp.org/ftp/Specs/archive/38_series/38.401/38401-h50.zip}{source page}. V17.5.0, Release 17.}
\afield{Cited in}{6 places.}
\begin{citedsentence}\textbf{\small Chapter 1.}\enspace The central unit is split again into control-plane and user-plane functions~\citemark{1, 2}.\end{citedsentence}
\begin{citedsentence}\textbf{\small Section 1.1.}\enspace In the disaggregated architecture of the 3rd Generation Partnership Project (3GPP) it is the control-plane half of the central unit, hosting RRC and managing the UE context lifecycle~\citemark{1, 4}, under the millisecond-scale procedure budgets the 5G system places on control signalling~\citemark{5} and at the event rates and connected-device counts a deployed network sustains.\end{citedsentence}
\begin{citedsentence}\textbf{\small Section 2.1.}\enspace The next-generation radio access network (NG-RAN) architecture splits the gNB into a central unit (gNB-CU) and distributed units, with the control plane (gNB-CU-CP) handling RRC and connection management for each UE~\citemark{1, 2, 4}.\end{citedsentence}
\begin{citedsentence}\textbf{\small Section 2.1.}\enspace These parts communicate over standardised interfaces: NG towards the 5G core, F1 between the CU and the DU, and E1 between the CU-CP and the CU-UP~\citemark{1, 2, 30}.\end{citedsentence}
\begin{citedsentence}\textbf{\small Section 2.1.}\enspace The DU terminates its lower layers and the CU-CP its RRC signalling~\citemark{1, 2}.\end{citedsentence}
\begin{citedsentence}\textbf{\small Section 6.2.}\enspace It is the RRC connection lifecycle the 3GPP specifications define for a UE, from establishment through in-session signalling and bearer setup to reconfiguration and release~\citemark{1, 4}.\end{citedsentence}
\afield{Claim it carries}{the CU/DU split, the CU-CP/CU-UP separation, the F1/E1/NG interfaces, and the gNB-CU-CP as the node hosting RRC}
\afield{Source says}{}
\begin{saidquote}(\texttt{ts\_138401v170500p.pdf}):\end{saidquote}
\begin{saidquote}\begin{itemize}[leftmargin=1.2em,itemsep=1pt,topsep=1pt]
\item p. 11, clause 3.1: \enquote{gNB-CU-Control Plane (gNB-CU-CP): a logical node hosting the RRC and the control plane part of the PDCP protocol of the gNB-CU \ldots{} terminates the E1 interface connected with the gNB-CU-UP and the F1-C interface connected with the gNB-DU.}
\item p. 16, clause 6.1.2 \enquote{Overall architecture for separation of gNB-CU-CP and gNB-CU-UP}, Figure 6.1.2-1.
\item Clause 10.4 \enquote{E1 interface}.
\end{itemize}\end{saidquote}
\afield{Metadata}{[verified] \texttt{38401-h50.zip} in the 3GPP archive is dated \textbf{2023-06-29}, matching \texttt{jun 2023}. The ETSI mirror is \emph{ETSI TS 138 401 V17.5.0 (2023-07)}. URL in the entry resolves.}
\afield{Verdict}{OK. This entry, not \texttt{3gpp38300}, is what carries the CP/UP separation claim.}
\afield{Record}{Sources\_evidence.md, entry record.}
\auditentry{2}{3gpp38300}{standard, level~A}
\afield{Reference}{\fullcite{3gpp38300}}
\afield{Entry}{\href{https://www.3gpp.org/ftp/Specs/archive/38_series/38.300/38300-h50.zip}{source page}. V17.5.0, Release 17.}
\afield{Cited in}{4 places.}
\begin{citedsentence}\textbf{\small Chapter 1.}\enspace The central unit is split again into control-plane and user-plane functions~\citemark{1, 2}.\end{citedsentence}
\begin{citedsentence}\textbf{\small Section 2.1.}\enspace The next-generation radio access network (NG-RAN) architecture splits the gNB into a central unit (gNB-CU) and distributed units, with the control plane (gNB-CU-CP) handling RRC and connection management for each UE~\citemark{1, 2, 4}.\end{citedsentence}
\begin{citedsentence}\textbf{\small Section 2.1.}\enspace These parts communicate over standardised interfaces: NG towards the 5G core, F1 between the CU and the DU, and E1 between the CU-CP and the CU-UP~\citemark{1, 2, 30}.\end{citedsentence}
\begin{citedsentence}\textbf{\small Section 2.1.}\enspace The DU terminates its lower layers and the CU-CP its RRC signalling~\citemark{1, 2}.\end{citedsentence}
\afield{Claim it carries}{alongside 38.401 for architecture and functional split}
\afield{Source says}{}
\begin{saidquote}(\texttt{ts\_138300v170500p.pdf}): p. 21, clause 4.1 \enquote{Overall Architecture} (an NG-RAN node is a gNB or an ng-eNB. Figure 4.1-1), p. 22, clause 4.2 \enquote{Functional Split}, listing gNB functions including Radio Resource Management and \enquote{Connection setup and release}.\end{saidquote}
\afield{Metadata}{[verified] \texttt{38300-h50.zip} dated \textbf{2023-06-30}. ETSI mirror \emph{V17.5.0 (2023-07)}.}
\afield{Note}{38.300 v17.5.0 does \textbf{not} describe the gNB-CU-CP / gNB-CU-UP separation, where it mentions CU functionality it defers to TS 38.401 (e.g. p. 19, p. 28). The pairing in the thesis is fine because 38.401 leads, but 38.300 should not be cited alone for the CP/UP split.}
\afield{Verdict}{OK with that note.}
\afield{Record}{Sources\_evidence.md, entry record.}
\auditentry{3}{polese2023oran}{peer reviewed, level~A}
\afield{Reference}{\fullcite{polese2023oran}}
\afield{Entry}{DOI \href{https://doi.org/10.1109/COMST.2023.3239220}{\texttt{10.1109/COMST.2023.3239220}}.}
\afield{Cited in}{4 places.}
\begin{citedsentence}\textbf{\small Chapter 1.}\enspace Open-RAN pushes the same idea further: the RAN becomes virtualised, multi-vendor software linked by open interfaces~\citemark{3}.\end{citedsentence}
\begin{citedsentence}\textbf{\small Section 1.1.}\enspace This shift is accelerated by Open-RAN disaggregation, which turns the base station into virtualised, multi-vendor software connected by open interfaces~\citemark{3}, so that the control plane a vendor ships increasingly competes on software-engineering quality (reliability, predictable timing, security, and the ease with which a large codebase can be evolved safely) and no longer on purpose-built hardware.\end{citedsentence}
\begin{citedsentence}\textbf{\small Section 2.1.}\enspace Open-RAN initiatives standardise these interfaces further so that each part can be supplied by a different vendor and run as software on commodity hardware~\citemark{3}.\end{citedsentence}
\begin{citedsentence}\textbf{\small Appendix B.}\enspace The fourth strand is the domain and the method: the tail-latency and microsecond-scale literature~\citemark{6, 7, 31, 106}, thread-per-core structuring moving the tail in its own right~\citemark{46}, O-RAN disaggregation~\citemark{3}, control-plane latency as an active target in the core network too~\citemark{107, 108}, adoption evidence pointing at the learning curve rather than run-time cost~\citemark{80, 81, 109}, and standard guidance on repeat-based evaluation and fault injection~\citemark{33, 60}.\end{citedsentence}
\afield{Claim it carries}{O-RAN turns the RAN into virtualised, multi-vendor software connected by open interfaces}
\afield{Source says}{}
\begin{saidquote}(arXiv 2202.01032, p. 1, abstract): \enquote{\textbf{O-RAN promotes virtualized RANs where disaggregated components are connected via open interfaces} and optimized by intelligent controllers \ldots{} \textbf{O-RAN networks can be built with multi-vendor, interoperable components}.}\end{saidquote}
\afield{Metadata}{[verified] Crossref matches every field. (The apostrophe in D'Oro differs typographically only.)}
\afield{Verdict}{OK, a close match to the thesis's phrasing.}
\afield{Record}{Sources\_evidence.md, entry record.}
\auditentry{4}{3gpp38331}{standard, level~C}
\afield{Reference}{\fullcite{3gpp38331}}
\afield{Entry}{\href{https://www.3gpp.org/ftp/Specs/archive/38_series/38.331/38331-h50.zip}{source page}. V17.5.0, Release 17.}
\afield{Cited in}{5 places.}
\begin{citedsentence}\textbf{\small Chapter 1.}\enspace It tracks every connected device, each of which the standards call a user equipment (UE), runs Radio Resource Control (RRC) signalling~\citemark{4}, and coordinates mobility and handovers, all within strict latency limits.\end{citedsentence}
\begin{citedsentence}\textbf{\small Section 1.1.}\enspace In the disaggregated architecture of the 3rd Generation Partnership Project (3GPP) it is the control-plane half of the central unit, hosting RRC and managing the UE context lifecycle~\citemark{1, 4}, under the millisecond-scale procedure budgets the 5G system places on control signalling~\citemark{5} and at the event rates and connected-device counts a deployed network sustains.\end{citedsentence}
\begin{citedsentence}\textbf{\small Section 2.1.}\enspace The next-generation radio access network (NG-RAN) architecture splits the gNB into a central unit (gNB-CU) and distributed units, with the control plane (gNB-CU-CP) handling RRC and connection management for each UE~\citemark{1, 2, 4}.\end{citedsentence}
\begin{citedsentence}\textbf{\small Section 3.1.}\enspace It is modelled on a real Nokia gNB-CU-CP service domain that owns the UE context lifecycle, RRC procedure execution~\citemark{4}, and mobility/handover, which are the functions the 3GPP architecture assigns to the control-plane part of a disaggregated gNB (Section~\refmark{sec:ranarch}).\end{citedsentence}
\begin{citedsentence}\textbf{\small Section 6.2.}\enspace It is the RRC connection lifecycle the 3GPP specifications define for a UE, from establishment through in-session signalling and bearer setup to reconfiguration and release~\citemark{1, 4}.\end{citedsentence}
\afield{Claim it carries}{RRC signalling procedures and the UE lifecycle from establishment to release}
\afield{Source says}{}
\begin{saidquote}The RRC connection control procedures are clause 5.3 (establishment, reconfiguration, re-establishment, release). The document itself was not opened for this audit, it is ~1,000 pages and the claim is structural rather than numeric.\end{saidquote}
\afield{Metadata}{[nit] \texttt{38331-h50.zip} is dated \textbf{2023-07-01}, the entry says June 2023. ETSI mirror is \emph{V17.5.0 (2023-07)}. One-month slip, cosmetic.}
\afield{Verdict}{Metadata: change month to \texttt{jul} if you want it exact.}
\afield{Record}{Sources\_evidence.md, entry record.}
\auditentry{5}{3gpp38913}{standard, level~A}
\afield{Reference}{\fullcite{3gpp38913}}
\afield{Entry}{\href{https://www.3gpp.org/ftp/Specs/archive/38_series/38.913/38913-i00.zip}{source page}. V18.0.0, Release 18. Clause 7.4 sets the control plane latency target, for the transition from a battery-efficient state to the start of continuous data transfer, at 10~ms.}
\afield{Cited in}{4 places.}
\begin{citedsentence}\textbf{\small Chapter 1.}\enspace A single site can run thousands of connection-management and mobility procedures per second, each of which has to complete inside the millisecond-scale budget the 5G system sets for control-plane signalling~\citemark{5}.\end{citedsentence}
\begin{citedsentence}\textbf{\small Section 1.1.}\enspace In the disaggregated architecture of the 3rd Generation Partnership Project (3GPP) it is the control-plane half of the central unit, hosting RRC and managing the UE context lifecycle~\citemark{1, 4}, under the millisecond-scale procedure budgets the 5G system places on control signalling~\citemark{5} and at the event rates and connected-device counts a deployed network sustains.\end{citedsentence}
\begin{citedsentence}\textbf{\small Section 2.1.}\enspace Such control-plane systems must satisfy demanding non-functional requirements: very high event rates (UE mobility, signalling storms), strict and predictable latency bounds that the 5G service and RAN specifications quantify as key performance indicators (KPIs)~\citemark{5, 29}, long-lived software with continuously evolving complexity, and heavily concurrent workloads.\end{citedsentence}
\begin{citedsentence}\textbf{\small Section 3.1.}\enspace It is core RAN layer-3 functionality, and it carries the properties that make production control-plane code difficult: per-UE mutable state, a strict operation lifecycle, high concurrency across UEs, and the latency and reliability targets the 5G system sets for control-plane procedures~\citemark{5, 29}.\end{citedsentence}
\afield{Claim it carries}{the millisecond-scale control-plane budget, and specifically the 10 ms target}
\afield{Source says}{}
\begin{saidquote}(\texttt{tr\_138913v180000p.pdf}, \textbf{p. 24}, clause 7.4 \enquote{Control plane latency}): \enquote{Control plane latency refers to the time to move from a battery efficient state (e.g., IDLE) to start of continuous data transfer (e.g., ACTIVE). \textbf{The target for control plane latency should be 10ms.}} This is exactly the claim at \texttt{Thesis\_v3.tex} and the wording already in the entry's \texttt{note} field.\end{saidquote}
\afield{Metadata}{[nit] month. \texttt{38913-i00.zip} is dated \textbf{2024-04-04}. ETSI mirror is \emph{ETSI TR 138 913 V18.0.0 (2024-05)}. The entry says March 2024. $\rightarrow$ \texttt{month = apr}.}
\afield{Verdict}{Claim verified, fix the month.}
\afield{Record}{Sources\_evidence.md, entry record.}
\auditentry{6}{dean2013tail}{peer reviewed, level~B}
\afield{Reference}{\fullcite{dean2013tail}}
\afield{Entry}{DOI \href{https://doi.org/10.1145/2408776.2408794}{\texttt{10.1145/2408776.2408794}}.}
\afield{Cited in}{7 places.}
\begin{citedsentence}\textbf{\small Chapter 1.}\enspace Every device the site serves is affected~\citemark{6, 7}.\end{citedsentence}
\begin{citedsentence}\textbf{\small Chapter 1.}\enspace emph{Tail latency}, the p99 and p99.9 percentiles rather than the mean, is the property that, more than average speed, decides whether such a system meets its deadlines~\citemark{6, 7}.\end{citedsentence}
\begin{citedsentence}\textbf{\small Section 1.2.}\enspace Second, what governs whether RAN deadlines are met is not average throughput but worst-case and tail behaviour under load and overload~\citemark{6}, a property that is easy to regress and hard to measure correctly.\end{citedsentence}
\begin{citedsentence}\textbf{\small Section 2.2.}\enspace At the scale of a network, tail latency dominates user-visible behaviour: when a single user action fans out into many internal operations, the slowest of them sets the response time, so rare per-operation delays become common per-request delays~\citemark{6}.\end{citedsentence}
\begin{citedsentence}\textbf{\small Section 2.4.}\enspace Java, C\# and Go achieve memory safety through a runtime and garbage collection, which is an excellent trade for much software and a poor one for a tail-latency-critical control plane: collection pauses are a recognised source of tail latency in large systems~\citemark{6}, and a managed runtime gives up the control over memory layout and scheduling this domain relies on.\end{citedsentence}
\begin{citedsentence}\textbf{\small Section 7.6.}\enspace RQ1 is answered with high percentiles and the maximum instead of the mean, which is the standard way of measuring predictability in this literature~\citemark{6}.\end{citedsentence}
\begin{citedsentence}\textbf{\small Appendix B.}\enspace The fourth strand is the domain and the method: the tail-latency and microsecond-scale literature~\citemark{6, 7, 31, 106}, thread-per-core structuring moving the tail in its own right~\citemark{46}, O-RAN disaggregation~\citemark{3}, control-plane latency as an active target in the core network too~\citemark{107, 108}, adoption evidence pointing at the learning curve rather than run-time cost~\citemark{80, 81, 109}, and standard guidance on repeat-based evaluation and fault injection~\citemark{33, 60}.\end{citedsentence}
\afield{Claim it carries}{fan-out makes a rare per-operation latency a common per-request latency}
\afield{Source says}{}
\begin{saidquote}Crossref-deposited abstract: \enquote{Software techniques that tolerate latency variability are vital to building responsive large-scale Web services.} The fan-out argument is the paper's central section (\enquote{The Tail at Scale}, \S{}\enquote{Why Variability Exists} and \S{}\enquote{Reducing Component Variability}).\end{saidquote}
\afield{Metadata}{[verified] Crossref matches every field.}
\afield{Verdict}{OK.}
\afield{Record}{Sources\_evidence.md, entry record.}
\auditentry{7}{barroso2017microseconds}{peer reviewed, level~B}
\afield{Reference}{\fullcite{barroso2017microseconds}}
\afield{Entry}{DOI \href{https://doi.org/10.1145/3015146}{\texttt{10.1145/3015146}}.}
\afield{Cited in}{4 places.}
\begin{citedsentence}\textbf{\small Chapter 1.}\enspace Every device the site serves is affected~\citemark{6, 7}.\end{citedsentence}
\begin{citedsentence}\textbf{\small Chapter 1.}\enspace emph{Tail latency}, the p99 and p99.9 percentiles rather than the mean, is the property that, more than average speed, decides whether such a system meets its deadlines~\citemark{6, 7}.\end{citedsentence}
\begin{citedsentence}\textbf{\small Section 2.2.}\enspace Modern data-centre and networking software increasingly operates in a regime where the relevant latencies are measured in microseconds, a range in which conventional operating-system and software overheads are no longer negligible~\citemark{7}.\end{citedsentence}
\begin{citedsentence}\textbf{\small Appendix B.}\enspace The fourth strand is the domain and the method: the tail-latency and microsecond-scale literature~\citemark{6, 7, 31, 106}, thread-per-core structuring moving the tail in its own right~\citemark{46}, O-RAN disaggregation~\citemark{3}, control-plane latency as an active target in the core network too~\citemark{107, 108}, adoption evidence pointing at the learning curve rather than run-time cost~\citemark{80, 81, 109}, and standard guidance on repeat-based evaluation and fault injection~\citemark{33, 60}.\end{citedsentence}
\afield{Claim it carries}{at microsecond scale, conventional OS and software overheads stop being negligible}
\afield{Source says}{}
\begin{saidquote}Crossref-deposited abstract: \enquote{\textbf{Microsecond-scale I/O means tension between performance and productivity} that will need new latency-mitigating ideas, including in hardware.}\end{saidquote}
\afield{Metadata}{[verified] Crossref matches.}
\afield{Verdict}{OK.}
\afield{Record}{Sources\_evidence.md, entry record.}
\auditentry{8}{msrc2019memory}{grey, level~A}
\afield{Reference}{\fullcite{msrc2019memory}}
\afield{Entry}{\href{https://github.com/microsoft/MSRC-Security-Research/tree/master/presentations/2019_02_BlueHatIL}{source page}. Approximately 70\% of Microsoft CVEs are memory-safety issues.}
\afield{Cited in}{5 places.}
\begin{citedsentence}\textbf{\small Chapter 1.}\enspace They also account for a large share of security vulnerabilities~\citemark{8, 9}.\end{citedsentence}
\begin{citedsentence}\textbf{\small Section 1.1.}\enspace Large software organisations report that the majority of their serious security vulnerabilities stem from memory-safety errors, and that adopting memory-safe languages measurably reduces them~\citemark{8, 9, 16, 17}.\end{citedsentence}
\begin{citedsentence}\textbf{\small Section 2.3.}\enspace Memory-safety issues account for roughly 70\,\% of the security vulnerabilities reported at Microsoft~\citemark{8}, and Google reports a sharp drop in such vulnerabilities as memory-safe languages are adopted in Android, to below a fifth of that platform's total vulnerabilities by 2025~\citemark{9, 16}.\end{citedsentence}
\begin{citedsentence}\textbf{\small Appendix B.}\enspace The safety strand is settled the same partial way, with machine-checked foundations for the safe fragment and the aliasing discipline \texttt{unsafe} code has to respect~\citemark{11, 40, 41}, confirmation that \texttt{unsafe} is used sparingly and that memory-safety CVEs always involve it~\citemark{27, 42, 43}, auditing that finds defects at scale even so and unsafety spreading through dependency chains~\citemark{99, 100, 92, 44}, and the reports motivating the shift set against a long record of memory corruption resisting runtime mitigation~\citemark{8, 9, 17, 19, 20, 34}.\end{citedsentence}
\begin{citedsentence}\textbf{\small Appendix G.}\enspace Adopting Rust removes use-after-free, double-free, iterator invalidation, and data races from the defect budget entirely, and those are the classes that dominate the security statistics~\citemark{8, 9}.\end{citedsentence}
\afield{Claim it carries}{most sharply at 751: \enquote{memory-safety issues account for roughly 70 \% of the security vulnerabilities reported at Microsoft}}
\afield{Source says}{}
\begin{saidquote}(deck \textbf{p. 11}): \enquote{\textbf{~70\% of the vulnerabilities addressed through a security update each year continue to be memory safety issues}}, above the chart \enquote{\% of memory safety vs. non-memory safety CVEs by patch year}, 2006--2018.\end{saidquote}
\afield{Metadata}{[verified] Matt Miller, MSRC, February 2019. The repository path in the URL contains the deck: \texttt{2019\_02 - BlueHatIL - Trends, challenge, and shifts in software vulnerability mitigation.pdf}.}
\afield{Verdict}{OK. The source says \enquote{addressed through a security update} rather than \enquote{reported}, the entry's \texttt{note} (\enquote{Approximately 70 \% of Microsoft CVEs are memory-safety issues}) matches the chart title.}
\afield{Record}{Sources\_evidence.md, entry record.}
\auditentry{9}{google2022android}{grey, level~A}
\afield{Reference}{\fullcite{google2022android}}
\afield{Entry}{\href{https://security.googleblog.com/2022/12/memory-safe-languages-in-android-13.html}{source page}.}
\afield{Cited in}{5 places.}
\begin{citedsentence}\textbf{\small Chapter 1.}\enspace They also account for a large share of security vulnerabilities~\citemark{8, 9}.\end{citedsentence}
\begin{citedsentence}\textbf{\small Section 1.1.}\enspace Large software organisations report that the majority of their serious security vulnerabilities stem from memory-safety errors, and that adopting memory-safe languages measurably reduces them~\citemark{8, 9, 16, 17}.\end{citedsentence}
\begin{citedsentence}\textbf{\small Section 2.3.}\enspace Memory-safety issues account for roughly 70\,\% of the security vulnerabilities reported at Microsoft~\citemark{8}, and Google reports a sharp drop in such vulnerabilities as memory-safe languages are adopted in Android, to below a fifth of that platform's total vulnerabilities by 2025~\citemark{9, 16}.\end{citedsentence}
\begin{citedsentence}\textbf{\small Appendix B.}\enspace The safety strand is settled the same partial way, with machine-checked foundations for the safe fragment and the aliasing discipline \texttt{unsafe} code has to respect~\citemark{11, 40, 41}, confirmation that \texttt{unsafe} is used sparingly and that memory-safety CVEs always involve it~\citemark{27, 42, 43}, auditing that finds defects at scale even so and unsafety spreading through dependency chains~\citemark{99, 100, 92, 44}, and the reports motivating the shift set against a long record of memory corruption resisting runtime mitigation~\citemark{8, 9, 17, 19, 20, 34}.\end{citedsentence}
\begin{citedsentence}\textbf{\small Appendix G.}\enspace Adopting Rust removes use-after-free, double-free, iterator invalidation, and data races from the defect budget entirely, and those are the classes that dominate the security statistics~\citemark{8, 9}.\end{citedsentence}
\afield{Source says}{}
\begin{saidquote}for identity and date, the specific \enquote{below a fifth by 2025} figure at \texttt{Ha\_Vien\_Mastersthesis\_2026.tex} is carried by \texttt{vanderstoep2025android} (below), which is cited alongside it.\end{saidquote}
\afield{Metadata}{[verified] Page title and URL confirm: \emph{Memory Safe Languages in Android 13}, Google Online Security Blog, December 2022, Jeff Vander Stoep.}
\afield{Verdict}{OK.}
\afield{Record}{Sources\_evidence.md, entry record.}
\auditentry{10}{klabnik2023rust}{book, level~C}
\afield{Reference}{\fullcite{klabnik2023rust}}
\afield{Entry}{\href{https://openlibrary.org/isbn/9781718503106}{source page}.}
\afield{Cited in}{2 places.}
\begin{citedsentence}\textbf{\small Chapter 1.}\enspace Its compiler enforces memory safety and data-race freedom at compile time, on foundations that have been checked by machine~\citemark{10, 11}.\end{citedsentence}
\begin{citedsentence}\textbf{\small Section 2.4.}\enspace Rust enforces memory safety and data-race freedom at compile time through its ownership and borrowing model~\citemark{10}, and ships with modern tooling in Cargo and integrated testing.\end{citedsentence}
\afield{Claim it carries}{ownership and borrowing, safety at compile time}
\afield{Source says}{}
\begin{saidquote}Book not opened. The relevant chapters are 4 (ownership) and 16 (fearless concurrency), add a chapter to the citation if you want it tighter.\end{saidquote}
\afield{Metadata}{[verified] OpenLibrary record for ISBN 9781718503106: \emph{Rust Programming Language, 2nd Edition}, Steve Klabnik and Carol Nichols, No Starch Press, 2023. Catalogue URL \texttt{https://openlibrary.org/isbn/9781718503106}.}
\afield{Verdict}{OK.}
\afield{Record}{Sources\_evidence.md, entry record.}
\auditentry{11}{jung2018rustbelt}{peer reviewed, level~A}
\afield{Reference}{\fullcite{jung2018rustbelt}}
\afield{Entry}{DOI \href{https://doi.org/10.1145/3158154}{\texttt{10.1145/3158154}}.}
\afield{Cited in}{3 places.}
\begin{citedsentence}\textbf{\small Chapter 1.}\enspace Its compiler enforces memory safety and data-race freedom at compile time, on foundations that have been checked by machine~\citemark{10, 11}.\end{citedsentence}
\begin{citedsentence}\textbf{\small Section 2.4.}\enspace The soundness of the safe fragment of the language, including the role of \texttt{unsafe} as an explicitly marked escape hatch, has been established formally by the RustBelt project~\citemark{11}, and the aliasing discipline that even \texttt{unsafe} code must respect has been given a precise operational semantics~\citemark{41}.\end{citedsentence}
\begin{citedsentence}\textbf{\small Appendix B.}\enspace The safety strand is settled the same partial way, with machine-checked foundations for the safe fragment and the aliasing discipline \texttt{unsafe} code has to respect~\citemark{11, 40, 41}, confirmation that \texttt{unsafe} is used sparingly and that memory-safety CVEs always involve it~\citemark{27, 42, 43}, auditing that finds defects at scale even so and unsafety spreading through dependency chains~\citemark{99, 100, 92, 44}, and the reports motivating the shift set against a long record of memory corruption resisting runtime mitigation~\citemark{8, 9, 17, 19, 20, 34}.\end{citedsentence}
\afield{Claim it carries}{machine-checked foundations for safe Rust}
\afield{Source says}{}
\begin{saidquote}(p. 1, abstract, restated p. 2): \enquote{we give the \textbf{first formal (and machine-checked) safety proof} for a language representing a realistic subset of Rust.}\end{saidquote}
\afield{Metadata}{[verified] The paper's own footer reads \enquote{Proceedings of the ACM on Programming Languages, Vol. 2, No. POPL, \textbf{Article 66. Publication date: January 2018}}, so \texttt{year = 2018} and \texttt{pages = 66:1--66:34} are both right. (Crossref's \enquote{2017} is the online-first date, do not \enquote{correct} the entry to it.)}
\afield{Verdict}{OK.}
\afield{Record}{Sources\_evidence.md, entry record.}
\auditentry{12}{balasubramanian2017rust}{peer reviewed, level~B}
\afield{Reference}{\fullcite{balasubramanian2017rust}}
\afield{Entry}{DOI \href{https://doi.org/10.1145/3102980.3103006}{\texttt{10.1145/3102980.3103006}}.}
\afield{Cited in}{7 places.}
\begin{citedsentence}\textbf{\small Chapter 1.}\enspace The intended performance is competitive with C++~\citemark{12}.\end{citedsentence}
\begin{citedsentence}\textbf{\small Section 1.2.}\enspace The available evidence is largely general-purpose and average-case~\citemark{12, 27}, leaving these three questions (predictability, prevented-defect scope, and migration cost) open for a representative subsystem.\end{citedsentence}
\begin{citedsentence}\textbf{\small Section 2.4.}\enspace Systems research argues Rust attains safety without the overhead of garbage collection or runtime checks, and even enables efficiency techniques that are impractical in C~\citemark{12}.\end{citedsentence}
\begin{citedsentence}\textbf{\small Section 2.4.}\enspace Rust is the first widely supported language to combine compile-time memory and thread safety with the absence of a garbage collector and throughput and memory use alongside C and C++~\citemark{12, 40, 45}.\end{citedsentence}
\begin{citedsentence}\textbf{\small Section 7.2.}\enspace The median result is consistent with the wider evidence that Rust's safety checks are largely absorbed by modern out-of-order CPUs~\citemark{12, 73}.\end{citedsentence}
\begin{citedsentence}\textbf{\small Section 7.5.}\enspace The ``zero-cost'' framing~\citemark{12} is therefore best read as neither established nor refuted in general, but as a property of a particular program on a particular workload.\end{citedsentence}
\begin{citedsentence}\textbf{\small Appendix B.}\enspace Compiled systems languages sit close on time and energy~\citemark{45, 93}, Rust's safety is cheap at run time~\citemark{12, 94}, and that holds where timing is tightest, in embedded hypervisors and in packet processing within a few percent of a C reference~\citemark{95, 73, 96}.\end{citedsentence}
\afield{Claim it carries}{the zero-cost framing, and as of the 31 August citation-attachment pass, Chapter 1's clause on performance intended to be competitive with C++}
\afield{Source says}{}
\begin{saidquote}(Crossref-deposited abstract): \enquote{Rust is unique in that it \textbf{enforces safety without runtime overhead}, most importantly, without the overhead of garbage collection. While \textbf{zero-cost safety} is remarkable on its own, we argue that the superpowers of Rust go beyond safety. In particular, Rust's linear type system enables capabilities that cannot be implemented efficiently in traditional languages, both safe and unsafe.}\end{saidquote}
\afield{Verdict}{OK, and this one now matters more than it did: the citation-attachment pass moved Chapter 1's performance clause onto this entry, and the deposited abstract carries that clause directly. \S{}2.4's stronger reading, that Rust \enquote{even enables efficiency techniques that are impractical in C}, is the abstract's own \enquote{capabilities that cannot be implemented efficiently in traditional languages}.}
\afield{Record}{Sources\_evidence.md, entry record.}
\auditentry{13}{itu2023imt2030}{standard, level~A}
\afield{Reference}{\fullcite{itu2023imt2030}}
\afield{Entry}{\href{https://www.itu.int/rec/R-REC-M.2160/en}{source page}. Approved 13 November 2023. Defines the IMT-2030 (``6G'') framework and its development timeline.}
\afield{Cited in}{2 places.}
\begin{citedsentence}\textbf{\small Section 1.1.}\enspace International Mobile Telecommunications-2030 (IMT-2030), the International Telecommunication Union (ITU) framework for 6G, lists artificial intelligence and communication among its six usage scenarios and ubiquitous intelligence among its four design principles~\citemark{13}.\end{citedsentence}
\begin{citedsentence}\textbf{\small Section 1.1.}\enspace IMT-2030 names security and resilience as one of its four design principles, expects the system to be secure by design, defines resilience as the ability to keep operating correctly during and after a disruptive event, and carries the ultra-reliable low-latency scenario forward in strengthened form~\citemark{13}.\end{citedsentence}
\afield{Source says}{}
\begin{saidquote}for the metadata. The ITU page gives M.2160-0 (11/2023), approved 2023-11, in force, title as in the entry\end{saidquote}
\afield{Record}{Sources\_evidence.md, 24 August verification pass.}
\auditentry{14}{polese2026airan}{peer reviewed, level~B}
\afield{Reference}{\fullcite{polese2026airan}}
\afield{Entry}{DOI \href{https://doi.org/10.1109/MCOM.001.2500438}{\texttt{10.1109/MCOM.001.2500438}}.}
\afield{Cited in}{1 place.}
\begin{citedsentence}\textbf{\small Section 1.1.}\enspace The architecture forming around that is AI-RAN, in which telecommunications and AI workloads are orchestrated together on shared, accelerated infrastructure, so the sites that run the radio also host distributed edge AI platforms~\citemark{14}.\end{citedsentence}
\afield{Source says}{}
\begin{saidquote}Crossref confirms IEEE Communications Magazine 64(7):90--95, 2026, five authors as listed\end{saidquote}
\afield{Record}{Sources\_evidence.md, 24 August verification pass.}
\auditentry{15}{airan2026orchestrator}{grey, level~A}
\afield{Reference}{\fullcite{airan2026orchestrator}}
\afield{Entry}{\href{https://ai-ran.org/documents/AI-RAN-Platform-Infrastructure-Orchestrator.pdf}{source page}. Version 10-rc1, clause 2.2. Defines AI-and-RAN as one common IT infrastructure accommodating both AI and RAN workloads, which the infrastructure is required to run concurrently, and notes that the two kinds of workload may have different service-level agreements. Contributors from Ericsson, Nokia, Nvidia, Red Hat, Samsung, SoftBank, zTouch Networks and Northeastern University. Accessed August 2026.}
\afield{Cited in}{1 place.}
\begin{citedsentence}\textbf{\small Section 1.1.}\enspace The AI-RAN Alliance states the operational requirement directly: one common infrastructure has to run both kinds of workload concurrently, and their service-level agreements may differ~\citemark{15}.\end{citedsentence}
\afield{Source says}{}
\begin{saidquote}clause 2.2: \enquote{The concept of 'AI-and-RAN' means using a common IT infrastructure to accommodate both AI and RAN workloads. Therefore, there is a fundamental requirement for the infrastructure to run AI and RAN workloads concurrently. It is important to note that AI and RAN workloads may have different SLAs.} The contributor list in the note matches the eight organisations on page 2. Only the date was wrong, see above\end{saidquote}
\afield{Record}{Sources\_evidence.md, 24 August verification pass.}
\auditentry{16}{vanderstoep2025android}{grey, level~A}
\afield{Reference}{\fullcite{vanderstoep2025android}}
\afield{Entry}{\href{https://blog.google/security/rust-in-android-move-fast-fix-things/}{source page}.}
\afield{Cited in}{2 places.}
\begin{citedsentence}\textbf{\small Section 1.1.}\enspace Large software organisations report that the majority of their serious security vulnerabilities stem from memory-safety errors, and that adopting memory-safe languages measurably reduces them~\citemark{8, 9, 16, 17}.\end{citedsentence}
\begin{citedsentence}\textbf{\small Section 2.3.}\enspace Memory-safety issues account for roughly 70\,\% of the security vulnerabilities reported at Microsoft~\citemark{8}, and Google reports a sharp drop in such vulnerabilities as memory-safe languages are adopted in Android, to below a fifth of that platform's total vulnerabilities by 2025~\citemark{9, 16}.\end{citedsentence}
\afield{Claim it carries}{memory-safety vulnerabilities falling below a fifth of Android's total by 2025}
\afield{Source says}{}
\begin{saidquote}(page metadata/description): \enquote{The 2025 data continues to validate the approach, with \textbf{memory safety vulnerabilities falling below 20\% of total vulnerabilities for the first time}.}\end{saidquote}
\afield{Metadata}{[verified] blog.google, author metadata \enquote{Jeff Vander Stoep}, November 2025.}
\afield{Verdict}{OK, the thesis's \enquote{below a fifth \ldots{} by 2025} is the source's own claim.}
\afield{Record}{Sources\_evidence.md, entry record.}
\auditentry{17}{google2024securebydesign}{grey, level~B}
\afield{Reference}{\fullcite{google2024securebydesign}}
\afield{Entry}{\href{https://research.google/pubs/secure-by-design-googles-perspective-on-memory-safety/}{source page}.}
\afield{Cited in}{2 places.}
\begin{citedsentence}\textbf{\small Section 1.1.}\enspace Large software organisations report that the majority of their serious security vulnerabilities stem from memory-safety errors, and that adopting memory-safe languages measurably reduces them~\citemark{8, 9, 16, 17}.\end{citedsentence}
\begin{citedsentence}\textbf{\small Appendix B.}\enspace The safety strand is settled the same partial way, with machine-checked foundations for the safe fragment and the aliasing discipline \texttt{unsafe} code has to respect~\citemark{11, 40, 41}, confirmation that \texttt{unsafe} is used sparingly and that memory-safety CVEs always involve it~\citemark{27, 42, 43}, auditing that finds defects at scale even so and unsafety spreading through dependency chains~\citemark{99, 100, 92, 44}, and the reports motivating the shift set against a long record of memory corruption resisting runtime mitigation~\citemark{8, 9, 17, 19, 20, 34}.\end{citedsentence}
\afield{Source says}{}
\begin{saidquote}Title and venue support the claim as used (vendor position on memory safety). No page-level quote taken.\end{saidquote}
\afield{Metadata}{[verified] research.google landing page: \emph{Secure by Design: Google's Perspective on Memory Safety}, \textbf{Google Security Engineering (2024)}, authors Alex Rebert and Christoph Kern.}
\afield{Verdict}{OK.}
\afield{Record}{Sources\_evidence.md, entry record.}
\auditentry{18}{watson2025memorysafety}{peer reviewed, level~A}
\afield{Reference}{\fullcite{watson2025memorysafety}}
\afield{Entry}{DOI \href{https://doi.org/10.1145/3708553}{\texttt{10.1145/3708553}}.}
\afield{Cited in}{1 place.}
\begin{citedsentence}\textbf{\small Section 1.1.}\enspace The case no longer rests on vendor reporting alone: a group spanning academia and industry has argued that memory safety now needs standardised terminology and requirements if it is to be specified and procured at all~\citemark{18}.\end{citedsentence}
\afield{Claim it carries}{memory safety needs standardised terminology and requirements}
\afield{Source says}{}
\begin{saidquote}(UCAM-CL-TR-996, p. 1 and p. 3, the extended version of the same paper): \enquote{It is time to standardize principles and practices for software memory safety \ldots{} we explore memory-safety standardization, which we argue is an essential step\ldots{}}. Crossref's deposited CACM abstract: \enquote{\ldots{}argue that standardization is an essential next step to achieving universal strong memory safety.}\end{saidquote}
\afield{Metadata}{[wrong] author list. Crossref and OpenAlex both give \textbf{21} authors, the CACM abstract itself opens \enquote{\textbf{Twenty-one authors}, spanning academia and industry\ldots{}}. The entry lists 22: the extra name is \textbf{Christoph Kern}, who appears on the cover of the \emph{extended} version (Cambridge Computer Laboratory Technical Report UCAM-CL-TR-996) but not on the CACM article. CACM also gives \enquote{Murali Vijayaraghavan} (entry: \enquote{Muralidaran}), and orders Chisnall before Chen and Richardson before Rebert. Volume, issue, pages, year and DOI are all correct.}
\afield{Verdict}{Claim verified. \textbf{fix the author list} (fix \#3).}
\afield{Record}{Sources\_evidence.md, entry record.}
\auditentry{19}{cisa2023memorysafe}{standard, level~A}
\afield{Reference}{\fullcite{cisa2023memorysafe}}
\afield{Entry}{\href{https://www.cisa.gov/resources-tools/resources/case-memory-safe-roadmaps}{source page}.}
\afield{Cited in}{3 places.}
\begin{citedsentence}\textbf{\small Section 1.1.}\enspace This concern has risen to the level of national policy: cybersecurity agencies now urge software manufacturers to publish ``memory-safe roadmaps'' and to migrate away from memory-unsafe languages such as C and C++~\citemark{19, 20}.\end{citedsentence}
\begin{citedsentence}\textbf{\small Section 1.4.}\enspace The intended readers are the vendor making the language decision, the standards discussion around memory-safe roadmaps~\citemark{19, 20}, and anyone repeating a comparison of this shape, who can reuse the specification, the pair, the harness and the catalogue.\end{citedsentence}
\begin{citedsentence}\textbf{\small Appendix B.}\enspace The safety strand is settled the same partial way, with machine-checked foundations for the safe fragment and the aliasing discipline \texttt{unsafe} code has to respect~\citemark{11, 40, 41}, confirmation that \texttt{unsafe} is used sparingly and that memory-safety CVEs always involve it~\citemark{27, 42, 43}, auditing that finds defects at scale even so and unsafety spreading through dependency chains~\citemark{99, 100, 92, 44}, and the reports motivating the shift set against a long record of memory corruption resisting runtime mitigation~\citemark{8, 9, 17, 19, 20, 34}.\end{citedsentence}
\afield{Claim it carries}{the \enquote{publish memory-safe roadmaps} claim}
\afield{Source says}{}
\begin{saidquote}(\textbf{p. 2}): \enquote{the agencies urge software manufacturers to \textbf{create and publish memory safe roadmaps} that detail how they will eliminate memory safety vulnerabilities in their products.}\end{saidquote}
\afield{Metadata}{[verified] for content, [nit] for the \texttt{author} field. Cover page (\textbf{p. 1}) gives \enquote{Publication: December 2023} and the eight authoring agencies in exactly the order the entry's \texttt{institution} field lists them: US CISA, US NSA, US FBI, Australian Signals Directorate's ACSC, Canadian Centre for Cyber Security, UK NCSC, New Zealand NCSC, CERT NZ. The \texttt{author} field, however, is a single braced string ending \enquote{and others} and will print as a run-on name.}
\afield{Verdict}{Claim verified, tidy the \texttt{author} field (fix \#9).}
\afield{Record}{Sources\_evidence.md, entry record.}
\auditentry{20}{oncd2024building}{standard, level~A}
\afield{Reference}{\fullcite{oncd2024building}}
\afield{Entry}{\href{https://bidenwhitehouse.archives.gov/wp-content/uploads/2024/02/Final-ONCD-Technical-Report.pdf}{source page}.}
\afield{Cited in}{3 places.}
\begin{citedsentence}\textbf{\small Section 1.1.}\enspace This concern has risen to the level of national policy: cybersecurity agencies now urge software manufacturers to publish ``memory-safe roadmaps'' and to migrate away from memory-unsafe languages such as C and C++~\citemark{19, 20}.\end{citedsentence}
\begin{citedsentence}\textbf{\small Section 1.4.}\enspace The intended readers are the vendor making the language decision, the standards discussion around memory-safe roadmaps~\citemark{19, 20}, and anyone repeating a comparison of this shape, who can reuse the specification, the pair, the harness and the catalogue.\end{citedsentence}
\begin{citedsentence}\textbf{\small Appendix B.}\enspace The safety strand is settled the same partial way, with machine-checked foundations for the safe fragment and the aliasing discipline \texttt{unsafe} code has to respect~\citemark{11, 40, 41}, confirmation that \texttt{unsafe} is used sparingly and that memory-safety CVEs always involve it~\citemark{27, 42, 43}, auditing that finds defects at scale even so and unsafety spreading through dependency chains~\citemark{99, 100, 92, 44}, and the reports motivating the shift set against a long record of memory corruption resisting runtime mitigation~\citemark{8, 9, 17, 19, 20, 34}.\end{citedsentence}
\afield{Claim it carries}{the policy push away from memory-unsafe languages}
\afield{Source says}{}
\begin{saidquote}(\textbf{p. 7}): \enquote{Experts have identified a few programming languages that both lack traits associated with memory safety and also have high proliferation across critical systems, such as C and C++}, and \enquote{in most cases, \textbf{using a memory safe programming language is the most efficient way to substantially improve software security}}.\end{saidquote}
\afield{Metadata}{[verified] Cover page: \enquote{BACK TO THE BUILDING BLOCKS: A PATH TOWARD SECURE AND MEASURABLE SOFTWARE, \textbf{FEBRUARY 2024}}.}
\afield{Verdict}{OK.}
\afield{Record}{Sources\_evidence.md, entry record.}
\auditentry{21}{eu2024cra}{regulation, level~A}
\afield{Reference}{\fullcite{eu2024cra}}
\afield{Entry}{\href{https://eur-lex.europa.eu/eli/reg/2024/2847/oj}{source page}.}
\afield{Cited in}{1 place.}
\begin{citedsentence}\textbf{\small Section 1.1.}\enspace The European Union (EU) Cyber Resilience Act places secure-by-design and lifecycle vulnerability-handling obligations on manufacturers of products with digital elements, with the vulnerability-reporting obligations applying from September 2026 and most of the remainder from December 2027~\citemark{21}.\end{citedsentence}
\afield{Claim it carries}{\enquote{with reporting obligations applying from September 2026 and the remainder from December 2027}}
\afield{Source says}{}
\begin{saidquote}(\textbf{Article 71(2)}, EUR-Lex consolidated text): \enquote{This Regulation shall apply from \textbf{11 December 2027}. However, \textbf{Article 14 shall apply from 11 September 2026} and \textbf{Chapter IV (Articles 35 to 51) shall apply from 11 June 2026}.} Recital (126) says the same in prose.\end{saidquote}
\afield{Metadata}{[verified] OJ L, 2024/2847, published 20 November 2024. ELI URL resolves.}
\afield{Verdict}{Dates are right, but \enquote{the remainder from December 2027} omits the June 2026 tranche (fix \#6).}
\afield{Record}{Sources\_evidence.md, entry record.}
\auditentry{22}{agache2020firecracker}{peer reviewed, level~A}
\afield{Reference}{\fullcite{agache2020firecracker}}
\afield{Entry}{\href{https://www.usenix.org/conference/nsdi20/presentation/agache}{source page}.}
\afield{Cited in}{1 place.}
\begin{citedsentence}\textbf{\small Section 1.1.}\enspace Production hypervisors such as Amazon Web Services (AWS) Firecracker, which underpins its serverless platform, are already written in Rust~\citemark{22}.\end{citedsentence}
\afield{Claim it carries}{production hypervisors such as AWS Firecracker, which underpins its serverless platform, are already written in Rust}
\afield{Source says}{}
\begin{saidquote}(\texttt{nsdi20-paper-agache.pdf}): p. 5. \enquote{Firecracker contains \textbf{approximately 50k lines of Rust code} (96\% fewer lines than QEMU)}, p. 13. \enquote{The use of virtualization, along with side-channel hardening, \textbf{implementation in Rust}, extensive testing and validation\ldots{}} and \enquote{We have been \textbf{running Lambda workloads in production on Firecracker since 2018}.}\end{saidquote}
\afield{Metadata}{[wrong] \textbf{missing author}. The USENIX proceedings page and the paper's title page both read: \enquote{Alexandru Agache, Marc Brooker, \textbf{Andreea Florescu}, Alexandra Iordache, Anthony Liguori, Rolf Neugebauer, Phil Piwonka, and Diana-Maria Popa, Amazon Web Services.} The entry has seven of those eight. Pages 419--434 confirmed by DBLP. (DBLP also omits Florescu, the proceedings page is the better authority.)}
\afield{Verdict}{Claim verified. \textbf{add Andreea Florescu} (fix \#2, applied). Appendix B's performance-strand cluster carried a second citation until 1 September 2026. It was the only member of that cluster with no language-controlled baseline (Rust-Shyper, Emmerich and rpkt each report against a C reference, Firecracker reports against QEMU), and the claim verified above is adoption rather than performance, so it was removed there on the author's call. The list item \enquote{production cloud infrastructure} went with it, since this was its only source. \textbf{Do not restore it to a performance cluster.}}
\afield{Record}{Sources\_evidence.md, entry record.}
\auditentry{23}{delloro2025ran3q}{grey, level~A}
\afield{Reference}{\fullcite{delloro2025ran3q}}
\afield{Entry}{\href{https://www.delloro.com/news/ran-mostly-stable-in-3q-2025/}{source page}. The five largest RAN suppliers by worldwide revenue (Huawei, Ericsson, Nokia, ZTE, Samsung) accounted for 96\% of the 1Q25--3Q25 RAN market, up from 95\% in 2024. Accessed August 2026.}
\afield{Cited in}{1 place.}
\begin{citedsentence}\textbf{\small Section 1.1.}\enspace That market is concentrated in very few hands: the five largest suppliers took 96\% of worldwide RAN revenue over the first three quarters of 2025, with concentration by Herfindahl--Hirschman Index at a ten-year high~\citemark{23, 24}.\end{citedsentence}
\afield{Source says}{}
\begin{saidquote}Press release of 20 November 2025. \enquote{The top five suppliers accounted for 96 percent of the 1Q25-3Q25 RAN market, up from 95 percent in 2024}, and the top five are Huawei, Ericsson, Nokia, ZTE and Samsung. Both numbers and the ordering match\end{saidquote}
\afield{Record}{Sources\_evidence.md, 24 August verification pass.}
\auditentry{24}{delloro2026ran}{grey, level~A}
\afield{Reference}{\fullcite{delloro2026ran}}
\afield{Entry}{\href{https://www.delloro.com/news/ran-market-stabilized-in-2025/}{source page}. Full-year 2025 preliminary findings: supplier rankings unchanged, and overall market concentration measured by the Herfindahl--Hirschman Index reached a ten-year high. Accessed August 2026.}
\afield{Cited in}{1 place.}
\begin{citedsentence}\textbf{\small Section 1.1.}\enspace That market is concentrated in very few hands: the five largest suppliers took 96\% of worldwide RAN revenue over the first three quarters of 2025, with concentration by Herfindahl--Hirschman Index at a ten-year high~\citemark{23, 24}.\end{citedsentence}
\afield{Source says}{}
\begin{saidquote}Press release of 17 February 2026. \enquote{Overall market concentration, as measured by the Herfindahl--Hirschman Index, reached a 10-year high in 2025}, and \enquote{Revenue rankings did not change in 2025}\end{saidquote}
\afield{Record}{Sources\_evidence.md, 24 August verification pass.}
\auditentry{25}{3gpp2024road6g}{grey, level~A}
\afield{Reference}{\fullcite{3gpp2024road6g}}
\afield{Entry}{\href{https://www.3gpp.org/news-events/3gpp-news/sa1-6g-road}{source page}. Written by the chair and the secretary of 3GPP WG SA1. Release 20 carries the 6G studies and Release 21 the normative work. ITU-R sets early 2029 for technology proposals and mid-2030 for the submitted specifications. Accessed August 2026.}
\afield{Cited in}{1 place.}
\begin{citedsentence}\textbf{\small Section 1.1.}\enspace The next one is being specified now: 3GPP has allocated Release 20 to the 6G studies and Release 21 to the normative work, against an ITU calendar that asks for technology proposals in early 2029 and the submitted specifications by mid-2030~\citemark{25}.\end{citedsentence}
\afield{Source says}{}
\begin{saidquote}Jose Almodovar (SA1 Chair) and Alain Sultan (SA1 Secretary), \enquote{First published December 2024, in Highlights Issue 09}. Release 20 for studies and Release 21 for the normative work, against an ITU calendar of early 2029 for technology proposals and mid-2030 for the specifications. The migration and interworking claim sits in the TR 22.870 section the article names. The web page is stamped 2 January 2025, and the entry follows the Highlights issue instead, which the note states\end{saidquote}
\afield{Record}{Sources\_evidence.md, 24 August verification pass.}
\auditentry{26}{cisa2025badpractices}{standard, level~A}
\afield{Reference}{\fullcite{cisa2025badpractices}}
\afield{Entry}{\href{https://www.cisa.gov/resources-tools/resources/product-security-bad-practices}{source page}. Version 2.0. Bad practice 1 covers the development of new product lines in memory-unsafe languages and the memory-safety roadmap expectation, from which only products with an announced end-of-support date before 1 January 2030 are exempt.}
\afield{Cited in}{1 place.}
\begin{citedsentence}\textbf{\small Section 1.1.}\enspace The Cybersecurity and Infrastructure Security Agency (CISA) and the Federal Bureau of Investigation (FBI) classify the development of new product lines in a memory-unsafe language as a bad practice where the product serves critical infrastructure, a scope mobile-network infrastructure sits squarely inside, and they exempt an existing product from the memory-safe roadmap expectation only when its announced support ends before January 2030~\citemark{26}.\end{citedsentence}
\afield{Source says}{}
\begin{saidquote}Version 2.0, January 2025, both confirmed on the revision table. Bad practice 1 covers new product lines for critical infrastructure in a memory-unsafe language. On the exemption the source says: \enquote{Publication of a memory safety roadmap does not apply to products that have an announced end-of-support date that is prior to Jan. 1, 2030.} That is the claim at \texttt{Ha\_Vien\_Mastersthesis\_2026.tex}, at the strength the thesis states it\end{saidquote}
\afield{Record}{Sources\_evidence.md, 24 August verification pass.}
\auditentry{27}{qin2020rust}{peer reviewed, level~A}
\afield{Reference}{\fullcite{qin2020rust}}
\afield{Entry}{DOI \href{https://doi.org/10.1145/3385412.3386036}{\texttt{10.1145/3385412.3386036}}.}
\afield{Cited in}{4 places.}
\begin{citedsentence}\textbf{\small Section 1.2.}\enspace The available evidence is largely general-purpose and average-case~\citemark{12, 27}, leaving these three questions (predictability, prevented-defect scope, and migration cost) open for a representative subsystem.\end{citedsentence}
\begin{citedsentence}\textbf{\small Section 2.4.}\enspace Empirical evidence supports the model in practice: in real codebases \texttt{unsafe} is used sparingly and in localised ways~\citemark{27, 42}, and a study of all Rust memory-safety entries in the Common Vulnerabilities and Exposures (CVE) list found that every one of them but a single compiler bug involved \texttt{unsafe} code~\citemark{43}.\end{citedsentence}
\begin{citedsentence}\textbf{\small Section 7.5.}\enspace Empirical studies of real-world Rust find that memory- and thread-safety bugs still occur in practice, with blocking bugs and misused synchronisation among them~\citemark{27, 92}.\end{citedsentence}
\begin{citedsentence}\textbf{\small Appendix B.}\enspace The safety strand is settled the same partial way, with machine-checked foundations for the safe fragment and the aliasing discipline \texttt{unsafe} code has to respect~\citemark{11, 40, 41}, confirmation that \texttt{unsafe} is used sparingly and that memory-safety CVEs always involve it~\citemark{27, 42, 43}, auditing that finds defects at scale even so and unsafety spreading through dependency chains~\citemark{99, 100, 92, 44}, and the reports motivating the shift set against a long record of memory corruption resisting runtime mitigation~\citemark{8, 9, 17, 19, 20, 34}.\end{citedsentence}
\afield{Source says}{}
\begin{saidquote}(author copy, p. 1, abstract): \enquote{We performed the \textbf{first empirical study of Rust by close, manual inspection of 850 unsafe code usages and 170 bugs} in five open-source Rust projects, five widely-used Rust libraries, two online security databases, and the Rust standard library.}\end{saidquote}
\afield{Metadata}{[verified] Crossref matches.}
\afield{Verdict}{OK, supports both the \enquote{used sparingly} claim and the \enquote{blocking bugs still occur} claim at 3744.}
\afield{Record}{Sources\_evidence.md, entry record.}
\auditentry{28}{meneely2025justuserust}{peer reviewed, level~A}
\afield{Reference}{\fullcite{meneely2025justuserust}}
\afield{Entry}{DOI \href{https://doi.org/10.1109/SVM66695.2025.00008}{\texttt{10.1109/SVM66695.2025.00008}}.}
\afield{Cited in}{2 places.}
\begin{citedsentence}\textbf{\small Section 1.3.}\enspace Every base entry of the Common Weakness Enumeration (CWE) has been mapped onto what safe Rust would prevent, and weighting those classes by the vulnerability history of 68 open-source C projects puts 58.2\% of the classified vulnerabilities into classes safe Rust would have made virtually impossible~\citemark{28}.\end{citedsentence}
\begin{citedsentence}\textbf{\small Section 7.2.}\enspace The taxonomy-scale mapping cited in Section~\refmark{sec:rqs} names that premise as its own first limitation, because a class counted as virtually impossible stays impossible only while the converted code stays safe~\citemark{28}.\end{citedsentence}
\afield{Claim it carries}{the boundary RQ2 sets out to locate has been estimated once at taxonomy scale, and that estimate assumes the converted code stays safe, which is the assumption the seam double-free withdraws}
\afield{Source says}{}
\begin{saidquote}(the paper as published, read from the Helsinki-licensed Xplore PDF, page numbers below are the printed folios 25--32, not PDF pages):\end{saidquote}
\begin{saidquote}\begin{itemize}[leftmargin=1.2em,itemsep=1pt,topsep=1pt]
\item The figure the thesis quotes, p. 25 (abstract) and again as a display claim on p. 30: \enquote{We estimate that \textbf{58.2\% of historical vulnerabilities in our selection of open source C projects would have been virtually impossible in Rust}.}
\item Its denominator, p. 30 (\S{}V-C): 4,507 vulnerabilities are associated with the 68 projects and \enquote{\textbf{2,257 had a CWE designation}}, which is what Figure 2 weights. The thesis therefore writes \emph{classified} vulnerabilities and not \emph{recorded} ones.
\item The mapping, p. 27 (\S{}IV-A): \enquote{\textbf{We counted 520 CWE Base entries} and had the first author map each one.} Every base entry is covered, which is what \enquote{at the scale of a whole weakness taxonomy} reports.
\item The assumption, p. 32 (\S{}VI \emph{Limitations}, first item, \emph{Safe Rust}): \enquote{all of the Virtually Impossible designations \textbf{require safe Rust}. Automated conversion tooling or careless implementation may simply mark some code as unsafe, which \textbf{brings back these vulnerabilities as possibilities}.} This is the clause the RQ2 discussion cites, and it is the authors' own first limitation rather than an inference from their result.
\item Scope, same section (\emph{Security Only}): \enquote{This work only focuses on security, \textbf{not performance}, correctness, maintainability}, so it bounds RQ2 and says nothing about RQ1 or RQ3.
\end{itemize}\end{saidquote}
\begin{saidquote}\textbf{Three cautions for anyone re-quoting this source.} It prints the Virtually Impossible share of the taxonomy as 14.8\% on p. 28 where its own counts give 72 of 520, so quote the counts and not that percentage. Its Summary on p. 32 says the study covers \enquote{75 prominent, mature open source projects} where the abstract, \S{}V-B and \S{}V-C all say 68. And \textbf{the CWE identifiers in its prose do not match the names beside them} (p. 28 gives CWE-401 as \enquote{Use of Pointer Subtraction to Determine Size}, which is CWE-469, and p. 29 gives CWE-165 for integer wraparound, which is CWE-190), while Table II's identifiers on p. 31 are correct. Take a CWE number from the table only. The thesis quotes 68 and 58.2\%, the pair the abstract and the results section agree on, and no CWE number from this source at all. \textbf{One cell of its mapping this thesis contradicts.} Table II classifies CWE-401 as Virtually Impossible where the \texttt{leak\_cycle} boundary class compiles in safe Rust and leaks. The disagreement is recorded in \texttt{LITERATURE\_NOTES.md} \S{}A.1 of the 8 September 2026 addition, with MITRE's own reference-count example for that CWE and the paper's own method note that predicts it. Nothing in the thesis currently states this, so it is an available addition and not an outstanding defect.\end{saidquote}
\afield{Metadata}{[verified] DOI \texttt{10.1109/SVM66695.2025.00008} resolves. IEEE/ACM 3rd International Workshop on Software Vulnerability Management, co-located with ICSE 2025, Ottawa, 3 May 2025. Author order taken from the publisher record and from the first author's ORCID, both of which give Meneely, Green, Jaafari, Fluet, Keller. \textbf{not} the reading-order of the paper's own two-column author block, which puts Fluet second. Pages 25--32 confirmed against the workshop's table of contents.}
\afield{Verdict}{Entry OK. All four citing sentences carry the best-case framing with the number, which is the source's own reading of it.}
\afield{Record}{Sources\_evidence.md, entry record.}
\auditentry{29}{3gpp22261}{standard, level~C}
\afield{Reference}{\fullcite{3gpp22261}}
\afield{Entry}{\href{https://www.3gpp.org/ftp/Specs/archive/22_series/22.261/22261-i60.zip}{source page}. V18.6.0, Release 18.}
\afield{Cited in}{2 places.}
\begin{citedsentence}\textbf{\small Section 2.1.}\enspace Such control-plane systems must satisfy demanding non-functional requirements: very high event rates (UE mobility, signalling storms), strict and predictable latency bounds that the 5G service and RAN specifications quantify as key performance indicators (KPIs)~\citemark{5, 29}, long-lived software with continuously evolving complexity, and heavily concurrent workloads.\end{citedsentence}
\begin{citedsentence}\textbf{\small Section 3.1.}\enspace It is core RAN layer-3 functionality, and it carries the properties that make production control-plane code difficult: per-UE mutable state, a strict operation lifecycle, high concurrency across UEs, and the latency and reliability targets the 5G system sets for control-plane procedures~\citemark{5, 29}.\end{citedsentence}
\afield{Claim it carries}{the 5G system's latency and reliability KPIs}
\afield{Source says}{}
\begin{saidquote}ETSI no longer mirrors v18.6.0 (only 18.13.0 and later), so the KPI clause was not opened. The performance requirements live in clause 7. If you want a page-level anchor here, the cleanest option is to lean on \texttt{3gpp38913} clause 7.4, which \emph{is} verified below.\end{saidquote}
\afield{Metadata}{[verified] \texttt{22261-i60.zip} (i = Release 18, 6.0) is dated \textbf{2022-04-01} in the 3GPP archive, consistent with \texttt{mar 2022} (SA plenary approval, posted 1 April). Re-checked 24 August 2026, date unchanged. Note the convention clash flagged in that pass: this entry follows the approval month while \texttt{3gpp38913} was moved to the posting month by this same audit, and the two situations are identical.}
\afield{Verdict}{Metadata OK, claim evidence unverified.}
\afield{Record}{Sources\_evidence.md, entry record.}
\auditentry{30}{3gpp38463}{standard, level~A}
\afield{Reference}{\fullcite{3gpp38463}}
\afield{Entry}{\href{https://www.3gpp.org/ftp/Specs/archive/38_series/38.463/38463-h00.zip}{source page}. V17.0.0, Release 17.}
\afield{Cited in}{1 place.}
\begin{citedsentence}\textbf{\small Section 2.1.}\enspace These parts communicate over standardised interfaces: NG towards the 5G core, F1 between the CU and the DU, and E1 between the CU-CP and the CU-UP~\citemark{1, 2, 30}.\end{citedsentence}
\afield{Claim it carries}{the E1 interface between CU-CP and CU-UP}
\afield{Source says}{}
\begin{saidquote}Title page of \texttt{ts\_138463v170000p.pdf}: \enquote{NG-RAN; E1 Application Protocol (E1AP) (3GPP TS 38.463 version 17.0.0 Release 17)}.\end{saidquote}
\afield{Metadata}{[verified] \texttt{38463-h00.zip} dated \textbf{2022-04-06}, matching \texttt{apr 2022}. ETSI mirror \emph{ETSI TS 138 463 V17.0.0 (2022-05)}.}
\afield{Verdict}{OK.}
\afield{Record}{Sources\_evidence.md, entry record.}
\auditentry{31}{li2014tales}{peer reviewed, level~A}
\afield{Reference}{\fullcite{li2014tales}}
\afield{Entry}{DOI \href{https://doi.org/10.1145/2670979.2670988}{\texttt{10.1145/2670979.2670988}}.}
\afield{Cited in}{2 places.}
\begin{citedsentence}\textbf{\small Section 2.2.}\enspace Fine-grained measurement studies show that, even at moderate utilisation, hardware, operating-system, and application-level effects such as background interference, scheduling, interrupt routing, and power-saving states inflate the tail far beyond what queueing models predict~\citemark{31}.\end{citedsentence}
\begin{citedsentence}\textbf{\small Appendix B.}\enspace The fourth strand is the domain and the method: the tail-latency and microsecond-scale literature~\citemark{6, 7, 31, 106}, thread-per-core structuring moving the tail in its own right~\citemark{46}, O-RAN disaggregation~\citemark{3}, control-plane latency as an active target in the core network too~\citemark{107, 108}, adoption evidence pointing at the learning curve rather than run-time cost~\citemark{80, 81, 109}, and standard guidance on repeat-based evaluation and fault injection~\citemark{33, 60}.\end{citedsentence}
\afield{Claim it carries}{hardware, OS and application-level effects inflate the tail beyond what queueing models predict}
\afield{Source says}{}
\begin{saidquote}(author copy, p. 1, title and abstract): \enquote{Tales of the Tail: \textbf{Hardware, OS, and Application-level Sources of Tail Latency}}, University of Washington.\end{saidquote}
\afield{Metadata}{[verified] Crossref matches, \texttt{pages = 9:1--9:14} is the article-number convention.}
\afield{Verdict}{OK.}
\afield{Record}{Sources\_evidence.md, entry record.}
\auditentry{32}{schroeder2006openclosed}{peer reviewed, level~A}
\afield{Reference}{\fullcite{schroeder2006openclosed}}
\afield{Entry}{\href{https://www.usenix.org/conference/nsdi-06/open-versus-closed-cautionary-tale}{source page}.}
\afield{Cited in}{2 places.}
\begin{citedsentence}\textbf{\small Section 2.2.}\enspace In a closed system a new request is triggered only by the completion of a previous one, whereas in an open system arrivals are independent of completions, and the two produce systematically different response-time behaviour on the same service process, with the closed model reporting the lower figures~\citemark{32}.\end{citedsentence}
\begin{citedsentence}\textbf{\small Section 6.2.}\enspace The send schedule is fixed and the histogram is corrected, so a stall introduces no coordinated-omission bias~\citemark{32, 64}.\end{citedsentence}
\afield{Claim it carries}{open and closed generators produce systematically different response times, the closed model reporting lower figures}
\afield{Source says}{}
\begin{saidquote}(\texttt{schroeder.pdf}, p. 1, abstract): \enquote{Workload generators may be classified as based on a \textbf{closed system model, where new job arrivals are only triggered by job completions} \ldots{} or an \textbf{open system model, where new jobs arrive independently of job completions} \ldots{} we illustrate that there is a \textbf{vast difference in behavior between open and closed models} in real-world settings.}\end{saidquote}
\afield{Metadata}{[verified] DBLP confirms NSDI 2006 and assigns \textbf{no page numbers}, the entry correctly omits them. USENIX page live.}
\afield{Verdict}{OK. This is the right primary citation for the coordinated-omission argument.}
\afield{Record}{Sources\_evidence.md, entry record.}
\auditentry{33}{georges2007rigorous}{peer reviewed, level~A}
\afield{Reference}{\fullcite{georges2007rigorous}}
\afield{Entry}{DOI \href{https://doi.org/10.1145/1297027.1297033}{\texttt{10.1145/1297027.1297033}}.}
\afield{Cited in}{3 places.}
\begin{citedsentence}\textbf{\small Section 2.2.}\enspace Sound measurement therefore requires an open-loop model with an intended schedule, correction for coordinated omission, high-dynamic-range (HDR) latency histograms that retain resolution across many orders of magnitude, and statistically rigorous reporting with repeated runs and confidence intervals rather than single point estimates~\citemark{33}.\end{citedsentence}
\begin{citedsentence}\textbf{\small Section 4.4.}\enspace This follows established guidance on statistically rigorous performance evaluation, which warns that reporting a single number invites conclusions that run-to-run variation does not support~\citemark{33}.\end{citedsentence}
\begin{citedsentence}\textbf{\small Appendix B.}\enspace The fourth strand is the domain and the method: the tail-latency and microsecond-scale literature~\citemark{6, 7, 31, 106}, thread-per-core structuring moving the tail in its own right~\citemark{46}, O-RAN disaggregation~\citemark{3}, control-plane latency as an active target in the core network too~\citemark{107, 108}, adoption evidence pointing at the learning curve rather than run-time cost~\citemark{80, 81, 109}, and standard guidance on repeat-based evaluation and fault injection~\citemark{33, 60}.\end{citedsentence}
\afield{Claim it carries}{reporting a single number invites conclusions run-to-run variation does not support}
\afield{Source says}{}
\begin{saidquote}(Ghent University copy, p. 1, abstract): \enquote{This paper shows that \textbf{prevalent methodologies can be misleading, and can even lead to incorrect conclusions}. The reason is that the data analysis is not statistically rigorous.}\end{saidquote}
\afield{Metadata}{[verified] Crossref matches, booktitle wording differs cosmetically.}
\afield{Verdict}{OK, near-verbatim support for the sentence at 1572.}
\afield{Record}{Sources\_evidence.md, entry record.}
\auditentry{34}{szekeres2013sok}{peer reviewed, level~A}
\afield{Reference}{\fullcite{szekeres2013sok}}
\afield{Entry}{DOI \href{https://doi.org/10.1109/SP.2013.13}{\texttt{10.1109/SP.2013.13}}.}
\afield{Cited in}{3 places.}
\begin{citedsentence}\textbf{\small Section 2.3.}\enspace Memory-corruption errors of this kind have been a persistent security problem for decades, and deployed mitigations remain only partial defences~\citemark{34}.\end{citedsentence}
\begin{citedsentence}\textbf{\small Section 3.2.}\enspace These are the same memory-corruption classes that decades of security research have struggled to contain in C and C++~\citemark{34}.\end{citedsentence}
\begin{citedsentence}\textbf{\small Appendix B.}\enspace The safety strand is settled the same partial way, with machine-checked foundations for the safe fragment and the aliasing discipline \texttt{unsafe} code has to respect~\citemark{11, 40, 41}, confirmation that \texttt{unsafe} is used sparingly and that memory-safety CVEs always involve it~\citemark{27, 42, 43}, auditing that finds defects at scale even so and unsafety spreading through dependency chains~\citemark{99, 100, 92, 44}, and the reports motivating the shift set against a long record of memory corruption resisting runtime mitigation~\citemark{8, 9, 17, 19, 20, 34}.\end{citedsentence}
\afield{Claim it carries}{memory corruption as a decades-old problem whose deployed mitigations are partial}
\afield{Source says}{}
\begin{saidquote}(p. 1, abstract): \enquote{This problem has existed for \textbf{more than 30 years} and a vast number of potential solutions have been proposed, yet memory corruption attacks continue to pose a serious threat. Real world exploits show that \textbf{all currently deployed protections can be defeated}.}\end{saidquote}
\afield{Metadata}{[verified] IEEE S\&P 2013, pp. 48--62. IEEE's record parses \enquote{Tao Wei} as a surname, the entry is correct.}
\afield{Verdict}{OK, this is a precise match for \enquote{deployed mitigations remain only partial defences}.}
\afield{Record}{Sources\_evidence.md, entry record.}
\auditentry{35}{bensalem2005deadlock}{peer reviewed, level~A}
\afield{Reference}{\fullcite{bensalem2005deadlock}}
\afield{Entry}{DOI \href{https://doi.org/10.1007/11678779_15}{\texttt{10.1007/11678779\_15}}.}
\afield{Cited in}{1 place.}
\begin{citedsentence}\textbf{\small Section 2.3.}\enspace Its deadlock detector inherits that limitation from the technique it implements, which infers deadlock potentials from the lock-acquisition graph of a single trace~\citemark{35}, and it ships as an experimental component covering only the Portable Operating System Interface (POSIX) mutex families~\citemark{36}.\end{citedsentence}
\afield{Source says}{}
\begin{saidquote}\enquote{Dynamic Deadlock Analysis of Multi-threaded Programs}. Saddek Bensalem, Klaus Havelund. Haifa Verification Conference 2005, LNCS 3875, pages 208--223. DOI 10.1007/11678779\_15. \emph{Note the date:} the conference is 2005, Springer's deposit records the LNCS volume as 2006. Cite the conference year and the volume, as DBLP does.\end{saidquote}
\begin{saidquote}\textbf{This is the peer-reviewed foundation for the mechanism \texttt{Ha\_Vien\_Mastersthesis\_2026.tex}:3009--3012 explains in its own voice.} The draft argues that TSan builds its lock-order graph from acquisitions that succeed, so a run in which the first inversion blocks yields one edge and never the second, and the detector stays silent at the moment the bug manifests. Bensalem and Havelund set out the technique that reasoning is about: lock-graph analysis over a single execution trace, in which a cycle signifies a deadlock \emph{potential}, and whose stated purpose is to \enquote{identify deadlock potentials even though no deadlocks occur in the examined execution trace}.\end{saidquote}
\begin{saidquote}Read together, the draft's finding stops being an oddity and becomes the expected behaviour of the technique. The analysis is designed for traces that \emph{survive}, so the trace that deadlocks is the one it has least to say about. That is a much stronger way to write the same paragraph, and it generalises the way the draft already claims it does.\end{saidquote}
\begin{saidquote}The paper is also about reducing false positives specifically, via thread labels, gate locks and segmentation, which is the complementary error mode and worth one clause for balance.\end{saidquote}
\afield{Record}{LITERATURE\_NOTES.md, section B.5 Bensalem and Havelund.}
\auditentry{36}{tsandeadlock}{tool record, level~A}
\afield{Reference}{\fullcite{tsandeadlock}}
\afield{Entry}{\href{https://github.com/google/sanitizers/wiki/ThreadSanitizerDeadlockDetector}{source page}. Project documentation. Accessed August 2026.}
\afield{Cited in}{1 place.}
\begin{citedsentence}\textbf{\small Section 2.3.}\enspace Its deadlock detector inherits that limitation from the technique it implements, which infers deadlock potentials from the lock-acquisition graph of a single trace~\citemark{35}, and it ships as an experimental component covering only the Portable Operating System Interface (POSIX) mutex families~\citemark{36}.\end{citedsentence}
\afield{Source says}{}
\begin{saidquote}Neither is peer-reviewed. Both are primary sources about the software, which is what is wanted here, and both must be labelled as such if cited.\end{saidquote}
\begin{saidquote}\begin{itemize}[leftmargin=1.2em,itemsep=1pt,topsep=1pt]
\item \textbf{ThreadSanitizerDeadlockDetector}, google/sanitizers wiki. \url{https://github.com/google/sanitizers/wiki/ThreadSanitizerDeadlockDetector}. States the mechanism in one sentence (\enquote{The deadlock detector maintains a directed graph of lock acquisitions. If a lock B is acquired while a lock A is being held by the same thread, a directed edge \texttt{A=>B} is added to the lock acquisition graph. A potential deadlock is reported when there is a cycle in the graph}), and describes the detector as \textbf{experimental}, supporting only \texttt{pthread\_mutex\_*} / \texttt{pthread\_rwlock\_*} / \texttt{pthread\_spin\_*}, and \enquote{tested only on tiny test cases}. It also confirms that detection needs \texttt{TSAN\_OPTIONS=detect\_deadlocks=1}, which is worth checking against what \texttt{faults/} actually sets.
\item \textbf{google/sanitizers issue \#1258.} Upstream record that there are two detector implementations and that \textbf{neither is actively maintained}, plus a contributor's diagnosis that version 1, the default, checks for the cycle in \texttt{DD::MutexBeforeLock} while the edge is only added in \texttt{DD::MutexAfterLock}. That is the draft's mechanism, described by the tool's own issue tracker, and it independently corroborates a paragraph the draft currently supports with its own measurements alone.
\end{itemize}\end{saidquote}
\afield{Record}{LITERATURE\_NOTES.md, section B.6 The two tool records, as tool records.}
\auditentry{37}{sutter2025profiles}{committee draft, level~A}
\afield{Reference}{\fullcite{sutter2025profiles}}
\afield{Entry}{\href{https://isocpp.org/files/papers/P3081R2.pdf}{source page}.}
\afield{Cited in}{1 place.}
\begin{citedsentence}\textbf{\small Section 2.3.}\enspace Since 2023 the standards committee has pursued \emph{profiles}: named sets of rules covering type, bounds, lifetime, and arithmetic safety that a conforming implementation must enforce wherever they are enabled~\citemark{37}.\end{citedsentence}
\afield{Claim it carries}{profiles a conforming implementation must enforce wherever enabled}
\afield{Source says}{}
\begin{saidquote}(p. 4, \S{}2.1 Strategy): \enquote{\textbf{Define standard enforced 'profiles' that a conforming C++ implementation must enforce when enabled, notably bounds, type, and lifetime.} \ldots{} Each rule must be \textbf{deterministically decidable at compile time (even if it results in injecting a check enforced at run time)} \ldots{} Rules are portable and enforced in the C++ implementation, not in a separate tool such as a static analyzer.} The four profiles appear as \S{}\S{}4--7 (type, bounds, lifetime, arithmetic).\end{saidquote}
\afield{Metadata}{[verified] Document header: \enquote{Core safety profiles for C++26, Document Number: \textbf{P3081 R2}, Date: \textbf{2025-02-04}, Reply-to: \textbf{Herb Sutter}, Audience: EWG, LEWG.}}
\afield{Verdict}{OK.}
\afield{Record}{Sources\_evidence.md, entry record.}
\auditentry{38}{baxter2024safecpp}{committee draft, level~A}
\afield{Reference}{\fullcite{baxter2024safecpp}}
\afield{Entry}{\href{https://safecpp.org/P3390R0.html}{source page}.}
\afield{Cited in}{1 place.}
\begin{citedsentence}\textbf{\small Section 2.3.}\enspace That direction was chosen over \emph{Safe C++}, a proposal for a safe subset with borrow checking and its own standard library, which is Rust's model expressed in C++~\citemark{38}.\end{citedsentence}
\afield{Claim it carries}{a proposal for a rigorously safe subset with borrow checking and an accompanying standard library}
\afield{Source says}{}
\begin{saidquote}(document body): the proposal's own compiler output shows \enquote{\textbf{borrow checking}: iterator.cxx:10:11} under \enquote{safety: during safety checking of int main() safe}, and the comparison section states \enquote{\textbf{C++ lacks traits, relocation and borrow checking}}.\end{saidquote}
\afield{Metadata}{[verified] Published at safecpp.org as WG21 P3390R0, authors Sean Baxter and Christian Mazakas. September 2024.}
\afield{Verdict}{OK.}
\afield{Record}{Sources\_evidence.md, entry record.}
\auditentry{39}{vandevoorde2026profiles}{committee draft, level~A}
\afield{Reference}{\fullcite{vandevoorde2026profiles}}
\afield{Entry}{\href{https://www.open-std.org/jtc1/sc22/wg21/docs/papers/2026/p3970r0.pdf}{source page}.}
\afield{Cited in}{3 places.}
\begin{citedsentence}\textbf{\small Section 2.3.}\enspace Profiles are not yet standardised, and in January 2026 their own authors were still calling for implementations and treating the design of the initial profiles as work still to be done~\citemark{39}.\end{citedsentence}
\begin{citedsentence}\textbf{\small Section 2.3.}\enspace The concurrency profile, the one that would bear on the defect classes catalogued in Chapter~\refmark{chapter:context}, is placed after that initial set and not within it: ``more profiles should be defined and implemented, notably a concurrency profile''~\citemark{39}.\end{citedsentence}
\begin{citedsentence}\textbf{\small Section 7.6.}\enspace That qualifies nothing measured here, because none of that work is standardised or generally implemented and the concurrency profile that would reach the boundary classes of RQ2 comes after an initial set still to be specified~\citemark{39}.\end{citedsentence}
\afield{Claim it carries}{the thesis quotes this document twice}
\afield{Source says}{}
\begin{saidquote}:\end{saidquote}
\begin{saidquote}\begin{itemize}[leftmargin=1.2em,itemsep=1pt,topsep=1pt]
\item p. 1: \enquote{Currently, \textbf{there is -- to the best of our knowledge -- a specification and an experimental implementation is being conducted in a major C++ compiler}. Also, an implementation guide is being written.} Then: \enquote{Beyond the framework implementation(s), \textbf{we need to design and precisely specify a few initial profiles}.} $\rightarrow$ supports the quotation at 779.
\item p. 2: \enquote{\textbf{After that, more profiles should be defined and implemented, notably a concurrency profile.}} $\rightarrow$ supports the quotation at 782 word for word.
\end{itemize}\end{saidquote}
\afield{Metadata}{[verified] Document header: \enquote{Doc. No. \textbf{P3970R0}, Date: \textbf{2026-01-16}}, authors David Vandevoorde, Jeff Garland, Paul E. McKenney, Roger Orr, Bjarne Stroustrup, Michael Wong, exactly the entry's list and order.}
\afield{Verdict}{OK. Both quoted fragments are verbatim.}
\afield{Record}{Sources\_evidence.md, entry record.}
\auditentry{40}{jung2021safe}{peer reviewed, level~A}
\afield{Reference}{\fullcite{jung2021safe}}
\afield{Entry}{DOI \href{https://doi.org/10.1145/3418295}{\texttt{10.1145/3418295}}.}
\afield{Cited in}{3 places.}
\begin{citedsentence}\textbf{\small Section 2.4.}\enspace The type system statically forbids mutation of shared state and confines the unavoidable exceptions to \texttt{unsafe} code kept behind safe APIs~\citemark{40}.\end{citedsentence}
\begin{citedsentence}\textbf{\small Section 2.4.}\enspace Rust is the first widely supported language to combine compile-time memory and thread safety with the absence of a garbage collector and throughput and memory use alongside C and C++~\citemark{12, 40, 45}.\end{citedsentence}
\begin{citedsentence}\textbf{\small Appendix B.}\enspace The safety strand is settled the same partial way, with machine-checked foundations for the safe fragment and the aliasing discipline \texttt{unsafe} code has to respect~\citemark{11, 40, 41}, confirmation that \texttt{unsafe} is used sparingly and that memory-safety CVEs always involve it~\citemark{27, 42, 43}, auditing that finds defects at scale even so and unsafety spreading through dependency chains~\citemark{99, 100, 92, 44}, and the reports motivating the shift set against a long record of memory corruption resisting runtime mitigation~\citemark{8, 9, 17, 19, 20, 34}.\end{citedsentence}
\afield{Claim it carries}{ownership and borrowing statically forbid mutation of shared state, exceptions confined to \texttt{unsafe} behind safe APIs}
\afield{Source says}{}
\begin{saidquote}(author copy, p. 1): \enquote{It tackles this challenge using a strong type system based on the ideas of \textbf{ownership and borrowing, which statically prohibits the mutation of shared state}}, and the following bullet on data types that \enquote{fundamentally depend on shared mutable state} being encapsulated.\end{saidquote}
\afield{Metadata}{[verified] Crossref matches on every field.}
\afield{Verdict}{OK.}
\afield{Record}{Sources\_evidence.md, entry record.}
\auditentry{41}{jung2020stacked}{peer reviewed, level~A}
\afield{Reference}{\fullcite{jung2020stacked}}
\afield{Entry}{DOI \href{https://doi.org/10.1145/3371109}{\texttt{10.1145/3371109}}.}
\afield{Cited in}{2 places.}
\begin{citedsentence}\textbf{\small Section 2.4.}\enspace The soundness of the safe fragment of the language, including the role of \texttt{unsafe} as an explicitly marked escape hatch, has been established formally by the RustBelt project~\citemark{11}, and the aliasing discipline that even \texttt{unsafe} code must respect has been given a precise operational semantics~\citemark{41}.\end{citedsentence}
\begin{citedsentence}\textbf{\small Appendix B.}\enspace The safety strand is settled the same partial way, with machine-checked foundations for the safe fragment and the aliasing discipline \texttt{unsafe} code has to respect~\citemark{11, 40, 41}, confirmation that \texttt{unsafe} is used sparingly and that memory-safety CVEs always involve it~\citemark{27, 42, 43}, auditing that finds defects at scale even so and unsafety spreading through dependency chains~\citemark{99, 100, 92, 44}, and the reports motivating the shift set against a long record of memory corruption resisting runtime mitigation~\citemark{8, 9, 17, 19, 20, 34}.\end{citedsentence}
\afield{Claim it carries}{the aliasing discipline \texttt{unsafe} code must respect, given a precise operational semantics}
\afield{Source says}{}
\begin{saidquote}(p. 1, abstract): \enquote{we propose \textbf{Stacked Borrows, an operational semantics for memory accesses in Rust}. Stacked Borrows defines an aliasing discipline and declares programs violating it to have undefined behavior.}\end{saidquote}
\afield{Metadata}{[verified] Footer: \enquote{Proc. ACM Program. Lang., Vol. 4, No. POPL, \textbf{Article 41. Publication date: January 2020}}. \texttt{pages = 41:1--41:32} correct.}
\afield{Verdict}{OK.}
\afield{Record}{Sources\_evidence.md, entry record.}
\auditentry{42}{astrauskas2020unsafe}{peer reviewed, level~B}
\afield{Reference}{\fullcite{astrauskas2020unsafe}}
\afield{Entry}{DOI \href{https://doi.org/10.1145/3428204}{\texttt{10.1145/3428204}}.}
\afield{Cited in}{2 places.}
\begin{citedsentence}\textbf{\small Section 2.4.}\enspace Empirical evidence supports the model in practice: in real codebases \texttt{unsafe} is used sparingly and in localised ways~\citemark{27, 42}, and a study of all Rust memory-safety entries in the Common Vulnerabilities and Exposures (CVE) list found that every one of them but a single compiler bug involved \texttt{unsafe} code~\citemark{43}.\end{citedsentence}
\begin{citedsentence}\textbf{\small Appendix B.}\enspace The safety strand is settled the same partial way, with machine-checked foundations for the safe fragment and the aliasing discipline \texttt{unsafe} code has to respect~\citemark{11, 40, 41}, confirmation that \texttt{unsafe} is used sparingly and that memory-safety CVEs always involve it~\citemark{27, 42, 43}, auditing that finds defects at scale even so and unsafety spreading through dependency chains~\citemark{99, 100, 92, 44}, and the reports motivating the shift set against a long record of memory corruption resisting runtime mitigation~\citemark{8, 9, 17, 19, 20, 34}.\end{citedsentence}
\afield{Claim it carries}{\texttt{unsafe} is used sparingly and in localised ways}
\afield{Source says}{}
\begin{saidquote}Crossref-deposited abstract: \enquote{an important assumption of the Rust language, which we dub the \textbf{Rust hypothesis}, is that programmers use Rust by following three main principles: \textbf{use unsafe code sparingly, make it easy to review, and hide it behind a safe abstraction}\ldots{} This paper studies empirically how unsafe code is used in practice.}\end{saidquote}
\afield{Metadata}{[verified] Crossref matches, \texttt{pages = 136:1--136:27} is the ACM article-number convention.}
\afield{Verdict}{OK.}
\afield{Record}{Sources\_evidence.md, entry record.}
\auditentry{43}{xu2021rustcve}{peer reviewed, level~A}
\afield{Reference}{\fullcite{xu2021rustcve}}
\afield{Entry}{DOI \href{https://doi.org/10.1145/3466642}{\texttt{10.1145/3466642}}.}
\afield{Cited in}{3 places.}
\begin{citedsentence}\textbf{\small Section 2.4.}\enspace Empirical evidence supports the model in practice: in real codebases \texttt{unsafe} is used sparingly and in localised ways~\citemark{27, 42}, and a study of all Rust memory-safety entries in the Common Vulnerabilities and Exposures (CVE) list found that every one of them but a single compiler bug involved \texttt{unsafe} code~\citemark{43}.\end{citedsentence}
\begin{citedsentence}\textbf{\small Section 7.2.}\enspace That fits the CVE evidence that Rust memory-safety bugs always involve \texttt{unsafe}~\citemark{43}, and the caution that unsafety can still spread through dependencies~\citemark{44}, which is why keeping the core \texttt{unsafe}-free matters.\end{citedsentence}
\begin{citedsentence}\textbf{\small Appendix B.}\enspace The safety strand is settled the same partial way, with machine-checked foundations for the safe fragment and the aliasing discipline \texttt{unsafe} code has to respect~\citemark{11, 40, 41}, confirmation that \texttt{unsafe} is used sparingly and that memory-safety CVEs always involve it~\citemark{27, 42, 43}, auditing that finds defects at scale even so and unsafety spreading through dependency chains~\citemark{99, 100, 92, 44}, and the reports motivating the shift set against a long record of memory corruption resisting runtime mitigation~\citemark{8, 9, 17, 19, 20, 34}.\end{citedsentence}
\afield{Source says}{}
\begin{saidquote}(arXiv 2003.03296, the authors' copy of the same study): p. 1, abstract. \enquote{Our analysis result shows that Rust can keep its promise that \textbf{all memory-safety bugs require unsafe code}}, and p. 2. \enquote{\textbf{all bugs in our dataset require unsafe code except one compiler bug}.}\end{saidquote}
\afield{Metadata}{[verified] TOSEM 31(1), article 3, issue dated January 2022 $\rightarrow$ \texttt{year = 2022} is right (Crossref's 2021 is online-first). Pages \texttt{3:1--3:25} correct.}
\afield{Verdict}{Entry OK. The thesis sentence at 805 (\enquote{every one of them involved \texttt{unsafe} code}) should carry the exception (fix \#5).}
\afield{Record}{Sources\_evidence.md, entry record.}
\auditentry{44}{evans2020rust}{peer reviewed, level~A}
\afield{Reference}{\fullcite{evans2020rust}}
\afield{Entry}{DOI \href{https://doi.org/10.1145/3377811.3380413}{\texttt{10.1145/3377811.3380413}}.}
\afield{Cited in}{5 places.}
\begin{citedsentence}\textbf{\small Section 2.4.}\enspace Fewer than a third of libraries use \texttt{unsafe} directly, but a much larger fraction cannot be fully statically checked because it hides somewhere in their call chain~\citemark{44}, which is a further argument for keeping the safety-critical core of a subsystem \texttt{unsafe}-free, as this thesis does.\end{citedsentence}
\begin{citedsentence}\textbf{\small Section 7.2.}\enspace That fits the CVE evidence that Rust memory-safety bugs always involve \texttt{unsafe}~\citemark{43}, and the caution that unsafety can still spread through dependencies~\citemark{44}, which is why keeping the core \texttt{unsafe}-free matters.\end{citedsentence}
\begin{citedsentence}\textbf{\small Section 7.5.}\enspace They also find that a substantial fraction of the ecosystem cannot be fully statically checked, because \texttt{unsafe} propagates through dependency chains~\citemark{44}.\end{citedsentence}
\begin{citedsentence}\textbf{\small Appendix B.}\enspace The safety strand is settled the same partial way, with machine-checked foundations for the safe fragment and the aliasing discipline \texttt{unsafe} code has to respect~\citemark{11, 40, 41}, confirmation that \texttt{unsafe} is used sparingly and that memory-safety CVEs always involve it~\citemark{27, 42, 43}, auditing that finds defects at scale even so and unsafety spreading through dependency chains~\citemark{99, 100, 92, 44}, and the reports motivating the shift set against a long record of memory corruption resisting runtime mitigation~\citemark{8, 9, 17, 19, 20, 34}.\end{citedsentence}
\begin{citedsentence}\textbf{\small Appendix G.}\enspace It keeps the dependency-propagation problem Evans et al.\ describe~\citemark{44} away from the safety-critical code, and it makes the reviewable surface of the migration a fixed, countable thing.\end{citedsentence}
\afield{Claim it carries}{a small fraction of crates use \texttt{unsafe} directly, but a much larger fraction cannot be fully statically checked}
\afield{Source says}{}
\begin{saidquote}(arXiv 2007.00752, p. 1, abstract): \enquote{software engineers use the keyword \textbf{unsafe in less than 30\% of Rust libraries, but more than half cannot be entirely statically checked} by the Rust compiler because of Unsafe Rust \textbf{hidden somewhere in a library's call chain}.}\end{saidquote}
\afield{Metadata}{[verified] Pages and year match. Crossref's booktitle is \enquote{Proceedings of the ACM/IEEE 42nd International Conference on Software Engineering}, the entry's shorter form is fine.}
\afield{Verdict}{OK, this is close to a word-for-word match with the thesis sentence at 810.}
\afield{Record}{Sources\_evidence.md, entry record.}
\auditentry{45}{bugden2022rust}{peer reviewed, level~A}
\afield{Reference}{\fullcite{bugden2022rust}}
\afield{Entry}{DOI \href{https://doi.org/10.48550/arXiv.2206.05503}{\texttt{10.48550/arXiv.2206.05503}}.}
\afield{Cited in}{2 places.}
\begin{citedsentence}\textbf{\small Section 2.4.}\enspace Rust is the first widely supported language to combine compile-time memory and thread safety with the absence of a garbage collector and throughput and memory use alongside C and C++~\citemark{12, 40, 45}.\end{citedsentence}
\begin{citedsentence}\textbf{\small Appendix B.}\enspace Compiled systems languages sit close on time and energy~\citemark{45, 93}, Rust's safety is cheap at run time~\citemark{12, 94}, and that holds where timing is tightest, in embedded hypervisors and in packet processing within a few percent of a C reference~\citemark{95, 73, 96}.\end{citedsentence}
\afield{Claim it carries}{cross-language benchmark surveys place Rust alongside C and C++}
\afield{Source says}{}
\begin{saidquote}(arXiv PDF, p. 1, abstract): \enquote{we \textbf{surveyed and analyzed recent research on Rust and research that benchmarks Rust with other available prominent programming languages} \ldots{} recent experimental results in benchmarking research show Rust's overall superiority\ldots{}}\end{saidquote}
\afield{Metadata}{[verified] arXiv 2206.05503, DOI 10.48550/arXiv.2206.05503 (DataCite), 2022.}
\afield{Verdict}{OK, but note this is still the only load-bearing arXiv-only citation, \texttt{ASSESSMENT.md} already flags replacing it as the last optional lift on criterion 2. Its use at 3782 is honest (the table marks it \enquote{General benchmarks / no / no / no}).}
\afield{Record}{Sources\_evidence.md, entry record.}
\auditentry{46}{enberg2019threadpercore}{peer reviewed, level~A}
\afield{Reference}{\fullcite{enberg2019threadpercore}}
\afield{Entry}{DOI \href{https://doi.org/10.1109/ANCS.2019.8901874}{\texttt{10.1109/ANCS.2019.8901874}}.}
\afield{Cited in}{2 places.}
\begin{citedsentence}\textbf{\small Section 2.5.}\enspace The choice also moves the quantity measured here: thread-per-core structuring changes application tail latency on its own~\citemark{46}, so the model is held fixed across both implementations instead of left to each language's idiom.\end{citedsentence}
\begin{citedsentence}\textbf{\small Appendix B.}\enspace The fourth strand is the domain and the method: the tail-latency and microsecond-scale literature~\citemark{6, 7, 31, 106}, thread-per-core structuring moving the tail in its own right~\citemark{46}, O-RAN disaggregation~\citemark{3}, control-plane latency as an active target in the core network too~\citemark{107, 108}, adoption evidence pointing at the learning curve rather than run-time cost~\citemark{80, 81, 109}, and standard guidance on repeat-based evaluation and fault injection~\citemark{33, 60}.\end{citedsentence}
\afield{Claim it carries}{the justification for holding the concurrency model fixed across both implementations}
\afield{Source says}{}
\begin{saidquote}(author copy, p. 1, abstract): \enquote{The \textbf{thread-per-core architecture has emerged to reduce these overheads, but it also has its challenges} from thread synchronization and OS interfaces. Applications can mitigate both issues with different techniques, but \textbf{their impact on application tail latency is an open question}.}\end{saidquote}
\afield{Metadata}{[verified] Crossref matches.}
\afield{Verdict}{OK.}
\afield{Record}{Sources\_evidence.md, entry record.}
\auditentry{47}{skoudlil2025asyncrust}{peer reviewed, level~A}
\afield{Reference}{\fullcite{skoudlil2025asyncrust}}
\afield{Entry}{DOI \href{https://doi.org/10.4230/LIPIcs.ECRTS.2025.1}{\texttt{10.4230/LIPIcs.ECRTS.2025.1}}.}
\afield{Cited in}{2 places.}
\begin{citedsentence}\textbf{\small Section 2.5.}\enspace Its timing behaviour belongs to the runtime, not to the language hosting it: a real-time study of asynchronous Rust finds one application meeting or missing its deadlines depending on the runtime and thread-priority mapping it is given, and the configuration that meets them is not the idiomatic work-stealing one~\citemark{47}.\end{citedsentence}
\begin{citedsentence}\textbf{\small Section 7.5.}\enspace v{S}koudlil et al.\ report that Tokio is ``designed to maximize throughput rather than time determinism'', and that restructuring onto prioritised threads yields response times that match classical analysis~\citemark{47}.\end{citedsentence}
\afield{Claim it carries}{the configuration that meets deadlines is not the idiomatic work-stealing one, tail behaviour is a property of the runtime}
\afield{Source says}{}
\begin{saidquote}:\end{saidquote}
\begin{saidquote}\begin{itemize}[leftmargin=1.2em,itemsep=1pt,topsep=1pt]
\item p. 5, \S{}2.2.2: \enquote{\textbf{Tokio} is a popular multi-threaded runtime. It seems to be \textbf{designed to maximize throughput rather than time determinism}. Their scheduling policies are complex and difficult to understand due to the use of abstraction layers.}
\item p. 15: \enquote{\textbf{futures-rt and rclcpp-rt are the only variants where all callbacks manage to complete before the deadlines} given by the publication periods. All other implemented variants fail to meet shorter deadlines for topics 1--3. This is expected for the default configurations of Rust async runtimes.}
\item p. 1, abstract: \enquote{the experimental results indicate that our proposed structure \textbf{achieves bounded response times for time-critical tasks}.}
\end{itemize}\end{saidquote}
\afield{Metadata}{[verified] The LIPIcs PDF's running heads confirm \textbf{1:1--1:21}, volume 335, ECRTS 2025, and the three authors as entered.}
\afield{Verdict}{OK, supports exactly the rebuttal the thesis builds on it.}
\afield{Record}{Sources\_evidence.md, entry record.}
\auditentry{48}{rivera2021galeed}{peer reviewed, level~B}
\afield{Reference}{\fullcite{rivera2021galeed}}
\afield{Entry}{DOI \href{https://doi.org/10.1145/3485832.3485903}{\texttt{10.1145/3485832.3485903}}.}
\afield{Cited in}{6 places.}
\begin{citedsentence}\textbf{\small Section 2.6.}\enspace This has been shown experimentally, and the proposed remedies isolate the two heaps in hardware or in the runtime instead of trusting the discipline of the interface~\citemark{48, 49}.\end{citedsentence}
\begin{citedsentence}\textbf{\small Section 7.2.}\enspace That is why the remedies proposed in the literature isolate the two heaps with hardware support instead of relying on interface discipline~\citemark{48, 49}, and the defect classes that arise at such boundaries have been catalogued in their own right~\citemark{50}.\end{citedsentence}
\begin{citedsentence}\textbf{\small Section 7.6.}\enspace RQ4 measures the seam and not the safety of the composed system: host and component share an address space~\citemark{48, 49}.\end{citedsentence}
\begin{citedsentence}\textbf{\small Chapter 8.}\enspace It can also compare Rust async runtimes against the thread-per-shard model, which would separate the tail divergence Arvidsson and Kaller attribute to \texttt{async}/\texttt{await}~\citemark{56}, measure in-process isolation of the two heaps at the seam~\citemark{48, 49}, and track defect density and review effort as Rust is adopted in a production RAN codebase.\end{citedsentence}
\begin{citedsentence}\textbf{\small Appendix B.}\enspace The interoperability strand qualifies the adoption argument, and Section~\refmark{sec:interop} states its central limit: confining \texttt{unsafe} to a seam does not protect the Rust side, because foreign code in the same address space can corrupt Rust-owned data anyway~\citemark{48, 49}.\end{citedsentence}
\begin{citedsentence}\textbf{\small Appendix G.}\enspace While the host remains C++ the two share a heap, so a team needing the stronger property has to look at in-process isolation~\citemark{48, 49} and not at interface discipline alone.\end{citedsentence}
\afield{Claim it carries}{proposed remedies isolate the two heaps rather than relying on interface discipline}
\afield{Source says}{}
\begin{saidquote}(deposited abstract): \enquote{these guarantees are violated when safe Rust interacts with unsafe code, most notably code written in other programming languages ... \textbf{a dangling pointer in a linked C/C++ library can access and overwrite memory allocated to Rust even when the Rust code is fully safe}.}\end{saidquote}
\afield{Verdict}{OK, and it is the sharpest statement in the bibliography of the limit \S{}7.6 draws around RQ4: confining \texttt{unsafe} bounds what must be reviewed, not what can be corrupted.}
\afield{Record}{Sources\_evidence.md, entry record.}
\auditentry{49}{kirth2022pkrusafe}{peer reviewed, level~B}
\afield{Reference}{\fullcite{kirth2022pkrusafe}}
\afield{Entry}{DOI \href{https://doi.org/10.1145/3492321.3519582}{\texttt{10.1145/3492321.3519582}}.}
\afield{Cited in}{6 places.}
\begin{citedsentence}\textbf{\small Section 2.6.}\enspace This has been shown experimentally, and the proposed remedies isolate the two heaps in hardware or in the runtime instead of trusting the discipline of the interface~\citemark{48, 49}.\end{citedsentence}
\begin{citedsentence}\textbf{\small Section 7.2.}\enspace That is why the remedies proposed in the literature isolate the two heaps with hardware support instead of relying on interface discipline~\citemark{48, 49}, and the defect classes that arise at such boundaries have been catalogued in their own right~\citemark{50}.\end{citedsentence}
\begin{citedsentence}\textbf{\small Section 7.6.}\enspace RQ4 measures the seam and not the safety of the composed system: host and component share an address space~\citemark{48, 49}.\end{citedsentence}
\begin{citedsentence}\textbf{\small Chapter 8.}\enspace It can also compare Rust async runtimes against the thread-per-shard model, which would separate the tail divergence Arvidsson and Kaller attribute to \texttt{async}/\texttt{await}~\citemark{56}, measure in-process isolation of the two heaps at the seam~\citemark{48, 49}, and track defect density and review effort as Rust is adopted in a production RAN codebase.\end{citedsentence}
\begin{citedsentence}\textbf{\small Appendix B.}\enspace The interoperability strand qualifies the adoption argument, and Section~\refmark{sec:interop} states its central limit: confining \texttt{unsafe} to a seam does not protect the Rust side, because foreign code in the same address space can corrupt Rust-owned data anyway~\citemark{48, 49}.\end{citedsentence}
\begin{citedsentence}\textbf{\small Appendix G.}\enspace While the host remains C++ the two share a heap, so a team needing the stronger property has to look at in-process isolation~\citemark{48, 49} and not at interface discipline alone.\end{citedsentence}
\afield{Claim it carries}{proposed remedies isolate the two heaps rather than relying on interface discipline}
\afield{Source says}{}
\begin{saidquote}(deposited abstract): \enquote{many memory-safe languages still have to \textbf{interface with unsafe code} to some extent, which opens up the possibility for attackers to exploit memory-corruption vulnerabilities in the unsafe part of the system and \textbf{subvert the safety guarantees provided by the memory-safe language}.} The mechanism is in the title: \enquote{Automatically Locking Down the Heap Between Safe and Unsafe Languages}.\end{saidquote}
\afield{Verdict}{OK.}
\afield{Record}{Sources\_evidence.md, entry record.}
\auditentry{50}{li2022crosslang}{peer reviewed, level~C}
\afield{Reference}{\fullcite{li2022crosslang}}
\afield{Entry}{DOI \href{https://doi.org/10.1007/978-3-031-17143-7_33}{\texttt{10.1007/978-3-031-17143-7\_33}}.}
\afield{Cited in}{3 places.}
\begin{citedsentence}\textbf{\small Section 2.6.}\enspace The defect classes specific to Rust/C boundaries have been catalogued separately from those of either language~\citemark{50}, and the binding generators themselves carry bugs~\citemark{51}.\end{citedsentence}
\begin{citedsentence}\textbf{\small Section 7.2.}\enspace That is why the remedies proposed in the literature isolate the two heaps with hardware support instead of relying on interface discipline~\citemark{48, 49}, and the defect classes that arise at such boundaries have been catalogued in their own right~\citemark{50}.\end{citedsentence}
\begin{citedsentence}\textbf{\small Appendix B.}\enspace The seam carries its own defect classes and its own tooling bugs, lighter crossing models exist for embedded use, aliasing patterns expressible in C but not in safe Rust bound translation in principle, interoperation is the central open problem in embedded Rust, \texttt{unsafe} concentrates at exactly these boundaries, and the two largest documented adoptions into existing C and C++ report the resulting friction empirically~\citemark{50, 51, 101, 102, 103, 104, 82, 83}.\end{citedsentence}
\afield{Claim it carries}{defect classes arising specifically at Rust/C boundaries have been catalogued in their own right}
\afield{Source says}{}
\begin{saidquote}Springer paywall, title-level, and the thesis uses it only for the existence of the category.\end{saidquote}
\afield{Metadata}{[verified] DBLP and Crossref both give pages \textbf{680--700}.}
\afield{Verdict}{OK.}
\afield{Record}{Sources\_evidence.md, entry record.}
\auditentry{51}{cai2026interop}{peer reviewed, level~B}
\afield{Reference}{\fullcite{cai2026interop}}
\afield{Entry}{DOI \href{https://doi.org/10.1145/3795532}{\texttt{10.1145/3795532}}.}
\afield{Cited in}{3 places.}
\begin{citedsentence}\textbf{\small Section 2.6.}\enspace The defect classes specific to Rust/C boundaries have been catalogued separately from those of either language~\citemark{50}, and the binding generators themselves carry bugs~\citemark{51}.\end{citedsentence}
\begin{citedsentence}\textbf{\small Appendix B.}\enspace The seam carries its own defect classes and its own tooling bugs, lighter crossing models exist for embedded use, aliasing patterns expressible in C but not in safe Rust bound translation in principle, interoperation is the central open problem in embedded Rust, \texttt{unsafe} concentrates at exactly these boundaries, and the two largest documented adoptions into existing C and C++ report the resulting friction empirically~\citemark{50, 51, 101, 102, 103, 104, 82, 83}.\end{citedsentence}
\begin{citedsentence}\textbf{\small Appendix G.}\enspace And the generated part of a seam is not above suspicion, because binding generators have their own defect history~\citemark{51}.\end{citedsentence}
\afield{Claim it carries}{binding generators carry their own defect history}
\afield{Source says}{}
\begin{saidquote}Crossref-deposited abstract: \enquote{we \textbf{mined 320 bugs in the popular interop tools, i.e., Bindgen, Cbindgen, and CXX}. We observe that the main obstacle is not tool crashes or performance issues, but the \textbf{inability to generate correct and functional code} \ldots{} Memory layout mismatch is the most common cause with Bindgen.}\end{saidquote}
\afield{Metadata}{[verified] Crossref matches authors, journal, year, DOI. Volume/issue not yet assigned, so the entry correctly omits them.}
\afield{Verdict}{OK, supports \enquote{review it as code rather than trusting it as tooling output} at 3558.}
\afield{Record}{Sources\_evidence.md, entry record.}
\auditentry{52}{pearce2026rustmc}{peer reviewed, level~A}
\afield{Reference}{\fullcite{pearce2026rustmc}}
\afield{Entry}{DOI \href{https://doi.org/10.1007/978-3-032-28187-6_14}{\texttt{10.1007/978-3-032-28187-6\_14}}.}
\afield{Cited in}{2 places.}
\begin{citedsentence}\textbf{\small Section 2.6.}\enspace A stateless model checker for Rust now explores the state space of a program together with its foreign dependencies, which puts data races arising from the interaction of the two languages within reach of a checker even though they sit outside what the borrow checker sees~\citemark{52}.\end{citedsentence}
\begin{citedsentence}\textbf{\small Chapter 8.}\enspace Beyond that, the framework can take captured anonymised traces in place of the synthetic workload, further control-plane subsystems or the data-heavy L2/PHY path, and a whole-program model checker covering foreign dependencies~\citemark{52}.\end{citedsentence}
\afield{Source says}{}
\begin{saidquote}\enquote{RustMC: Automated Verification of Real-World Concurrent Rust}. Oliver Pearce, Julien Lange, Dan O'Keeffe (Royal Holloway, University of London). \emph{Formal Techniques for Distributed Objects, Components, and Systems} (FORTE 2026), Urbino, 8--12 June 2026, LNCS, pages 235--255. DOI 10.1007/978-3-032-28187-6\_14. (A February 2025 preprint exists as arXiv 2502.06293 under the earlier title \enquote{RustMC: Extending the GenMC stateless model checker to Rust}. Cite the FORTE version.)\end{saidquote}
\begin{saidquote}\textbf{This is the best find of the revised sweep and it lands on two of the thesis's claims at once.} A stateless model checker for concurrent Rust built on GenMC, and, from the abstract, it verifies \enquote{Rust code and any foreign function dependencies}, finding \enquote{concurrency bugs in unsafe Rust code, FFI calls to C/C++, and atomic operations}. The authors state it is \enquote{the only available framework capable of exhaustively exploring the state space of Rust programs and their dependencies}.\end{saidquote}
\afield{Record}{LITERATURE\_NOTES.md, section The anchor: Pearce, Lange and O'Keeffe, FORTE 2026.}
\auditentry{53}{karlback2021ran}{thesis, level~A}
\afield{Reference}{\fullcite{karlback2021ran}}
\afield{Entry}{\href{https://urn.kb.se/resolve?urn=urn:nbn:se:liu:diva-176307}{source page}. Carried out at Ericsson AB. Report LIU-IDA/LITH-EX-A--2021/023--SE.}
\afield{Cited in}{7 places.}
\begin{citedsentence}\textbf{\small Section 2.7.}\enspace One weighs runtime cost against development efficiency across C++, Rust, and Go~\citemark{53}.\end{citedsentence}
\begin{citedsentence}\textbf{\small Section 2.7.}\enspace Three of them vary the language together with the runtime and the concurrency model, so a measured difference cannot be attributed to the language rather than to the framework built on it~\citemark{53, 55, 56}.\end{citedsentence}
\begin{citedsentence}\textbf{\small Section 7.2.}\enspace They call the pattern surprising and attribute it tentatively to the initial stack size of the green threads~\citemark{53}.\end{citedsentence}
\begin{citedsentence}\textbf{\small Section 7.5.}\enspace They ``aimed for providing the best solution in each language, instead of trying to make the code as identical as possible between the languages''~\citemark{53}.\end{citedsentence}
\begin{citedsentence}\textbf{\small Section 7.5, table cell.}\enspace Or\"o, 2021~\citemark{53}\end{citedsentence}
\begin{citedsentence}\textbf{\small Section 7.5.}\enspace None is unoccupied on its own, as the matched implementations~\citemark{73, 84, 85, 86, 87, 91, 57}, the percentile sweep stopping at a latency requirement~\citemark{56} and the effort measurements~\citemark{53, 80} all show.\end{citedsentence}
\begin{citedsentence}\textbf{\small Section 7.6.}\enspace Peak resident set size sits between the two groups, with the sign found here reported by two neighbouring theses and reversed by a third~\citemark{54, 79, 53}.\end{citedsentence}
\afield{Source says}{}
\begin{saidquote}, claim by claim (page numbers are PDF pages of \texttt{FULLTEXT01.pdf}):\end{saidquote}
\begin{saidquote}\begin{itemize}[leftmargin=1.2em,itemsep=1pt,topsep=1pt]
\item \emph{\enquote{conclude that C++ remains the choice where performance is critical}} (\texttt{Ha\_Vien\_Mastersthesis\_2026.tex}), supported, p. 79: \enquote{It is quite clear that \textbf{C++ is the superior language if low CPU usage is important}}, alongside the mixed processing-time results on the same page.
\item \emph{Memory: \enquote{find Rust the lowest of the three languages and C++ the highest, with the gap widening as concurrency rises}} (\texttt{Ha\_Vien\_Mastersthesis\_2026.tex}), supported, p. 80: \enquote{\textbf{We see that C++ has the highest memory usage, especially when looking at larger amount of sequences, and Rust has the least memory usage.}}
\item \emph{\enquote{a pattern they flag as unexplained}} (\texttt{Ha\_Vien\_Mastersthesis\_2026.tex}). \textbf{overstated.} Same passage, p. 80: \enquote{\textbf{What is surprising with these measurements is how little memory Rust uses. The main theory for this is that the initial stack size for the Rust green threads could be the smallest}, or that they adapt dynamically fast[er]\ldots{}}. They call it surprising and offer a theory, and their theory is the same green-thread mechanism the thesis then argues for, so saying so strengthens the thesis (fix \#7).
\item \textbf{The quotation at \texttt{Ha\_Vien\_Mastersthesis\_2026.tex–3697}, not found in the source.} The thesis prints: \enquote{an effort \ldots{} to utilize the different strengths of each language, and stay away from using exactly the same implementation}. The source, \S{}6.2.5 Implementation, p. 84 of the PDF (\textbf{printed page 74}), says: \enquote{This way we hopefully utilized the languages built-in support for the different parts of the code and gave each language an as representative solution as possible and \textbf{stayed away from using the exact same solution in each language}. We aimed for providing the best solution in each language, instead of trying to make the code as identical as possible between the languages.} And later on the same page: \enquote{\ldots{}a real world scenario where a programmer most likely would, and should, \textbf{utilize the strengths of each language} when implementing the solution.} The phrase \enquote{an effort was made} does occur on that page, but in a different sentence about a different subject: \enquote{\textbf{Although an effort was made to minimize this}, multiple implementation details could probably have been done better}, where \enquote{this} is their own inexperience skewing results in C++'s favour. \textbf{Suggested replacement} (accurate, and makes the same point): > The authors record that they deliberately did not write the same program three times: they > \enquote{aimed for providing the best solution in each language, instead of trying to make the code as > identical as possible between the languages}, and \enquote{stayed away from using the exact same > solution in each language}~\textbackslash{}cite\{karlback2021ran\}.
\end{itemize}\end{saidquote}
\afield{Metadata}{[verified] mostly. The DiVA record (\texttt{urn:nbn:se:liu:diva-176307} $\rightarrow$ \texttt{diva2:1563049}) is the right thesis: cover page reads \enquote{Linköping University, Department of Computer and Information Science, Master's thesis, 30 ECTS, Datateknik, \textbf{2021}, \textbf{LIU-IDA/LITH-EX-A--2021/023--SE}, \emph{Holistic view on alternative programming languages for Radio Access Network applications in cloud and embedded deployments}}. One nuance: the cover lists \textbf{Anton Orö before Rasmus Karlbäck}, the entry and the thesis text use \enquote{Karlbäck and Orö}. DiVA's own citation order is alphabetical, so this is defensible, but be aware the title page reads the other way.}
\afield{Verdict}{\textbf{FIX}, the substance is right, the quotation marks are not (fix \#1).}
\afield{Record}{Sources\_evidence.md, entry record.}
\auditentry{54}{olsson2023ranmgmt}{thesis, level~A}
\afield{Reference}{\fullcite{olsson2023ranmgmt}}
\afield{Entry}{\href{https://urn.kb.se/resolve?urn=urn:nbn:se:kth:diva-325822}{source page}. Carried out at Ericsson AB.}
\afield{Cited in}{6 places.}
\begin{citedsentence}\textbf{\small Section 2.7.}\enspace One asks whether Rust could replace C in RAN node management~\citemark{54}.\end{citedsentence}
\begin{citedsentence}\textbf{\small Section 2.7.}\enspace The remaining two hold the design closer to fixed but report execution time and peak memory rather than tail behaviour under rising load~\citemark{54, 57}, and the closest of them, Emberg and Wadsten, measures a stateless transformation rather than a stateful, concurrent subsystem.\end{citedsentence}
\begin{citedsentence}\textbf{\small Section 7.2.}\enspace A second Ericsson-hosted study of RAN node management software found Rust's peak memory usage generally higher than C's~\citemark{54}, and a matched pair of user-space device drivers outside telecom put the Rust version's footprint about 29\% above the C version's~\citemark{79}.\end{citedsentence}
\begin{citedsentence}\textbf{\small Section 7.5.}\enspace Olsson's later management-software comparison agrees with this study on two of three findings, including that resource result, and does not report percentiles or drive the software past capacity~\citemark{54}.\end{citedsentence}
\begin{citedsentence}\textbf{\small Section 7.5, table cell.}\enspace Olsson, 2023~\citemark{54}\end{citedsentence}
\begin{citedsentence}\textbf{\small Section 7.6.}\enspace Peak resident set size sits between the two groups, with the sign found here reported by two neighbouring theses and reversed by a third~\citemark{54, 79, 53}.\end{citedsentence}
\afield{Claim it carries}{Rust's peak memory usage generally higher than C's, and execution speed on par with or better than C's}
\afield{Source says}{}
\begin{saidquote}(\texttt{FULLTEXT01.pdf}, p. 4, abstract): \enquote{The results show that the execution speed of Rust applications are on par or better than their C counterpart, while the peak memory usage generally is higher in Rust. It was also shown that significant improvements in safety can be achieved, even in cases where the application is fairly straight forward.}\end{saidquote}
\afield{Metadata}{[verified] DiVA \texttt{diva2:1751095}, full text \texttt{FULLTEXT01.pdf} retrieved.}
\afield{Verdict}{OK. Both halves of the claim are the source's own summary sentence, and this is the neighbouring study whose peak-memory sign agrees with this thesis.}
\afield{Record}{Sources\_evidence.md, entry record.}
\auditentry{55}{olsson2023asyncpackets}{thesis, level~A}
\afield{Reference}{\fullcite{olsson2023asyncpackets}}
\afield{Entry}{\href{https://urn.kb.se/resolve?urn=urn:nbn:se:liu:diva-197029}{source page}. Carried out at Ericsson AB. Report LIU-IDA/LITH-EX-A--23/052--SE.}
\afield{Cited in}{4 places.}
\begin{citedsentence}\textbf{\small Section 2.7.}\enspace One builds a thread-per-core packet processor on asynchronous Rust and measures it against the Data Plane Development Kit (DPDK) application it mirrors~\citemark{55}.\end{citedsentence}
\begin{citedsentence}\textbf{\small Section 2.7.}\enspace Three of them vary the language together with the runtime and the concurrency model, so a measured difference cannot be attributed to the language rather than to the framework built on it~\citemark{53, 55, 56}.\end{citedsentence}
\begin{citedsentence}\textbf{\small Section 7.5.}\enspace Olsson and Stenstr\"om's asynchronous Rust at a fifth to a half of DPDK's throughput is the sharpest illustration of why the model has to be held fixed: an effect of two to five times is not the scale at which code generation operates~\citemark{55}.\end{citedsentence}
\begin{citedsentence}\textbf{\small Section 7.5, table cell.}\enspace Stenstr\"om, 2023~\citemark{55}\end{citedsentence}
\afield{Claim it carries}{asynchronous Rust at a fifth to a half of the throughput of the DPDK application it mirrors, and judged not currently viable}
\afield{Source says}{}
\begin{saidquote}(\texttt{FULLTEXT02.pdf}, p. 3, abstract): \enquote{The results show that the Rust application with the decentralized scheduler has a throughput that is between 20\% and 50\% of its DPDK counterpart, while its latency is far worse.} And: \enquote{the result of this thesis indicates that using Rust async for packet processing systems is currently not a viable approach.}\end{saidquote}
\afield{Metadata}{[verified] DiVA \texttt{diva2:1789208}, full text \texttt{FULLTEXT02.pdf} retrieved.}
\afield{Verdict}{OK, with the scope note above: the figure is the decentralized arm.}
\afield{Record}{Sources\_evidence.md, entry record.}
\auditentry{56}{arvidsson2025async}{thesis, level~A}
\afield{Reference}{\fullcite{arvidsson2025async}}
\afield{Entry}{\href{https://urn.kb.se/resolve?urn=urn:nbn:se:liu:diva-214746}{source page}. Carried out at Ericsson AB. Report LIU-IDA/LITH-EX-A--25/055--SE.}
\afield{Cited in}{9 places.}
\begin{citedsentence}\textbf{\small Section 2.7.}\enspace One measures median and 99th-percentile latency for control-plane procedures across five asynchronous programming frameworks~\citemark{56}.\end{citedsentence}
\begin{citedsentence}\textbf{\small Section 2.7.}\enspace Three of them vary the language together with the runtime and the concurrency model, so a measured difference cannot be attributed to the language rather than to the framework built on it~\citemark{53, 55, 56}.\end{citedsentence}
\begin{citedsentence}\textbf{\small Section 6.4.2.}\enspace The closest prior work reports a comparable degradation past a threshold request rate and attributes it to the asynchronous runtime~\citemark{56}.\end{citedsentence}
\begin{citedsentence}\textbf{\small Section 7.2.}\enspace They trace it to stackless \texttt{async}/\texttt{await}, which permits a context switch immediately before an I/O operation begins~\citemark{56}.\end{citedsentence}
\begin{citedsentence}\textbf{\small Section 7.5.}\enspace Across five asynchronous frameworks on 5G control-plane procedures, Boost.Fibers shows the lowest latency at every load, and while the Rust runtimes track it at the median, their 99th percentile rises more steeply beyond a certain request rate~\citemark{56}.\end{citedsentence}
\begin{citedsentence}\textbf{\small Section 7.5, table cell.}\enspace Kaller, 2025~\citemark{56}\end{citedsentence}
\begin{citedsentence}\textbf{\small Section 7.5.}\enspace None is unoccupied on its own, as the matched implementations~\citemark{73, 84, 85, 86, 87, 91, 57}, the percentile sweep stopping at a latency requirement~\citemark{56} and the effort measurements~\citemark{53, 80} all show.\end{citedsentence}
\begin{citedsentence}\textbf{\small Section 7.6.}\enspace No study located by the protocol of Appendix~\refmark{appendix:search} reports a load-dependent reversal between the two languages, and the one reporting the same shape attributes it to a runtime this design excludes~\citemark{56}.\end{citedsentence}
\begin{citedsentence}\textbf{\small Chapter 8.}\enspace It can also compare Rust async runtimes against the thread-per-shard model, which would separate the tail divergence Arvidsson and Kaller attribute to \texttt{async}/\texttt{await}~\citemark{56}, measure in-process isolation of the two heaps at the seam~\citemark{48, 49}, and track defect density and review effort as Rust is adopted in a production RAN codebase.\end{citedsentence}
\afield{Source says}{}
\begin{saidquote}, claim by claim (PDF pages):\end{saidquote}
\begin{saidquote}\begin{itemize}[leftmargin=1.2em,itemsep=1pt,topsep=1pt]
\item \emph{Frameworks compared and Boost.Fibers fastest}, p. 3, abstract: \enquote{\textbf{C++ with Boost.Fibers, Rust with both Tokio and Monoio, Go with Goroutines, and Erlang} \ldots{} The research reveals that \textbf{C++ with Boost.Fibers achieves the highest performance, exhibiting the lowest latency and highest throughput, while Rust frameworks present a competitive alternative.}}
\item \emph{\enquote{Boost.Fibers shows the lowest latency at every load}} (\texttt{Ha\_Vien\_Mastersthesis\_2026.tex}), p. 50: \enquote{\textbf{it is evident that Boost.Fibers consistently exhibits the lowest latency}} across the load sweep, same page: \enquote{as the load, sequence frequency, increases, latency also rises}.
\item \emph{\enquote{the Rust runtimes track it at the median, their 99th percentile rises more steeply beyond a certain request rate}}, supported by \S{}9.1.2 (p. 77) plus the load-sweep figures: \enquote{C++ with Boost.Fibers displays a \textbf{minimal difference between} [median and p99]\ldots{} In contrast, \textbf{both Rust frameworks exhibit a much greater disparity between their median and 99th percentile latencies}}, with the effect \enquote{especially under high-load conditions} (p. 78). The \emph{shape} claim (rate threshold) is read off their figures rather than a sentence, accurate, but if challenged, point to the figures.
\item \emph{\enquote{the authors attribute the effect to stackless \texttt{async}/\texttt{await} permitting a context switch immediately before an I/O operation is initiated}} (\texttt{Ha\_Vien\_Mastersthesis\_2026.tex–3715}). \textbf{exact}, \S{}9.1.2, pp. 77--78: \enquote{\textbf{Rust utilizes stackless async/await} for user-level threading, whereas C++ with Boost.Fibers employs stackful fibers \ldots{} \textbf{both Rust frameworks may perform a context switch at any await point} \ldots{} \textbf{Scenario 2 can occur in Rust, where a context switch happens just before an I/O operation is initiated} \ldots{} This behavior leads to increased latency, \textbf{particularly noticeable in the 99th percentile latency measurements}.}
\end{itemize}\end{saidquote}
\afield{Metadata}{[verified] The DiVA record (\texttt{urn:nbn:se:liu:diva-214746} $\rightarrow$ \texttt{diva2:1970356}) matches on every field: \enquote{Master's thesis, 30 ECTS, Computer Science and Engineering, \textbf{2025}, \textbf{LIU-IDA/LITH-EX-A--25/055--SE}, \emph{Evaluating Asynchronous Programming Frameworks for 5G Control Plane Procedures}, \textbf{Albin Arvidsson, Martin Kaller}}, external supervisor at Ericsson.}
\afield{Verdict}{OK, the strongest rival result in the thesis is represented fairly and attributed in the source's own terms.}
\afield{Record}{Sources\_evidence.md, entry record.}
\auditentry{57}{emberg2023asn1}{thesis, level~A}
\afield{Reference}{\fullcite{emberg2023asn1}}
\afield{Entry}{\href{https://urn.kb.se/resolve?urn=urn:nbn:se:liu:diva-195552}{source page}. Carried out at Ericsson, Link{\"o}ping. Report LIU-IDA/LITH-EX-A--23/078--SE.}
\afield{Cited in}{6 places.}
\begin{citedsentence}\textbf{\small Section 2.7.}\enspace The last builds matched C++ and Rust implementations of an Abstract Syntax Notation One (ASN.1) decompiler and compares them statistically~\citemark{57}.\end{citedsentence}
\begin{citedsentence}\textbf{\small Section 2.7.}\enspace The remaining two hold the design closer to fixed but report execution time and peak memory rather than tail behaviour under rising load~\citemark{54, 57}, and the closest of them, Emberg and Wadsten, measures a stateless transformation rather than a stateful, concurrent subsystem.\end{citedsentence}
\begin{citedsentence}\textbf{\small Section 7.5.}\enspace Emberg and Wadsten are the closest of the five in method: a matched pair, an FFI seam, a statistical comparison and an effort measurement, two years earlier~\citemark{57}.\end{citedsentence}
\begin{citedsentence}\textbf{\small Section 7.5, table cell.}\enspace Wadsten, 2023~\citemark{57}\end{citedsentence}
\begin{citedsentence}\textbf{\small Section 7.5.}\enspace None is unoccupied on its own, as the matched implementations~\citemark{73, 84, 85, 86, 87, 91, 57}, the percentile sweep stopping at a latency requirement~\citemark{56} and the effort measurements~\citemark{53, 80} all show.\end{citedsentence}
\begin{citedsentence}\textbf{\small Section 7.6.}\enspace That the threat is real and not merely conceivable is documented in the same industry, by the two benchmarking sessions of Emberg and Wadsten that reversed their own verdict~\citemark{57}.\end{citedsentence}
\afield{Claim it carries}{matched C++ and Rust ASN.1 decompilers reaching third-party C++ through a C wrapper, compared statistically, whose two benchmarking sessions reversed their own verdict}
\afield{Source says}{}
\begin{saidquote}(\texttt{FULLTEXT01.pdf}, p. 3, abstract): \enquote{Two ASN.1 decompiler implementations were made, one in C++ and one in Rust. Both decompilers invoked third-party C++ functions: in Rust, this was done by first wrapping them in C. The performance of both decompilers was measured during two benchmarking sessions, between which additional optimization was performed on both. The results were then compared statistically using the Mann-Whitney rank sum test method. The first benchmarking session showed similar performance for both decompilers ... In the second benchmarking session, the Mann-Whitney test showed that the Rust decompiler performed worse, with a mean difference of 85.19-91.36\%.} The C-wrapper route is drawn in Figure 5.2, p. 42.\end{saidquote}
\afield{Metadata}{[verified] DiVA \texttt{diva2:1772502}, full text \texttt{FULLTEXT01.pdf} retrieved.}
\afield{Verdict}{OK. Every clause the thesis attributes to this source is in its own abstract, and the 85 to 91 per cent range recorded in \texttt{LITERATURE\_NOTES.md} is exact at 85.19 to 91.36.}
\afield{Record}{Sources\_evidence.md, entry record.}
\auditentry{58}{gjomemo2016mesc}{peer reviewed, level~B}
\afield{Reference}{\fullcite{gjomemo2016mesc}}
\afield{Entry}{DOI \href{https://doi.org/10.1109/QRS.2016.44}{\texttt{10.1109/QRS.2016.44}}. Two of the six authors are affiliated with Nokia.}
\afield{Cited in}{1 place.}
\begin{citedsentence}\textbf{\small Section 2.7.}\enspace The published record on the other side is emptier: there is no study of Rust among the indexed publications of the vendor hosting this work, and the nearest peer-reviewed position from that side argues for making memory-error sanitisation practical rather than for replacing it with a compile-time guarantee~\citemark{58}, which is the strategy this thesis tests against.\end{citedsentence}
\afield{Claim it carries}{argues for making memory-error sanitisation compilers practical on the premise that certifying compilation is not}
\afield{Source says}{}
\begin{saidquote}(publisher-deposited abstract via OpenAlex, no open copy exists): \enquote{There has been considerable interest in having a compiler to ensure the safety of compiled code either through static verification or through instrumented runtime checks. \textbf{While certifying compilation has shown much promise, it has not been practical, leaving code instrumentation as the next best strategy} for compilation. We term such compilers Memory Error Sanitization Compilers (MESCs) ... The benchmarks demonstrate an improvement in runtime performance that \textbf{makes incorporation of runtime checks a viable option for defenses}.}\end{saidquote}
\afield{Metadata}{[verified] Crossref: six authors as in the entry, pages 323-334, August 2016. \textbf{Two Nokia affiliations confirmed} through OpenAlex institution data: Edmund W. Ballou and Kedar S. Namjoshi are both Nokia (United States). The other four are University of Illinois Chicago and University of Dayton. That settles the plan's \enquote{co-authored by two Nokia researchers}.}
\afield{Verdict}{OK, and an exact match on both halves. The premise clause and the practicality clause are both the source's own sentences. Level B only because the full text is paywalled and the abstract already settles it.}
\afield{Record}{Sources\_evidence.md, entry record.}
\auditentry{59}{michaelides2026ranlatency}{peer reviewed, level~A}
\afield{Reference}{\fullcite{michaelides2026ranlatency}}
\afield{Entry}{DOI \href{https://doi.org/10.1145/3765613.3797453}{\texttt{10.1145/3765613.3797453}}.}
\afield{Cited in}{1 place.}
\begin{citedsentence}\textbf{\small Section 3.3.}\enspace A measurement study of a disaggregated 5G RAN times the CU-CP-to-DU UE Context Setup procedure over F1-C, and the CU-CP-to-CU-UP bearer context setup over E1, under different interface-protection schemes, and finds the choice worth hundreds of microseconds: on F1-C, Internet Protocol Security (IPsec) adds on the order of $120\,\mu$s on average to UE context setup, against roughly $430\,\mu$s for Datagram Transport Layer Security (DTLS)~\citemark{59}.\end{citedsentence}
\afield{Claim it carries}{\enquote{IPsec adds roughly 120 $\mu$s on average to UE context setup, against roughly 430 $\mu$s for DTLS}}
\afield{Source says}{}
\begin{saidquote}(arXiv 2602.12059, the authors' copy): p. 4. \enquote{Over the F1-C, we monitor the \textbf{UE Context Setup procedure} \ldots{} To ensure statistical confidence, we \textbf{repeat the experiment 1,000 times}}, p. 5. \enquote{On the CP, specifically the F1-C interface, \textbf{all IPsec configurations add on average $\sim$120 $\mu$s to UE context setup, compared to $\sim$430 $\mu$s for DTLS}. The fastest IPsec configuration (AES256-HMAC) adds only $\sim$40 $\mu$s \ldots{} Over the E1 interface, IPsec configurations introduce on average $\sim$80 $\mu$s.}\end{saidquote}
\afield{Metadata}{[verified] Crossref matches authors, venue, pages, year and DOI.}
\afield{Verdict}{OK, both numbers and the procedure are exactly as the thesis states.}
\afield{Record}{Sources\_evidence.md, entry record.}
\auditentry{60}{natella2016fault}{peer reviewed, level~B}
\afield{Reference}{\fullcite{natella2016fault}}
\afield{Entry}{DOI \href{https://doi.org/10.1145/2841425}{\texttt{10.1145/2841425}}.}
\afield{Cited in}{2 places.}
\begin{citedsentence}\textbf{\small Section 4.2.2.}\enspace The method is therefore adversarial, following the established practice of software fault injection, where faults are introduced deliberately to see how a system behaves in their presence~\citemark{60}.\end{citedsentence}
\begin{citedsentence}\textbf{\small Appendix B.}\enspace The fourth strand is the domain and the method: the tail-latency and microsecond-scale literature~\citemark{6, 7, 31, 106}, thread-per-core structuring moving the tail in its own right~\citemark{46}, O-RAN disaggregation~\citemark{3}, control-plane latency as an active target in the core network too~\citemark{107, 108}, adoption evidence pointing at the learning curve rather than run-time cost~\citemark{80, 81, 109}, and standard guidance on repeat-based evaluation and fault injection~\citemark{33, 60}.\end{citedsentence}
\afield{Source says}{}
\begin{saidquote}Crossref-deposited abstract: \enquote{\textbf{Software Fault Injection is a method to anticipate worst-case scenarios caused by faulty software through the deliberate injection of software faults.} This survey provides a comprehensive overview of the state of the art\ldots{}}\end{saidquote}
\afield{Metadata}{[verified] Crossref matches, \texttt{pages = 44:1--44:55} is the article-number convention.}
\afield{Verdict}{OK, matches the thesis's definition at 1426 closely.}
\afield{Record}{Sources\_evidence.md, entry record.}
\auditentry{61}{efron1993bootstrap}{book, level~B}
\afield{Reference}{\fullcite{efron1993bootstrap}}
\afield{Cited in}{1 place.}
\begin{citedsentence}\textbf{\small Section 4.4.}\enspace Results are summarised by the median with a nonparametric bootstrap $95\%$ confidence interval~\citemark{61} rather than by the mean with a standard deviation, because latency distributions are heavy-tailed and skewed and a distribution-free interval on the median makes no normality assumption.\end{citedsentence}
\afield{Claim it carries}{nonparametric bootstrap 95 \% confidence intervals}
\afield{Source says}{}
\begin{saidquote}Standard reference, the percentile and BCa interval construction the thesis uses is Chapters 12--14. Book not opened for this audit.\end{saidquote}
\afield{Metadata}{[nit] \textbf{superseded by the 24 August 2026 pass.} Crossref holds the record (Efron and Tibshirani, 1993, DOI 10.1007/978-1-4899-4541-9) and the ISBN matches, which is why this entry read OK on 15 August. The DOI itself \textbf{does not resolve}: \texttt{doi.org} returns 404 and so does the Springer book page. A field comparison against Crossref cannot see that, only a fetch can. $\rightarrow$ \texttt{doi} and \texttt{url} removed, ISBN kept, withdrawn handle recorded in a BibTeX comment above the entry so that it does not print.}
\afield{Verdict}{Claim fine, link dead, fixed.}
\afield{Record}{Sources\_evidence.md, entry record.}
\auditentry{62}{rfc3393}{standard, level~A}
\afield{Reference}{\fullcite{rfc3393}}
\afield{Entry}{DOI \href{https://doi.org/10.17487/RFC3393}{\texttt{10.17487/RFC3393}}. Standards Track. Defines packet delay variation as the difference between consecutive one-way delays.}
\afield{Cited in}{2 places.}
\begin{citedsentence}\textbf{\small Section 4.4.}\enspace That is the packet-delay-variation quantity of Request for Comments (RFC)~3393~\citemark{62}, and it matches the sense in which this thesis uses the term: variation from one operation to the next, not spread of the distribution as a whole.\end{citedsentence}
\begin{citedsentence}\textbf{\small Section 6.4.1.}\enspace They do not describe how far consecutive operations differ from each other, which RFC~3393 names inter-operation delay variation (Section~\refmark{sec:stats})~\citemark{62}.\end{citedsentence}
\afield{Claim it carries}{jitter defined as the absolute difference between consecutive latencies, \enquote{the packet-delay-variation quantity of RFC 3393}}
\afield{Source says}{}
\begin{saidquote}(\texttt{rfc3393.txt}, abstract and \S{}1): \enquote{The metric is based on the difference in the One-Way-Delay of selected packets. This difference in delay is called 'IP Packet Delay Variation (ipdv)'.}\end{saidquote}
\afield{Metadata}{[verified] Crossref: Demichelis and Chimento, IETF, November 2002, DOI 10.17487/RFC3393. (Crossref does not carry the RFC number in a comparable field, the entry's \texttt{number = 3393} is correct.)}
\afield{Verdict}{OK. The thesis's usage, variation from one operation to the next, not spread of the distribution, is the RFC's sense.}
\afield{Record}{Sources\_evidence.md, entry record.}
\auditentry{63}{holm1979simple}{peer reviewed, level~B}
\afield{Reference}{\fullcite{holm1979simple}}
\afield{Entry}{\href{https://www.jstor.org/stable/4615733}{source page}.}
\afield{Cited in}{1 place.}
\begin{citedsentence}\textbf{\small Section 4.4.}\enspace Because the same comparison is made for several metrics and both allocators, the family of tests is corrected with the Holm step-down procedure~\citemark{63}, which controls the family-wise error rate without assuming independence between the tests.\end{citedsentence}
\afield{Claim it carries}{the Holm step-down procedure controlling family-wise error without assuming independence}
\afield{Source says}{}
\begin{saidquote}Metadata verified, the paper itself is behind JSTOR. The property the thesis attributes to it (FWER control without independence assumptions) is the paper's own result.\end{saidquote}
\afield{Metadata}{[verified] OpenAlex: \emph{Scandinavian Journal of Statistics}, volume 6, pages \textbf{65--70}, 1979, Sture Holm. The JSTOR stable ID in the entry's URL (4615733) matches. \textbf{Do not add a DOI here.} OpenAlex advertises \texttt{10.2307/4615733}, but that handle is unregistered: \texttt{doi.org} returns 404 and the handle API reports \enquote{not found}. It was added during this audit and then reverted. The entry should stay DOI-less, with the JSTOR URL as the locator.}
\afield{Verdict}{OK.}
\afield{Record}{Sources\_evidence.md, entry record.}
\auditentry{64}{tene\_hdrhistogram}{tool record, level~A}
\afield{Reference}{\fullcite{tene_hdrhistogram}}
\afield{Entry}{\href{https://github.com/HdrHistogram/HdrHistogram}{source page}. Section ``Corrected vs.\ raw value recording calls'' documents the \texttt{recordValueWithExpectedInterval} correction applied when a response time exceeds the expected interval between requests. Accessed August 2026.}
\afield{Cited in}{2 places.}
\begin{citedsentence}\textbf{\small Section 5.3.}\enspace Latency passes through the same HDR histogram algorithm~\citemark{64} on both sides, configured the same way (range $[1\,\text{ns}, 3.6\,\text{s}]$, three significant figures).\end{citedsentence}
\begin{citedsentence}\textbf{\small Section 6.2.}\enspace The send schedule is fixed and the histogram is corrected, so a stall introduces no coordinated-omission bias~\citemark{32, 64}.\end{citedsentence}
\afield{Source says}{}
\begin{saidquote}(\texttt{hdrhistogram.github.io}): \enquote{A Histogram that supports \textbf{recording and analyzing sampled data value counts across a configurable integer value range with configurable value precision} within the range. Value precision is expressed as the number of significant digits\ldots{}}, this fully supports the use at \texttt{Ha\_Vien\_Mastersthesis\_2026.tex} (identical HDR configuration, three significant figures).\end{saidquote}
\afield{Metadata}{[verified] Gil Tene is credited on the project page (Authors, Contributors and License section). Accessed August 2026 as stated.}
\afield{Issue}{The page does \textbf{not} discuss coordinated omission, which is what the citation is being asked to carry at 2402, 3922, \texttt{Thesis\_v3.tex} and \texttt{Thesis\_Plan\_v2.tex}. In all four places \texttt{schroeder2006openclosed} is cited alongside and does carry it (fix \#8).}
\afield{Verdict}{OK for the histogram, re-point or drop for coordinated omission.}
\afield{Record}{Sources\_evidence.md, entry record.}
\auditentry{65}{lizard}{tool record, level~A}
\afield{Reference}{\fullcite{lizard}}
\afield{Entry}{\href{https://github.com/terryyin/lizard}{source page}. Version 1.23.0. Accessed August 2026.}
\afield{Cited in}{2 places.}
\begin{citedsentence}\textbf{\small Section 5.4.4.}\enspace Table~\refmark{tab:complexity} therefore adds cyclomatic complexity, measured by \texttt{lizard}~\citemark{65} over both cores.\end{citedsentence}
\begin{citedsentence}\textbf{\small Section 5.4.4, caption.}\enspace Size and cyclomatic-complexity proxies for the two cores, measured with \texttt{lizard}~\citemark{65} applied identically to both languages.\end{citedsentence}
\afield{Claim it carries}{cyclomatic complexity measured with one tool applied identically to both languages}
\afield{Source says}{}
\begin{saidquote}Registry record as above.\end{saidquote}
\afield{Metadata}{[verified] PyPI: author \textbf{Terry Yin}, project home \texttt{lizard.ws}, GitHub \texttt{terryyin/lizard}, version \textbf{1.23.0} exists and is the current release (uploaded 2 June 2026), consistent with \enquote{Version 1.23.0. Accessed August 2026} and \texttt{year = 2026}.}
\afield{Verdict}{OK. Pinning the version in the entry is exactly right for reproducibility.}
\afield{Record}{Sources\_evidence.md, entry record.}
\auditentry{66}{zhou2025oltp}{peer reviewed, level~A}
\afield{Reference}{\fullcite{zhou2025oltp}}
\afield{Entry}{\href{https://www.vldb.org/cidrdb/papers/2025/p17-zhou.pdf}{source page}.}
\afield{Cited in}{2 places.}
\begin{citedsentence}\textbf{\small Section 6.4.1.}\enspace On a shared-nothing database partitioned one per worker thread like this service, Zhou et al.\ report that once the per-operation work is optimised what dominates the rest is communication across isolation boundaries and not the operation~\citemark{66}.\end{citedsentence}
\begin{citedsentence}\textbf{\small Section 7.6.}\enspace Three findings have published company: median equivalence below saturation, with matched pairs in three settings putting safe Rust within single-digit percentages of its C or C++ counterpart~\citemark{73, 84, 87}, a cross-thread handoff dominating the end-to-end cost~\citemark{66}, and most of RQ2's boundary classes, which the literature already documents.\end{citedsentence}
\afield{Claim it carries}{once per-operation work is optimised, cost is dominated by communication across isolation boundaries}
\afield{Source says}{}
\begin{saidquote}(p. 1, abstract): \enquote{we revisit these assumptions and demonstrate that \textbf{communication overhead is now the dominant factor affecting OLTP performance} \ldots{} revealing that \textbf{messaging costs play a critical role in end-to-end performance} \ldots{} \textbf{increased isolation can lead to higher communication costs}.}\end{saidquote}
\afield{Metadata}{[verified] DBLP confirms CIDR 2025 and assigns no pages, the entry correctly omits them. PDF retrieved from vldb.org.}
\afield{Verdict}{OK, this is the strongest external corroboration in the bibliography for the thesis's handoff-dominance result, and it is used exactly once, in the right place.}
\afield{Record}{Sources\_evidence.md, entry record.}
\auditentry{67}{clangtidyuseaftermove}{tool record, level~A}
\afield{Reference}{\fullcite{clangtidyuseaftermove}}
\afield{Entry}{\href{https://clang.llvm.org/extra/clang-tidy/checks/bugprone/use-after-move.html}{source page}. Extra Clang Tools documentation. Accessed August 2026.}
\afield{Cited in}{1 place.}
\begin{citedsentence}\textbf{\small Section 6.4.5.}\enspace A static analyser configured for it would plausibly flag the pattern, and \texttt{clang-tidy} ships such a check, whose documentation records that it follows control flow but not paths and considers only local variables and function parameters~\citemark{67}.\end{citedsentence}
\afield{Source says}{}
\begin{saidquote}\url{https://clang.llvm.org/extra/clang-tidy/checks/bugprone/use-after-move.html}, read 24 August 2026.\end{saidquote}
\begin{saidquote}\texttt{use\_after\_move} is the one in-guarantee class no C++ sanitiser reports, and \texttt{AGENT.md} requires the clang-tidy caveat wherever that is quoted. The caveat can now be cited rather than asserted, and in the tool's own words: the check \enquote{currently only considers calls of \texttt{std::move} on local variables or function parameters. It does not check moves of member variables or global variables}, it is flow-sensitive but not path-sensitive, and it flags a use-after-move even on types that define no move operations, deliberately. Cite as project documentation exactly as \texttt{tsandeadlock} is cited, and keep the sentence's claim to what the tool documents about itself.\end{saidquote}
\afield{Record}{LITERATURE\_NOTES.md, section C.9 clang-tidy `bugprone-use-after-move` documentation.}
\auditentry{68}{hughes2025gcrust}{peer reviewed, level~B}
\afield{Reference}{\fullcite{hughes2025gcrust}}
\afield{Entry}{DOI \href{https://doi.org/10.1145/3763179}{\texttt{10.1145/3763179}}.}
\afield{Cited in}{1 place.}
\begin{citedsentence}\textbf{\small Section 6.4.5.}\enspace The safe-Rust cycle leaks in silence, because leaking is safe in Rust by explicit design~\citemark{68, 69}.\end{citedsentence}
\afield{Source says}{}
\begin{saidquote}Hughes \& Tratt, PACMPL 9(OOPSLA2):3588--3614, 2025. \enquote{Garbage Collection for Rust: The Finalizer Frontier}. Opens with the fact the RQ2 boundary needs: Rust's lack of GC makes shared-ownership data structures \enquote{awkward, inefficient, or both}. Peer-reviewed support for the reference-cycle leak class, one of the two boundary classes that came back with the C++ tooling ahead.\end{saidquote}
\afield{Record}{LITERATURE\_NOTES.md, peer-reviewed additions of the third sweep.}
\auditentry{69}{cui2024rcanary}{peer reviewed, level~A}
\afield{Reference}{\fullcite{cui2024rcanary}}
\afield{Entry}{DOI \href{https://doi.org/10.1109/TSE.2024.3443624}{\texttt{10.1109/TSE.2024.3443624}}.}
\afield{Cited in}{1 place.}
\begin{citedsentence}\textbf{\small Section 6.4.5.}\enspace The safe-Rust cycle leaks in silence, because leaking is safe in Rust by explicit design~\citemark{68, 69}.\end{citedsentence}
\afield{Source says}{}
\begin{saidquote}\enquote{rCanary: Detecting Memory Leaks Across Semi-Automated Memory Management Boundary in Rust}. \emph{IEEE Transactions on Software Engineering} 50(9), pages 2472-2484, 2024. DOI 10.1109/TSE.2024.3443624. Crossref and OpenAlex records read, abstract read in full.\end{saidquote}
\begin{saidquote}The abstract states the thesis's own \texttt{leak\_cycle} conclusion in the authors' words. Rust's ownership model \enquote{is anticipated to eliminate memory leaks. However, we observed that user intervention drives it into semi-automated memory management and makes it error-prone to cause leaks. In contrast to violating memory-safety guarantees restricted by the \texttt{unsafe} keyword, \textbf{the boundary of leaking memory is implicit, and the compiler would not emit any warnings for developers}.} Applied to more than 1,200 real crates, it found 19 with leaks.\end{saidquote}
\begin{saidquote}\textbf{Why this is the single best find in the sweep.} The reference-cycle leak is the boundary class an examiner is most likely to call artificial, and this is a 2024 journal paper whose entire premise is that the class is real, silent and common enough to be worth a static checker. It also supplies the exact framing the thesis wants: the leak boundary is implicit, where the \texttt{unsafe} boundary is marked. Cite it in \S{}\texttt{sec:faultcat}, in the RQ2 discussion where LeakSanitizer catches what Rust misses, and in the future-work sentence about tooling, since rCanary is the Rust-side answer the thesis notes is missing. It pairs with \texttt{hughes2025gcrust} exactly as sweep 5 \S{}E.1 predicted.\end{saidquote}
\afield{Record}{LITERATURE\_NOTES.md, section C.3 Cui, Xu, Tian and Zhou, TSE 2024.}
\auditentry{70}{dietz2015overflow}{peer reviewed, level~A}
\afield{Reference}{\fullcite{dietz2015overflow}}
\afield{Entry}{DOI \href{https://doi.org/10.1145/2743019}{\texttt{10.1145/2743019}}.}
\afield{Cited in}{1 place.}
\begin{citedsentence}\textbf{\small Section 6.4.5.}\enspace Defined-but-wrong is easier to debug than undefined-and-wrong and no less incorrect~\citemark{70}.\end{citedsentence}
\afield{Source says}{}
\begin{saidquote}\enquote{Understanding Integer Overflow in C/C++}. \emph{ACM Transactions on Software Engineering and Methodology} 25(1), 2015. DOI 10.1145/2743019. \textbf{One field to settle before use: Crossref's deposit gives pages 1-29 while the journal numbers it as an article. Check the venue's expected form.}\end{saidquote}
\begin{saidquote}From the abstract: \enquote{intentional uses of wraparound behaviors are more common than is widely believed}, with \enquote{over 200 distinct locations in the SPEC CINT2000 benchmarks where overflow occurs}, and, orthogonally, \enquote{overflows are found in both well-defined and undefined flavors} with applications having \enquote{been broken by improvements in compiler optimizations}.\end{saidquote}
\begin{saidquote}\textbf{Why the thesis needs it.} The \texttt{integer\_overflow} row is one of the two boundary classes where the C++ tooling came back ahead, and the thesis's own framing is that Rust \emph{defines} release-mode wrapping while C++ leaves signed overflow undefined, which is why UndefinedBehaviorSanitizer has something to report. Dietz et al. is the peer-reviewed source for both halves of that: wrapping is often deliberate, so a checker cannot treat every overflow as a bug, and undefined overflow has broken real programs when optimisers improved. It also justifies why the class belongs in a defect catalogue at all rather than being dismissed as a contrived case.\end{saidquote}
\afield{Record}{LITERATURE\_NOTES.md, section C.2 Dietz, Li, Regehr and Adve, TOSEM 2015.}
\auditentry{71}{bornholt2021shardstore}{peer reviewed, level~A}
\afield{Reference}{\fullcite{bornholt2021shardstore}}
\afield{Entry}{DOI \href{https://doi.org/10.1145/3477132.3483540}{\texttt{10.1145/3477132.3483540}}.}
\afield{Cited in}{1 place.}
\begin{citedsentence}\textbf{\small Section 6.4.5.}\enspace Two Rust model checkers answer it, in the division of labour a production Rust system has been validated with: Loom explores interleavings soundly under bounded partial-order reduction, and Shuttle samples them probabilistically to reach cases beyond Loom~\citemark{71}.\end{citedsentence}
\afield{Source says}{}
\begin{saidquote}\enquote{Using Lightweight Formal Methods to Validate a Key-Value Storage Node in Amazon S3}. James Bornholt, Rajeev Joshi, Vytautas Astrauskas, Brendan Cully, Bernhard Kragl, Seth Markle, Kyle Sauri, Drew Schleit, Grant Slatton, Serdar Tasiran, Jacob Van Geffen, Andrew Warfield. SOSP '21, pages 836--850. DOI 10.1145/3477132.3483540.\end{saidquote}
\begin{saidquote}\textbf{This is the find of the sweep.} It is a peer-reviewed account of a production Rust system validated with \emph{the same two model checkers this thesis uses, for the same reason}, by the team that wrote one of them. Read the paper's own description against \texttt{Ha\_Vien\_Mastersthesis\_2026.tex}:2932--3047 and the draft is doing the same thing on a smaller object.\end{saidquote}
\begin{saidquote}Three passages are directly usable, quoted from the paper's section on checking concurrent executions:\end{saidquote}
\begin{saidquote}\begin{itemize}[leftmargin=1.2em,itemsep=1pt,topsep=1pt]
\item Loom \enquote{implements the CDSChecker algorithm for sound model checking in the release/acquire memory model, and uses bounded partial-order reduction to reduce state explosion.}
\item \enquote{we developed and open-sourced the stateless model checker Shuttle, which implements randomized algorithms such as probabilistic concurrency testing.}
\item \enquote{We use Loom to soundly check all interleavings of small, correctness-critical code such as custom concurrency primitives, and Shuttle to randomly check interleavings of larger test harnesses to which Loom does not scale.}
\end{itemize}\end{saidquote}
\begin{saidquote}That third sentence is the division of labour \texttt{faults/run\_model\_checking.sh} implements, stated independently in a top venue.\end{saidquote}
\afield{Record}{LITERATURE\_NOTES.md, section B.1 Bornholt et al., SOSP 2021.}
\auditentry{72}{mccormack2025ffi}{peer reviewed, level~A}
\afield{Reference}{\fullcite{mccormack2025ffi}}
\afield{Entry}{DOI \href{https://doi.org/10.1109/ICSE55347.2025.00167}{\texttt{10.1109/ICSE55347.2025.00167}}.}
\afield{Cited in}{1 place.}
\begin{citedsentence}\textbf{\small Section 6.4.5.}\enspace Driving Miri together with an interpreter for the foreign side is a published research result rather than something a migration can adopt, and it found 46 instances of undefined or undesired behaviour in 37 libraries that call foreign functions, one of them in a library the Rust project maintains itself~\citemark{72}.\end{citedsentence}
\afield{Claim it carries}{no shipped tool checks the C++ side of the seam}
\afield{Source says}{}
\begin{saidquote}(abstract): \enquote{Miri is currently the only dynamic analysis tool that can validate applications against these models, but \textbf{it does not support finding bugs in foreign functions} \ldots{} We used Miri and an LLVM interpreter to jointly execute applications that call foreign functions, where we \textbf{found 46 instances of undefined or undesired behavior in 37 libraries} \ldots{} and \textbf{one was found in a library maintained by the Rust Project}.} \textbf{What it may not be made to carry.} The three high-download bugs are \enquote{libraries that had more than 10,000 daily downloads on average during our observation period} and \textbf{not} libraries that had been audited. Do not upgrade the download figure into a review claim. \textbf{One thing it changed in the thesis beyond the citation.} \S{}7.4 read that Miri \enquote{interprets the program against the Stacked Borrows aliasing model} and reports the second release. The artefact does not support the attribution: \texttt{faults/run\_miri.sh} sets no \texttt{MIRIFLAGS} at all, so the run took Miri's defaults, and \texttt{faults/results/miri-seam\_double\_free.log} reports a \texttt{Box} validity error (\enquote{encountered a dangling box (use-after-free)}) rather than an aliasing violation. The model naming is removed. This paper is what made it matter, since it reports \enquote{the latest Tree Borrows model was significantly more permissive than the earlier Stacked Borrows model}, so naming a model implies a configuration this study never set.\end{saidquote}
\afield{Metadata}{[verified] Crossref matches the three-author list, the pages and the venue. Same first author and group as \texttt{mccormack2026unsafe}, which is the mixed-methods companion to this measurement.}
\afield{Verdict}{OK.}
\afield{Record}{Sources\_evidence.md, entry record.}
\auditentry{73}{emmerich2019drivers}{peer reviewed, level~A}
\afield{Reference}{\fullcite{emmerich2019drivers}}
\afield{Entry}{DOI \href{https://doi.org/10.1109/ANCS.2019.8901892}{\texttt{10.1109/ANCS.2019.8901892}}.}
\afield{Cited in}{5 places.}
\begin{citedsentence}\textbf{\small Section 7.2.}\enspace The median result is consistent with the wider evidence that Rust's safety checks are largely absorbed by modern out-of-order CPUs~\citemark{12, 73}.\end{citedsentence}
\begin{citedsentence}\textbf{\small Section 7.5.}\enspace Its safety checks are absorbed by out-of-order execution, not literally free~\citemark{73}.\end{citedsentence}
\begin{citedsentence}\textbf{\small Section 7.5.}\enspace None is unoccupied on its own, as the matched implementations~\citemark{73, 84, 85, 86, 87, 91, 57}, the percentile sweep stopping at a latency requirement~\citemark{56} and the effort measurements~\citemark{53, 80} all show.\end{citedsentence}
\begin{citedsentence}\textbf{\small Section 7.6.}\enspace Three findings have published company: median equivalence below saturation, with matched pairs in three settings putting safe Rust within single-digit percentages of its C or C++ counterpart~\citemark{73, 84, 87}, a cross-thread handoff dominating the end-to-end cost~\citemark{66}, and most of RQ2's boundary classes, which the literature already documents.\end{citedsentence}
\begin{citedsentence}\textbf{\small Appendix B.}\enspace Compiled systems languages sit close on time and energy~\citemark{45, 93}, Rust's safety is cheap at run time~\citemark{12, 94}, and that holds where timing is tightest, in embedded hypervisors and in packet processing within a few percent of a C reference~\citemark{95, 73, 96}.\end{citedsentence}
\afield{Claim it carries}{most sharply at : \enquote{running 4 \% slower than its C reference while executing 63 \% more instructions per packet}}
\afield{Source says}{}
\begin{saidquote}(arXiv 1909.06344, p. 1, abstract): \enquote{Out-of-order CPUs mitigate the cost of safety checks: \textbf{Our Rust driver executes 63\% more instructions per packet but is only 4\% slower than a reference C implementation.}} Table 9 (p. 11) gives the normalised results (Rust 98 \% of C).\end{saidquote}
\afield{Metadata}{[verified] Author list follows the paper (IEEE's record drops \enquote{García} from Esaú García Sánchez-Torija, the entry is right).}
\afield{Verdict}{OK, both numbers are the source's own, in the source's own order.}
\afield{Record}{Sources\_evidence.md, entry record.}
\auditentry{74}{zhang2022runtime}{peer reviewed, level~A}
\afield{Reference}{\fullcite{zhang2022runtime}}
\afield{Entry}{DOI \href{https://doi.org/10.1145/3551349.3559494}{\texttt{10.1145/3551349.3559494}}.}
\afield{Cited in}{4 places.}
\begin{citedsentence}\textbf{\small Section 7.2.}\enspace The prediction also named where to look first if that happened: bounds and overflow checking~\citemark{74, 75}.\end{citedsentence}
\begin{citedsentence}\textbf{\small Section 7.2.}\enspace The mechanisms the literature has localised~\citemark{74, 75} are candidates to test, not conclusions to import.\end{citedsentence}
\begin{citedsentence}\textbf{\small Section 7.6.}\enspace The mechanisms most likely responsible have been isolated elsewhere~\citemark{74, 75}.\end{citedsentence}
\begin{citedsentence}\textbf{\small Appendix B.}\enspace Where a cost does appear it has been pinned down: to bounds, overflow and division checks, to that same check set against hardware-enforced capabilities, and to recurring anti-patterns~\citemark{74, 75, 97}.\end{citedsentence}
\afield{Claim it carries}{the literature localises Rust's cost to bounds, overflow and division checks, which is where the thesis's registered prediction said to look first}
\afield{Source says}{}
\begin{saidquote}(author copy, \S{}1 and \S{}4 with Figure 2): \enquote{the major source of performance overhead is the run-time checks inserted by the compiler to ensure memory safety. For instance, \textbf{out-of-bounds checks account for 52.7\% of all the overhead}.} And: \enquote{With all run-time checks enabled, Rust incurs an 1.96x overhead, which drops to 1.77x, 1.27x, and 1.23x after we in turn disable \textbf{integer-overflow checks, out-of-bound checks, and division-by-zero checks}.}\end{saidquote}
\afield{Metadata}{[verified] Crossref. Author copy at \texttt{softsec.cs.utah.edu/papers/RustPerf.pdf}, and ACM publishes a full-HTML version.}
\afield{Verdict}{OK, and precise. The source names exactly the three check kinds the thesis attributes to it, in that order, which is the strongest available support for \S{}7.2's account of what the registered prediction got wrong.}
\afield{Record}{Sources\_evidence.md, entry record.}
\auditentry{75}{harris2023morello}{peer reviewed, level~A}
\afield{Reference}{\fullcite{harris2023morello}}
\afield{Entry}{DOI \href{https://doi.org/10.4230/LIPIcs.ECOOP.2023.39}{\texttt{10.4230/LIPIcs.ECOOP.2023.39}}.}
\afield{Cited in}{4 places.}
\begin{citedsentence}\textbf{\small Section 7.2.}\enspace The prediction also named where to look first if that happened: bounds and overflow checking~\citemark{74, 75}.\end{citedsentence}
\begin{citedsentence}\textbf{\small Section 7.2.}\enspace The mechanisms the literature has localised~\citemark{74, 75} are candidates to test, not conclusions to import.\end{citedsentence}
\begin{citedsentence}\textbf{\small Section 7.6.}\enspace The mechanisms most likely responsible have been isolated elsewhere~\citemark{74, 75}.\end{citedsentence}
\begin{citedsentence}\textbf{\small Appendix B.}\enspace Where a cost does appear it has been pinned down: to bounds, overflow and division checks, to that same check set against hardware-enforced capabilities, and to recurring anti-patterns~\citemark{74, 75, 97}.\end{citedsentence}
\afield{Claim it carries}{the bounds check specifically, set against hardware-enforced capabilities}
\afield{Source says}{}
\begin{saidquote}(\texttt{LIPIcs.ECOOP.2023.39.pdf}, p. 3, contributions list): \enquote{A performance analysis of Rust on Morello (\S{}4, \S{}5): \ldots{} \textbf{Analysis of the benefits of bounds checking elision on capability hardware (\S{}5.1)}.} Abstract (p. 1): \enquote{a comparison of the performance of Rust's runtime safety checks and the hardware-supported checks of Morello}.\end{saidquote}
\afield{Metadata}{[verified] The LIPIcs PDF's own running heads confirm \textbf{39:1--39:27} and volume 263. DataCite gives the container as \enquote{LIPIcs, Volume 263, ECOOP 2023}, the entry's booktitle is the conference name, which is conventional.}
\afield{Verdict}{OK.}
\afield{Record}{Sources\_evidence.md, entry record.}
\auditentry{76}{zhang2026deadlock}{peer reviewed, level~A}
\afield{Reference}{\fullcite{zhang2026deadlock}}
\afield{Entry}{DOI \href{https://doi.org/10.1007/s10515-026-00614-z}{\texttt{10.1007/s10515-026-00614-z}}. Extends the preprint arXiv:2401.01114.}
\afield{Cited in}{1 place.}
\begin{citedsentence}\textbf{\small Section 7.2.}\enspace Static deadlock detection for Rust is an active research problem in its own right~\citemark{76}, and model checkers for Rust kernel code are motivated explicitly by the observation that ownership and borrowing leave data races, deadlocks, and livelocks reachable~\citemark{77}.\end{citedsentence}
\afield{Claim it carries}{static deadlock detection for Rust is an active research problem}
\afield{Source says}{}
\begin{saidquote}(the preprint the entry's \texttt{note} points to, arXiv 2401.01114, p. 1): \enquote{Although Rust claims to enable fearless concurrency, \textbf{concurrency errors remain prevalent in practical applications, with deadlocks particularly prominent and harmful}. Current static deadlock detection techniques either do not cater to Rust programs \ldots{} or solely concentrate on resource deadlocks, disregarding communication deadlocks.}\end{saidquote}
\afield{Metadata}{[verified] Crossref resolves the DOI and matches authors, journal, volume, issue and year.}
\afield{Verdict}{OK. Good that the journal version is cited rather than the preprint.}
\afield{Record}{Sources\_evidence.md, entry record.}
\auditentry{77}{tang2025converos}{peer reviewed, level~A}
\afield{Reference}{\fullcite{tang2025converos}}
\afield{Entry}{\href{https://www.usenix.org/conference/atc25/presentation/tang}{source page}.}
\afield{Cited in}{1 place.}
\begin{citedsentence}\textbf{\small Section 7.2.}\enspace Static deadlock detection for Rust is an active research problem in its own right~\citemark{76}, and model checkers for Rust kernel code are motivated explicitly by the observation that ownership and borrowing leave data races, deadlocks, and livelocks reachable~\citemark{77}.\end{citedsentence}
\afield{Claim it carries}{model checkers for Rust kernel code are motivated by ownership and borrowing leaving data races, deadlocks and livelocks reachable}
\afield{Source says}{}
\begin{saidquote}(\texttt{atc25-tang.pdf}, p. 2): \enquote{\textbf{Although Rust greatly enhances memory safety and prevents data races in safe code, OS kernel development inevitably requires unsafe, which may still introduce concurrency-related undefined behaviors such as data races. Logical errors and flawed designs can lead to deadlocks, livelocks, and functional violations.}} And, in the abstract: \enquote{uncovering 20 bugs that led to issues such as data races, deadlocks, livelocks, and kernel panics.}\end{saidquote}
\afield{Metadata}{[verified] DBLP confirms pages \textbf{1143--1159} and the six-author list.}
\afield{Verdict}{OK, an unusually exact match to the sentence it supports.}
\afield{Record}{Sources\_evidence.md, entry record.}
\auditentry{78}{lu2008mistakes}{peer reviewed, level~A}
\afield{Reference}{\fullcite{lu2008mistakes}}
\afield{Entry}{DOI \href{https://doi.org/10.1145/1346281.1346323}{\texttt{10.1145/1346281.1346323}}.}
\afield{Cited in}{1 place.}
\begin{citedsentence}\textbf{\small Section 7.2.}\enspace In a field study of $105$ real concurrency bugs, about a third of the non-deadlock cases violated the order the programmers intended, not their locking, and the authors record that such intentions are not readily expressible through locks at all~\citemark{78}.\end{citedsentence}
\afield{Source says}{}
\begin{saidquote}\enquote{Learning from Mistakes: A Comprehensive Study on Real World Concurrency Bug Characteristics}. ASPLOS 2008, pages 329-339. DOI 10.1145/1346281.1346323. Crossref record read, abstract read in full.\end{saidquote}
\begin{saidquote}105 concurrency bugs from MySQL, Apache, Mozilla and OpenOffice. The finding that matters here, quoted from the abstract: \enquote{around one third of the examined non-deadlock concurrency bugs are caused by violation to programmers' order intentions, which may not be easily expressed via synchronization primitives like locks and transactional memories}. A second one is directly about detection: \enquote{around 34\% of the examined non-deadlock concurrency bugs involve multiple variables, which are not well addressed by existing bug detection tools}.\end{saidquote}
\begin{saidquote}\textbf{Why the thesis needs it.} The \texttt{lost\_update} boundary class is currently defended by the thesis's own reasoning: every access is correctly synchronised, so there is no unsynchronised access for ThreadSanitizer to find and nothing for the borrow checker to reject. This is a canonical, heavily cited paper saying that a third of real concurrency bugs are of exactly that shape, and that they are not expressible as a locking discipline. It converts the thesis's most easily doubted RQ2 row into an instance of a documented category. Cite it in \S{}\texttt{sec:faultcat} beside \texttt{natella2016fault} and in the RQ2 discussion where the lost update is explained.\end{saidquote}
\begin{saidquote}\textbf{Caveat to state in the same sentence.} The sample is 2008-era C and C++ desktop and server software, so it establishes the category and not its present-day frequency.\end{saidquote}
\afield{Record}{LITERATURE\_NOTES.md, section C.1 Lu, Park, Seo and Zhou, ASPLOS 2008.}
\auditentry{79}{sundvisson2026drivers}{thesis, level~A}
\afield{Reference}{\fullcite{sundvisson2026drivers}}
\afield{Entry}{\href{https://urn.kb.se/resolve?urn=urn:nbn:se:bth-29775}{source page}.}
\afield{Cited in}{3 places.}
\begin{citedsentence}\textbf{\small Section 7.2.}\enspace A second Ericsson-hosted study of RAN node management software found Rust's peak memory usage generally higher than C's~\citemark{54}, and a matched pair of user-space device drivers outside telecom put the Rust version's footprint about 29\% above the C version's~\citemark{79}.\end{citedsentence}
\begin{citedsentence}\textbf{\small Section 7.5.}\enspace Two 2026 bachelor's theses from outside the vendor set apply parts of the method to smaller objects: matched C and Rust user-space drivers, with a fixed memory overhead and no measurable run-time cost~\citemark{79}, and a CWE staging study across C++, Rust and Python~\citemark{89}.\end{citedsentence}
\begin{citedsentence}\textbf{\small Section 7.6.}\enspace Peak resident set size sits between the two groups, with the sign found here reported by two neighbouring theses and reversed by a third~\citemark{54, 79, 53}.\end{citedsentence}
\afield{Claim it carries}{matched C and Rust user-space drivers with the Rust footprint about 29 per cent above C, a fixed overhead and no measurable run-time cost}
\afield{Source says}{}
\begin{saidquote}Located in two places.\end{saidquote}
\begin{saidquote}\begin{itemize}[leftmargin=1.2em,itemsep=1pt,topsep=1pt]
\item p. 3, abstract: \enquote{The C driver was approximately 0.178 \% faster, a difference detectable at 99 \% confidence but well within a $\pm$1 \% equivalence zone. The Rust driver consumed approximately 29 \% more memory and contained five small unsafe blocks wrapping unavoidable system calls.}
\item p. 35: \enquote{Rust maintains a constant footprint of 1800 kB, whereas C exhibits 1400 kB, corresponding to an approximately 29 \% increase.}
\end{itemize}\end{saidquote}
\afield{Metadata}{[verified] DiVA \texttt{diva2:2071063}, full text \texttt{FULLTEXT01.pdf} retrieved.}
\afield{Verdict}{OK, and stronger than the thesis needs. \enquote{Constant footprint} is the source's own word for the fixed overhead, and the $\pm$1 per cent equivalence zone is what licenses \enquote{no measurable run-time cost}.}
\afield{Record}{Sources\_evidence.md, entry record.}
\auditentry{80}{fulton2021adopting}{peer reviewed, level~A}
\afield{Reference}{\fullcite{fulton2021adopting}}
\afield{Entry}{\href{https://www.usenix.org/conference/soups2021/presentation/fulton}{source page}.}
\afield{Cited in}{4 places.}
\begin{citedsentence}\textbf{\small Section 7.2.}\enspace The dominant cost is therefore human and not machine: designing a C-compatible surface and, as the adoption studies report~\citemark{80, 81}, climbing Rust's learning curve.\end{citedsentence}
\begin{citedsentence}\textbf{\small Section 7.5.}\enspace None is unoccupied on its own, as the matched implementations~\citemark{73, 84, 85, 86, 87, 91, 57}, the percentile sweep stopping at a latency requirement~\citemark{56} and the effort measurements~\citemark{53, 80} all show.\end{citedsentence}
\begin{citedsentence}\textbf{\small Appendix B.}\enspace The fourth strand is the domain and the method: the tail-latency and microsecond-scale literature~\citemark{6, 7, 31, 106}, thread-per-core structuring moving the tail in its own right~\citemark{46}, O-RAN disaggregation~\citemark{3}, control-plane latency as an active target in the core network too~\citemark{107, 108}, adoption evidence pointing at the learning curve rather than run-time cost~\citemark{80, 81, 109}, and standard guidance on repeat-based evaluation and fault injection~\citemark{33, 60}.\end{citedsentence}
\begin{citedsentence}\textbf{\small Appendix G.}\enspace The dominant cost identified here is human, which matches what adoption studies report~\citemark{80}.\end{citedsentence}
\afield{Claim it carries}{the dominant cost is human, not run-time}
\afield{Source says}{}
\begin{saidquote}(\texttt{soups2021-fulton.pdf}, p. 2, abstract): \enquote{a range of positive features \ldots{} \textbf{as well as drawbacks including a steep learning curve, limited library support, and concerns about the ability to hire additional Rust developers}}, p. 3: \enquote{\textbf{most seriously a steep learning curve} to adjust to the paradigms that enforce security guarantees} and \enquote{much of the \textbf{cost of adoption occurs up front}, while benefits tend to accrue later}.\end{saidquote}
\afield{Metadata}{[verified] DBLP confirms pages \textbf{597--616} and the five-author list.}
\afield{Verdict}{OK.}
\afield{Record}{Sources\_evidence.md, entry record.}
\auditentry{81}{crichton2024learning}{peer reviewed, level~A}
\afield{Reference}{\fullcite{crichton2024learning}}
\afield{Entry}{DOI \href{https://doi.org/10.1145/3649812}{\texttt{10.1145/3649812}}.}
\afield{Cited in}{2 places.}
\begin{citedsentence}\textbf{\small Section 7.2.}\enspace The dominant cost is therefore human and not machine: designing a C-compatible surface and, as the adoption studies report~\citemark{80, 81}, climbing Rust's learning curve.\end{citedsentence}
\begin{citedsentence}\textbf{\small Appendix B.}\enspace The fourth strand is the domain and the method: the tail-latency and microsecond-scale literature~\citemark{6, 7, 31, 106}, thread-per-core structuring moving the tail in its own right~\citemark{46}, O-RAN disaggregation~\citemark{3}, control-plane latency as an active target in the core network too~\citemark{107, 108}, adoption evidence pointing at the learning curve rather than run-time cost~\citemark{80, 81, 109}, and standard guidance on repeat-based evaluation and fault injection~\citemark{33, 60}.\end{citedsentence}
\afield{Claim it carries}{the learning curve as the dominant adoption cost}
\afield{Source says}{}
\begin{saidquote}(arXiv 2401.01257, p. 1, abstract, restated p. 2): \enquote{Over 13 months, \textbf{62,526 readers answered questions 1,140,202 times} \ldots{} we find that \textbf{many readers drop-out of the book early when faced with difficult language concepts like Rust's ownership types}.}\end{saidquote}
\afield{Metadata}{[verified] Crossref matches every field.}
\afield{Verdict}{OK.}
\afield{Record}{Sources\_evidence.md, entry record.}
\auditentry{82}{li2026rustlinux}{peer reviewed, level~B}
\afield{Reference}{\fullcite{li2026rustlinux}}
\afield{Entry}{DOI \href{https://doi.org/10.1145/3789260}{\texttt{10.1145/3789260}}. Extended version of the authors' USENIX ATC '24 study of Rust-for-Linux.}
\afield{Cited in}{3 places.}
\begin{citedsentence}\textbf{\small Section 7.2.}\enspace That makes migration risk an organisational concern and not a performance one, and the two largest documented adoptions into existing C and C++ codebases report the same ordering of costs~\citemark{82, 83}.\end{citedsentence}
\begin{citedsentence}\textbf{\small Section 7.5.}\enspace It finds that Rust's safety features introduce runtime overhead and development complexity when they are not carefully managed~\citemark{82}.\end{citedsentence}
\begin{citedsentence}\textbf{\small Appendix B.}\enspace The seam carries its own defect classes and its own tooling bugs, lighter crossing models exist for embedded use, aliasing patterns expressible in C but not in safe Rust bound translation in principle, interoperation is the central open problem in embedded Rust, \texttt{unsafe} concentrates at exactly these boundaries, and the two largest documented adoptions into existing C and C++ report the resulting friction empirically~\citemark{50, 51, 101, 102, 103, 104, 82, 83}.\end{citedsentence}
\afield{Source says}{}
\begin{saidquote}Li, Guo, Yang, Wang \& Xu, ACM TOCS 44(2), 2026. \enquote{Rust Meets Linux: Lessons from an Evolving Experiment}. Same author list as \texttt{li2024rustforlinux} (USENIX ATC '24): this is the journal version, and it sharpens the conclusion. \enquote{If not properly managed, Rust's safety features even introduce notable runtime overhead and increased development complexity.} Citation \textbf{upgrade}, and a contrary result the thesis should carry rather than soften. \emph{Article number not confirmed. Crossref reports pages 1--33, check before citing with a page range.}\end{saidquote}
\afield{Record}{LITERATURE\_NOTES.md, peer-reviewed additions of the third sweep.}
\auditentry{83}{fan2026android}{peer reviewed, level~B}
\afield{Reference}{\fullcite{fan2026android}}
\afield{Entry}{DOI \href{https://doi.org/10.1145/3803437.3805227}{\texttt{10.1145/3803437.3805227}}.}
\afield{Cited in}{2 places.}
\begin{citedsentence}\textbf{\small Section 7.2.}\enspace That makes migration risk an organisational concern and not a performance one, and the two largest documented adoptions into existing C and C++ codebases report the same ordering of costs~\citemark{82, 83}.\end{citedsentence}
\begin{citedsentence}\textbf{\small Appendix B.}\enspace The seam carries its own defect classes and its own tooling bugs, lighter crossing models exist for embedded use, aliasing patterns expressible in C but not in safe Rust bound translation in principle, interoperation is the central open problem in embedded Rust, \texttt{unsafe} concentrates at exactly these boundaries, and the two largest documented adoptions into existing C and C++ report the resulting friction empirically~\citemark{50, 51, 101, 102, 103, 104, 82, 83}.\end{citedsentence}
\afield{Claim it carries}{one of the two largest documented adoptions into an existing C and C++ codebase}
\afield{Source says}{}
\begin{saidquote}(deposited abstract): \enquote{the first longitudinal and systematic empirical study of \textbf{Rust in AOSP}. By analyzing 16 quarterly snapshots from January 2021 to October 2024, we track adoption trends and map the distribution of Rust code across 67 components ... we characterize platform-level integration mechanisms and examine unsafe governance.}\end{saidquote}
\afield{Verdict}{OK for the adoption and integration claim. The \enquote{same ordering of costs} that \S{}7.2 draws from pairing this with \texttt{li2026rustlinux} is the thesis's own synthesis of the two, which is how it is written.}
\afield{Record}{Sources\_evidence.md, entry record.}
\auditentry{84}{li2021ipc}{peer reviewed, level~A}
\afield{Reference}{\fullcite{li2021ipc}}
\afield{Entry}{DOI \href{https://doi.org/10.1145/3477113.3487275}{\texttt{10.1145/3477113.3487275}}.}
\afield{Cited in}{3 places.}
\begin{citedsentence}\textbf{\small Section 7.5.}\enspace Four systems studies come closest in method, though not in domain, because they too build the same component twice and attribute the difference: inter-process communication (IPC) mechanisms for one network function~\citemark{84}, RedLeaf against purpose-built C baselines~\citemark{85}, a concurrent translation of the Stanford Transactional Applications for Multi-Processing (STAMP) benchmarks that is the closest published match to this design~\citemark{86}, and a kernel User Datagram Protocol (UDP) tunnelling driver~\citemark{87}.\end{citedsentence}
\begin{citedsentence}\textbf{\small Section 7.5.}\enspace None is unoccupied on its own, as the matched implementations~\citemark{73, 84, 85, 86, 87, 91, 57}, the percentile sweep stopping at a latency requirement~\citemark{56} and the effort measurements~\citemark{53, 80} all show.\end{citedsentence}
\begin{citedsentence}\textbf{\small Section 7.6.}\enspace Three findings have published company: median equivalence below saturation, with matched pairs in three settings putting safe Rust within single-digit percentages of its C or C++ counterpart~\citemark{73, 84, 87}, a cross-thread handoff dominating the end-to-end cost~\citemark{66}, and most of RQ2's boundary classes, which the literature already documents.\end{citedsentence}
\afield{Claim it carries}{a matched C++ and Rust pair in a network-function setting putting safe Rust within single-digit percentages of unsafe C++}
\afield{Source says}{}
\begin{saidquote}(author copy, abstract and \S{}4 Conclusions): \enquote{we develop two implementations (one in C++ and one in Rust) of a network function processing framework similar to Netbricks ... We find that while high-level Rust abstractions can introduce significant overhead, a carefully planned (with respect to performance) Rust implementation remains fast and \textbf{executes within 4-8\% of unsafe C++}.} Restated in the conclusions: \enquote{on realistic workloads Rust incurs only minor overhead of 4-8\% compared to unsafe C++.}\end{saidquote}
\afield{Metadata}{[verified] Crossref, pages 53-61. Author copies at \texttt{mars-research.github.io} and \texttt{users.cs.utah.edu/\textasciitilde{}aburtsev}.}
\afield{Verdict}{OK. 4 to 8 per cent is single-digit percentages, and the matched-pair design is the source's own. Two scope notes above: the figure is conditional on a carefully planned implementation, and the paper builds several network functions rather than one.}
\afield{Record}{Sources\_evidence.md, entry record.}
\auditentry{85}{narayanan2020redleaf}{peer reviewed, level~A}
\afield{Reference}{\fullcite{narayanan2020redleaf}}
\afield{Entry}{\href{https://www.usenix.org/conference/osdi20/presentation/narayanan-vikram}{source page}.}
\afield{Cited in}{2 places.}
\begin{citedsentence}\textbf{\small Section 7.5.}\enspace Four systems studies come closest in method, though not in domain, because they too build the same component twice and attribute the difference: inter-process communication (IPC) mechanisms for one network function~\citemark{84}, RedLeaf against purpose-built C baselines~\citemark{85}, a concurrent translation of the Stanford Transactional Applications for Multi-Processing (STAMP) benchmarks that is the closest published match to this design~\citemark{86}, and a kernel User Datagram Protocol (UDP) tunnelling driver~\citemark{87}.\end{citedsentence}
\begin{citedsentence}\textbf{\small Section 7.5.}\enspace None is unoccupied on its own, as the matched implementations~\citemark{73, 84, 85, 86, 87, 91, 57}, the percentile sweep stopping at a latency requirement~\citemark{56} and the effort measurements~\citemark{53, 80} all show.\end{citedsentence}
\afield{Claim it carries}{evaluated against C baselines constructed for the purpose}
\afield{Source says}{}
\begin{saidquote}(\texttt{osdi20-narayanan\_vikram.pdf}): p. 2. \enquote{we develop efficient versions of a 10Gbps Intel ixgbe network and NVMe solid-state disk device drivers that \textbf{match the performance of the fastest DPDK and SPDK equivalents}}, p. 12, Figure 5. \enquote{C vs Rust performance comparison}, and the ixgbe driver comparison against Linux and DPDK.\end{saidquote}
\afield{Metadata}{[verified] DBLP confirms pages \textbf{21--39} and the seven-author list as entered.}
\afield{Verdict}{OK.}
\afield{Record}{Sources\_evidence.md, entry record.}
\auditentry{86}{suchert2023stamp}{peer reviewed, level~B}
\afield{Reference}{\fullcite{suchert2023stamp}}
\afield{Entry}{DOI \href{https://doi.org/10.1007/978-3-031-31180-2_10}{\texttt{10.1007/978-3-031-31180-2\_10}}.}
\afield{Cited in}{3 places.}
\begin{citedsentence}\textbf{\small Section 7.5.}\enspace Four systems studies come closest in method, though not in domain, because they too build the same component twice and attribute the difference: inter-process communication (IPC) mechanisms for one network function~\citemark{84}, RedLeaf against purpose-built C baselines~\citemark{85}, a concurrent translation of the Stanford Transactional Applications for Multi-Processing (STAMP) benchmarks that is the closest published match to this design~\citemark{86}, and a kernel User Datagram Protocol (UDP) tunnelling driver~\citemark{87}.\end{citedsentence}
\begin{citedsentence}\textbf{\small Section 7.5.}\enspace Idiomatic safe Rust implementations of the STAMP transactional benchmarks run slower than both the C originals and machine-transpiled Rust~\citemark{86}, and a general study of Rust against C reports it significantly slower in some cases~\citemark{90}.\end{citedsentence}
\begin{citedsentence}\textbf{\small Section 7.5.}\enspace None is unoccupied on its own, as the matched implementations~\citemark{73, 84, 85, 86, 87, 91, 57}, the percentile sweep stopping at a latency requirement~\citemark{56} and the effort measurements~\citemark{53, 80} all show.\end{citedsentence}
\afield{Source says}{}
\begin{saidquote}Suchert \& Castrillon, Bench 2022 (LNCS 13852:160--175, publ. 2023). Manually translated the \textbf{STAMP} transactional benchmark suite to Rust, plus a \texttt{c2rust}-transpiled version. Finding: \enquote{idiomatic implementations using safe Rust are generally slower than both C and transpiled code, they guarantee memory safety and improve code quality.} A matched pair on a \textbf{concurrent} workload, contrary result, peer-reviewed. Better neighbour for \texttt{tab:related} than any of the general-purpose benchmark rows.\end{saidquote}
\afield{Record}{LITERATURE\_NOTES.md, peer-reviewed additions of the third sweep.}
\auditentry{87}{gonzalez2023rustudp}{peer reviewed, level~B}
\afield{Reference}{\fullcite{gonzalez2023rustudp}}
\afield{Entry}{DOI \href{https://doi.org/10.1145/3623759.3624547}{\texttt{10.1145/3623759.3624547}}.}
\afield{Cited in}{3 places.}
\begin{citedsentence}\textbf{\small Section 7.5.}\enspace Four systems studies come closest in method, though not in domain, because they too build the same component twice and attribute the difference: inter-process communication (IPC) mechanisms for one network function~\citemark{84}, RedLeaf against purpose-built C baselines~\citemark{85}, a concurrent translation of the Stanford Transactional Applications for Multi-Processing (STAMP) benchmarks that is the closest published match to this design~\citemark{86}, and a kernel User Datagram Protocol (UDP) tunnelling driver~\citemark{87}.\end{citedsentence}
\begin{citedsentence}\textbf{\small Section 7.5.}\enspace None is unoccupied on its own, as the matched implementations~\citemark{73, 84, 85, 86, 87, 91, 57}, the percentile sweep stopping at a latency requirement~\citemark{56} and the effort measurements~\citemark{53, 80} all show.\end{citedsentence}
\begin{citedsentence}\textbf{\small Section 7.6.}\enspace Three findings have published company: median equivalence below saturation, with matched pairs in three settings putting safe Rust within single-digit percentages of its C or C++ counterpart~\citemark{73, 84, 87}, a cross-thread handoff dominating the end-to-end cost~\citemark{66}, and most of RQ2's boundary classes, which the literature already documents.\end{citedsentence}
\afield{Source says}{}
\begin{saidquote}Gonzalez, Mvondo \& Bromberg, PLOS@SOSP 2023 (pp. 18--25). Rust UDP tunnelling driver in the Linux kernel against an equivalent C driver: latency \textbf{+0.1906\%} (p = 1.464e-15), throughput \textbf{$-$0.00090\%} (p = 6.004e-5). A matched networking pair whose difference is statistically overwhelming and practically nil, the cleanest published instance of the distinction the thesis's pre-committed interval-vs-\emph{p} rule (\S{}5.5) exists to enforce.\end{saidquote}
\afield{Record}{LITERATURE\_NOTES.md, peer-reviewed additions of the third sweep.}
\auditentry{88}{zhuang2026flowprioritize}{peer reviewed, level~B}
\afield{Reference}{\fullcite{zhuang2026flowprioritize}}
\afield{Entry}{DOI \href{https://doi.org/10.1016/j.softx.2026.102771}{\texttt{10.1016/j.softx.2026.102771}}.}
\afield{Cited in}{1 place.}
\begin{citedsentence}\textbf{\small Section 7.5.}\enspace Priority scheduling is something an asynchronous Rust runtime has to be given by a library~\citemark{88}.\end{citedsentence}
\afield{Source says}{}
\begin{saidquote}Zhuang, SoftwareX 35:102771, 2026. FlowPrioritize, a priority-scheduling library for async Rust runtimes. Supporting evidence, alongside \texttt{skoudlil2025asyncrust}, that priority and determinism are things async Rust runtimes must be \emph{given}, not properties they have.\end{saidquote}
\afield{Record}{LITERATURE\_NOTES.md, peer-reviewed additions of the third sweep.}
\auditentry{89}{andersson2026stdmemsafety}{thesis, level~A}
\afield{Reference}{\fullcite{andersson2026stdmemsafety}}
\afield{Entry}{\href{https://urn.kb.se/resolve?urn=urn:nbn:se:liu:diva-226479}{source page}. Report LIU-IDA/LITH-EX-G--26/057--SE.}
\afield{Cited in}{1 place.}
\begin{citedsentence}\textbf{\small Section 7.5.}\enspace Two 2026 bachelor's theses from outside the vendor set apply parts of the method to smaller objects: matched C and Rust user-space drivers, with a fixed memory overhead and no measurable run-time cost~\citemark{79}, and a CWE staging study across C++, Rust and Python~\citemark{89}.\end{citedsentence}
\afield{Claim it carries}{a CWE staging study across C++, Rust and Python}
\afield{Source says}{}
\begin{saidquote}(\texttt{FULLTEXT01.pdf}, p. 4): \enquote{the manual memory management of C++, the ownership-based compile-time model of Rust, and the automatic garbage collection of Python. These languages are evaluated by reproducing a suite of Common Weakness Enumeration (CWE) vulnerabilities, including buffer overflows, use-after-free, and null pointer dereference conditions to determine the stage at which faults are detected: compile-time, runtime, or silent execution.}\end{saidquote}
\afield{Metadata}{[verified] DiVA \texttt{diva2:2091281}, full text \texttt{FULLTEXT01.pdf} retrieved.}
\afield{Verdict}{OK. \enquote{Stage at which faults are detected: compile-time, runtime, or silent execution} is the same three-outcome partition this thesis uses for RQ2, in the source's own words.}
\afield{Record}{Sources\_evidence.md, entry record.}
\auditentry{90}{okawa2024ruststudy}{peer reviewed, level~B}
\afield{Reference}{\fullcite{okawa2024ruststudy}}
\afield{Entry}{DOI \href{https://doi.org/10.1109/CANDARW64572.2024.00066}{\texttt{10.1109/CANDARW64572.2024.00066}}.}
\afield{Cited in}{1 place.}
\begin{citedsentence}\textbf{\small Section 7.5.}\enspace Idiomatic safe Rust implementations of the STAMP transactional benchmarks run slower than both the C originals and machine-transpiled Rust~\citemark{86}, and a general study of Rust against C reports it significantly slower in some cases~\citemark{90}.\end{citedsentence}
\afield{Source says}{}
\begin{saidquote}Okawa \& Yamaguchi, CANDARW 2024 (pp. 355--359), finds Rust \enquote{significantly lower} in some cases and proposes mitigation. Minor, small venue, use only if a second general contrary result is wanted next to \texttt{emmerich2019drivers}.\end{saidquote}
\afield{Record}{LITERATURE\_NOTES.md, peer-reviewed additions of the third sweep.}
\auditentry{91}{pirelli2023safelowlevel}{peer reviewed, level~B}
\afield{Reference}{\fullcite{pirelli2023safelowlevel}}
\afield{Entry}{DOI \href{https://doi.org/10.1109/ICSE48619.2023.00183}{\texttt{10.1109/ICSE48619.2023.00183}}.}
\afield{Cited in}{2 places.}
\begin{citedsentence}\textbf{\small Section 7.5.}\enspace Safe low-level code without run-time overhead has been demonstrated for network device drivers, the same class of program as Emmerich's, when the safety obligations are discharged through type invariants rather than checks~\citemark{91}.\end{citedsentence}
\begin{citedsentence}\textbf{\small Section 7.5.}\enspace None is unoccupied on its own, as the matched implementations~\citemark{73, 84, 85, 86, 87, 91, 57}, the percentile sweep stopping at a latency requirement~\citemark{56} and the effort measurements~\citemark{53, 80} all show.\end{citedsentence}
\afield{Source says}{}
\begin{saidquote}Pirelli \& Candea, ICSE 2023 (pp. 2173--2184). Argues safety without run-time overhead is achievable through type invariants, and shows \textbf{safe network device drivers with no run-time overhead} in extended Rust and C\#. The peer-reviewed counterweight to \texttt{emmerich2019drivers}: the thesis currently presents the \enquote{not literally free} side of this question only.\end{saidquote}
\afield{Record}{LITERATURE\_NOTES.md, peer-reviewed additions of the third sweep.}
\auditentry{92}{qin2024safety}{peer reviewed, level~B}
\afield{Reference}{\fullcite{qin2024safety}}
\afield{Entry}{DOI \href{https://doi.org/10.1109/TSE.2024.3380393}{\texttt{10.1109/TSE.2024.3380393}}.}
\afield{Cited in}{2 places.}
\begin{citedsentence}\textbf{\small Section 7.5.}\enspace Empirical studies of real-world Rust find that memory- and thread-safety bugs still occur in practice, with blocking bugs and misused synchronisation among them~\citemark{27, 92}.\end{citedsentence}
\begin{citedsentence}\textbf{\small Appendix B.}\enspace The safety strand is settled the same partial way, with machine-checked foundations for the safe fragment and the aliasing discipline \texttt{unsafe} code has to respect~\citemark{11, 40, 41}, confirmation that \texttt{unsafe} is used sparingly and that memory-safety CVEs always involve it~\citemark{27, 42, 43}, auditing that finds defects at scale even so and unsafety spreading through dependency chains~\citemark{99, 100, 92, 44}, and the reports motivating the shift set against a long record of memory corruption resisting runtime mitigation~\citemark{8, 9, 17, 19, 20, 34}.\end{citedsentence}
\afield{Claim it carries}{memory- and thread-safety bugs still occur in practice, with blocking bugs and misused synchronisation among them}
\afield{Source says}{}
\begin{saidquote}(deposited abstract): \enquote{a comprehensive, empirical study of Rust safety issues by close, manual inspection of \textbf{70 memory bugs, 100 concurrency bugs, and 110 programming errors leading to unexpected execution panics} from five open-source Rust projects, five widely-used Rust libraries, and two online security databases.}\end{saidquote}
\afield{Verdict}{OK.}
\afield{Record}{Sources\_evidence.md, entry record.}
\auditentry{93}{pereira2017energy}{peer reviewed, level~A}
\afield{Reference}{\fullcite{pereira2017energy}}
\afield{Entry}{DOI \href{https://doi.org/10.1145/3136014.3136031}{\texttt{10.1145/3136014.3136031}}.}
\afield{Cited in}{1 place.}
\begin{citedsentence}\textbf{\small Appendix B.}\enspace Compiled systems languages sit close on time and energy~\citemark{45, 93}, Rust's safety is cheap at run time~\citemark{12, 94}, and that holds where timing is tightest, in embedded hypervisors and in packet processing within a few percent of a C reference~\citemark{95, 73, 96}.\end{citedsentence}
\afield{Claim it carries}{compiled systems languages sit close together on time and energy, from general-purpose benchmarks reporting average-case behaviour}
\afield{Source says}{}
\begin{saidquote}(Table 4, \enquote{Normalized global results for Energy, Time, and Memory}, p. 263): C is the 1.00 baseline on both axes, \textbf{Rust is 1.03 energy and 1.04 time}, Go 1.05, C++ 1.34 energy and 1.56 time. The study covers \enquote{twenty seven well-known software languages} over \enquote{ten different programming problems}, which is the general-purpose average-case design the thesis attributes to it.\end{saidquote}
\afield{Metadata}{[verified] Crossref. Green open access, and Table 4 is reproduced in several independent places.}
\afield{Verdict}{OK. See the scope note above: the table supports the thesis's point more strongly than the sentence claims, and the thesis is right not to quote the figures.}
\afield{Record}{Sources\_evidence.md, entry record.}
\auditentry{94}{levy2017kernel}{peer reviewed, level~B}
\afield{Reference}{\fullcite{levy2017kernel}}
\afield{Entry}{DOI \href{https://doi.org/10.1145/3124680.3124717}{\texttt{10.1145/3124680.3124717}}.}
\afield{Cited in}{1 place.}
\begin{citedsentence}\textbf{\small Appendix B.}\enspace Compiled systems languages sit close on time and energy~\citemark{45, 93}, Rust's safety is cheap at run time~\citemark{12, 94}, and that holds where timing is tightest, in embedded hypervisors and in packet processing within a few percent of a C reference~\citemark{95, 73, 96}.\end{citedsentence}
\afield{Claim it carries}{inside a list, that Rust's safety is cheap at run time in kernel settings}
\afield{Source says}{}
\begin{saidquote}(deposited abstract): \enquote{Previous efforts towards this end concluded that writing a kernel requires changing Rust. This paper reaches a different conclusion, that \textbf{no changes to Rust are needed and a kernel can be implemented with a very small amount of unsafe code}.}\end{saidquote}
\afield{Verdict}{OK for feasibility with little \texttt{unsafe}, which is what the list uses it for. The run-time cost half of that list sentence is carried by \texttt{balasubramanian2017rust} beside it, and this entry should not be quoted alone for a cost claim.}
\afield{Record}{Sources\_evidence.md, entry record.}
\auditentry{95}{mo2023rustshyper}{peer reviewed, level~C}
\afield{Reference}{\fullcite{mo2023rustshyper}}
\afield{Entry}{DOI \href{https://doi.org/10.1016/j.sysarc.2023.102948}{\texttt{10.1016/j.sysarc.2023.102948}}.}
\afield{Cited in}{1 place.}
\begin{citedsentence}\textbf{\small Appendix B.}\enspace Compiled systems languages sit close on time and energy~\citemark{45, 93}, Rust's safety is cheap at run time~\citemark{12, 94}, and that holds where timing is tightest, in embedded hypervisors and in packet processing within a few percent of a C reference~\citemark{95, 73, 96}.\end{citedsentence}
\afield{Source says}{}
\begin{saidquote}Elsevier paywall, title-level support for \enquote{embedded hypervisor in Rust}.\end{saidquote}
\afield{Metadata}{[verified] Crossref matches, including the article number as page.}
\afield{Verdict}{OK.}
\afield{Record}{Sources\_evidence.md, entry record.}
\auditentry{96}{xiao2024rpkt}{peer reviewed, level~C}
\afield{Reference}{\fullcite{xiao2024rpkt}}
\afield{Entry}{DOI \href{https://doi.org/10.1007/978-981-96-2864-3_36}{\texttt{10.1007/978-981-96-2864-3\_36}}.}
\afield{Cited in}{1 place.}
\begin{citedsentence}\textbf{\small Appendix B.}\enspace Compiled systems languages sit close on time and energy~\citemark{45, 93}, Rust's safety is cheap at run time~\citemark{12, 94}, and that holds where timing is tightest, in embedded hypervisors and in packet processing within a few percent of a C reference~\citemark{95, 73, 96}.\end{citedsentence}
\afield{Claim it carries}{packet processing within a few percent of a C reference}
\afield{Source says}{}
\begin{saidquote}Springer paywall, title-level support (\enquote{Generic, Safe, and Efficient Userspace Packet Processing Library in Rust}). The \enquote{few percent} half of the sentence is carried by \texttt{emmerich2019drivers}, which is verified above and cited alongside.\end{saidquote}
\afield{Metadata}{[verified] Pages \textbf{454--465} confirmed by DBLP and Crossref. Note the split year: DBLP lists the conference as \textbf{NPC 2024}, Springer published the LNCS proceedings in \textbf{2025}, which is what the entry's \texttt{year} field says. Optional addition: \texttt{series = \{Lecture Notes in Computer Science\}}.}
\afield{Verdict}{OK.}
\afield{Record}{Sources\_evidence.md, entry record.}
\auditentry{97}{cui2025perfbugs}{peer reviewed, level~B}
\afield{Reference}{\fullcite{cui2025perfbugs}}
\afield{Entry}{DOI \href{https://doi.org/10.1109/ISSRE66568.2025.00045}{\texttt{10.1109/ISSRE66568.2025.00045}}.}
\afield{Cited in}{1 place.}
\begin{citedsentence}\textbf{\small Appendix B.}\enspace Where a cost does appear it has been pinned down: to bounds, overflow and division checks, to that same check set against hardware-enforced capabilities, and to recurring anti-patterns~\citemark{74, 75, 97}.\end{citedsentence}
\afield{Claim it carries}{recurring performance anti-patterns}
\afield{Source says}{}
\begin{saidquote}(deposited abstract): \enquote{we have mined the repository of 100 well-known projects on GitHub and collected 114 performance bugs. We manually audit each case and find \textbf{8 optimization patterns unique to Rust}, including three types of checking, two types of cloning, and three types of data collection.}\end{saidquote}
\afield{Verdict}{OK. \enquote{Recurring anti-patterns} is the abstract's eight patterns.}
\afield{Record}{Sources\_evidence.md, entry record.}
\auditentry{98}{naylor2026cherisimt}{peer reviewed, level~A}
\afield{Reference}{\fullcite{naylor2026cherisimt}}
\afield{Entry}{DOI \href{https://doi.org/10.1145/3760250.3762234}{\texttt{10.1145/3760250.3762234}}.}
\afield{Cited in}{1 place.}
\begin{citedsentence}\textbf{\small Appendix B.}\enspace One like-for-like port puts the bounds-check share at 34\% of a 46\% geometric-mean overhead, and cautions that code credited with cheap checks has sometimes declared its hot functions \texttt{unsafe}~\citemark{98}.\end{citedsentence}
\afield{Source says}{}
\begin{saidquote}\enquote{CHERI-SIMT: Implementing Capability Memory Protection in GPUs}. Matthew Naylor, Alexandre Joannou, A. Theodore Markettos, Paul Metzger, Simon W. Moore, Timothy M. Jones. ASPLOS '26, pages 50--66. DOI 10.1145/3760250.3762234. \textbf{PARTIAL on the year: Crossref's deposit records 2025 while the event field records ASPLOS '26. Confirm the citable year before use.}\end{saidquote}
\begin{saidquote}The GPU subject looks off-topic and the relevant part is not the subject. The authors build a \textbf{like-for-like C++-to-Rust port} of a benchmark suite and report \enquote{a geometric mean overhead of 46\% from using Rust}, of which \enquote{bounds checking alone accounts for a 34\% overhead}.\end{saidquote}
\begin{saidquote}\textbf{The reason to add it is the paragraph after that number.} It engages \texttt{harris2023morello}, which this thesis already cites, and disputes how it should be read: prior work \enquote{reports that the Rust compiler is very good at eliminating bounds checks, resulting in low overhead ... However, care is needed when interpreting these results as it is unclear how many bounds checks are avoided due to the use of unsafe code blocks. For example, in the matrixmultiply benchmark they use, the key functions are declared as unsafe and operate directly on raw pointers.}\end{saidquote}
\afield{Record}{LITERATURE\_NOTES.md, section F.0b Naylor et al., ASPLOS 2026.}
\auditentry{99}{bae2021rudra}{peer reviewed, level~B}
\afield{Reference}{\fullcite{bae2021rudra}}
\afield{Entry}{DOI \href{https://doi.org/10.1145/3477132.3483570}{\texttt{10.1145/3477132.3483570}}.}
\afield{Cited in}{1 place.}
\begin{citedsentence}\textbf{\small Appendix B.}\enspace The safety strand is settled the same partial way, with machine-checked foundations for the safe fragment and the aliasing discipline \texttt{unsafe} code has to respect~\citemark{11, 40, 41}, confirmation that \texttt{unsafe} is used sparingly and that memory-safety CVEs always involve it~\citemark{27, 42, 43}, auditing that finds defects at scale even so and unsafety spreading through dependency chains~\citemark{99, 100, 92, 44}, and the reports motivating the shift set against a long record of memory corruption resisting runtime mitigation~\citemark{8, 9, 17, 19, 20, 34}.\end{citedsentence}
\afield{Claim it carries}{ecosystem-scale auditing finds defects nonetheless}
\afield{Source says}{}
\begin{saidquote}(deposited abstract): \enquote{Rust's safety guarantee relies on the soundness of \textbf{all unsafe code in the program as well as the standard and external libraries}, making it hard to reason about their correctness. In other words, \textbf{a single bug in any unsafe code breaks the whole program's safety guarantee.}}\end{saidquote}
\afield{Verdict}{OK for the premise. The ecosystem-scale audit result itself, the package-set scan and its defect count, sits in the body rather than the deposited abstract, so this is the one entry in this group whose claim is only partly carried by what was read. It is used inside a list beside \texttt{zheng2024security}, \texttt{qin2024safety} and \texttt{evans2020rust}, all of which carry it. A gold open-access PDF exists at ACM for anyone wanting to close it properly.}
\afield{Record}{Sources\_evidence.md, entry record.}
\auditentry{100}{zheng2024security}{peer reviewed, level~A}
\afield{Reference}{\fullcite{zheng2024security}}
\afield{Entry}{DOI \href{https://doi.org/10.1145/3624738}{\texttt{10.1145/3624738}}.}
\afield{Cited in}{1 place.}
\begin{citedsentence}\textbf{\small Appendix B.}\enspace The safety strand is settled the same partial way, with machine-checked foundations for the safe fragment and the aliasing discipline \texttt{unsafe} code has to respect~\citemark{11, 40, 41}, confirmation that \texttt{unsafe} is used sparingly and that memory-safety CVEs always involve it~\citemark{27, 42, 43}, auditing that finds defects at scale even so and unsafety spreading through dependency chains~\citemark{99, 100, 92, 44}, and the reports motivating the shift set against a long record of memory corruption resisting runtime mitigation~\citemark{8, 9, 17, 19, 20, 34}.\end{citedsentence}
\afield{Source says}{}
\begin{saidquote}(arXiv 2308.15046, p. 1): \enquote{we first compile a dataset of \textbf{433 vulnerabilities, 300 vulnerable code repositories, and 218 vulnerability fix commits} in the Rust ecosystem, spanning over 7 years.}\end{saidquote}
\afield{Metadata}{[verified] TOSEM 33(2), issue dated February 2024 $\rightarrow$ \texttt{year = 2024} right (Crossref online-first 2023). Pages \texttt{34:1--34:30} correct.}
\afield{Verdict}{OK.}
\afield{Record}{Sources\_evidence.md, entry record.}
\auditentry{101}{schuermann2023encapsulated}{peer reviewed, level~B}
\afield{Reference}{\fullcite{schuermann2023encapsulated}}
\afield{Entry}{DOI \href{https://doi.org/10.1145/3625275.3625397}{\texttt{10.1145/3625275.3625397}}.}
\afield{Cited in}{1 place.}
\begin{citedsentence}\textbf{\small Appendix B.}\enspace The seam carries its own defect classes and its own tooling bugs, lighter crossing models exist for embedded use, aliasing patterns expressible in C but not in safe Rust bound translation in principle, interoperation is the central open problem in embedded Rust, \texttt{unsafe} concentrates at exactly these boundaries, and the two largest documented adoptions into existing C and C++ report the resulting friction empirically~\citemark{50, 51, 101, 102, 103, 104, 82, 83}.\end{citedsentence}
\afield{Claim it carries}{lighter crossing models proposed for embedded use}
\afield{Source says}{}
\begin{saidquote}(deposited abstract): \enquote{practical systems must interact with code written in memory-unsafe languages. This is especially true in \textbf{security and safety-critical embedded systems}, where subsystems such as cryptographic implementations are subject to industrial and governmental certification requirements ... \textbf{Any single bug in either the foreign code or the cross-language interactions may arbitrarily violate the memory safety of the wrapping language.}}\end{saidquote}
\afield{Verdict}{OK.}
\afield{Record}{Sources\_evidence.md, entry record.}
\auditentry{102}{emre2023aliasing}{peer reviewed, level~B}
\afield{Reference}{\fullcite{emre2023aliasing}}
\afield{Entry}{DOI \href{https://doi.org/10.1145/3586046}{\texttt{10.1145/3586046}}.}
\afield{Cited in}{1 place.}
\begin{citedsentence}\textbf{\small Appendix B.}\enspace The seam carries its own defect classes and its own tooling bugs, lighter crossing models exist for embedded use, aliasing patterns expressible in C but not in safe Rust bound translation in principle, interoperation is the central open problem in embedded Rust, \texttt{unsafe} concentrates at exactly these boundaries, and the two largest documented adoptions into existing C and C++ report the resulting friction empirically~\citemark{50, 51, 101, 102, 103, 104, 82, 83}.\end{citedsentence}
\afield{Claim it carries}{aliasing patterns expressible in C but not in safe Rust bound translation in principle}
\afield{Source says}{}
\begin{saidquote}Crossref-deposited abstract: \enquote{existing translation methods are \textbf{severely limited by a lack of precision in the Rust compiler's safety checker, causing many safe pointer manipulations to be labeled as potentially unsafe} \ldots{} increasing the number of pointers that can be translated to safe Rust references by 75\% over the baseline (from 12\% to 21\% of all pointers).}\end{saidquote}
\afield{Metadata}{[verified] Crossref matches every field.}
\afield{Verdict}{OK.}
\afield{Record}{Sources\_evidence.md, entry record.}
\auditentry{103}{sharma2024embedded}{peer reviewed, level~A}
\afield{Reference}{\fullcite{sharma2024embedded}}
\afield{Entry}{DOI \href{https://doi.org/10.1145/3658644.3690275}{\texttt{10.1145/3658644.3690275}}.}
\afield{Cited in}{1 place.}
\begin{citedsentence}\textbf{\small Appendix B.}\enspace The seam carries its own defect classes and its own tooling bugs, lighter crossing models exist for embedded use, aliasing patterns expressible in C but not in safe Rust bound translation in principle, interoperation is the central open problem in embedded Rust, \texttt{unsafe} concentrates at exactly these boundaries, and the two largest documented adoptions into existing C and C++ report the resulting friction empirically~\citemark{50, 51, 101, 102, 103, 104, 82, 83}.\end{citedsentence}
\afield{Claim it carries}{interoperation is the central open problem in embedded Rust}
\afield{Source says}{}
\begin{saidquote}(arXiv 2311.05063, p. 1, abstract): \enquote{\textbf{Rust interoperability enables developing Rust applications on top of existing C codebases} \ldots{} we found that existing Rust software support is inadequate, SAST tools cannot handle certain features \ldots{} and \textbf{the prevalence of advanced types in existing Rust software makes it challenging to engineer interoperable code}.}\end{saidquote}
\afield{Metadata}{[verified] Crossref matches (booktitle wording differs cosmetically).}
\afield{Verdict}{OK.}
\afield{Record}{Sources\_evidence.md, entry record.}
\auditentry{104}{mccormack2026unsafe}{peer reviewed, level~A}
\afield{Reference}{\fullcite{mccormack2026unsafe}}
\afield{Entry}{DOI \href{https://doi.org/10.1145/3798277}{\texttt{10.1145/3798277}}.}
\afield{Cited in}{1 place.}
\begin{citedsentence}\textbf{\small Appendix B.}\enspace The seam carries its own defect classes and its own tooling bugs, lighter crossing models exist for embedded use, aliasing patterns expressible in C but not in safe Rust bound translation in principle, interoperation is the central open problem in embedded Rust, \texttt{unsafe} concentrates at exactly these boundaries, and the two largest documented adoptions into existing C and C++ report the resulting friction empirically~\citemark{50, 51, 101, 102, 103, 104, 82, 83}.\end{citedsentence}
\afield{Claim it carries}{\texttt{unsafe} concentrates at exactly these boundaries}
\afield{Source says}{}
\begin{saidquote}(arXiv 2404.02230, p. 1, abstract): \enquote{we study \textbf{how developers reason about foreign function calls}, the limitations of the tools that they currently use \ldots{} \textbf{limited tooling support for foreign function calls made participants uncertain about their design choices, and certain foreign aliasing and concurrency patterns were difficult to encapsulate}.}\end{saidquote}
\afield{Metadata}{[verified] Crossref matches, no volume/issue yet, correctly omitted.}
\afield{Verdict}{OK.}
\afield{Record}{Sources\_evidence.md, entry record.}
\auditentry{105}{hong2025ctorust}{peer reviewed, level~B}
\afield{Reference}{\fullcite{hong2025ctorust}}
\afield{Entry}{DOI \href{https://doi.org/10.1145/3737696}{\texttt{10.1145/3737696}}.}
\afield{Cited in}{1 place.}
\begin{citedsentence}\textbf{\small Appendix B.}\enspace Automating the rewrite is not a solved problem either, since current translators still emit code whose safety the compiler cannot validate~\citemark{105}.\end{citedsentence}
\afield{Source says}{}
\begin{saidquote}Hong \& Ryu, CACM 68(11):58--65, 2025. \enquote{Automatically Translating C to Rust}. Peer-reviewed home for the RQ4 translation story, in place of leaning on preprints (\S{}6.2).\end{saidquote}
\afield{Record}{LITERATURE\_NOTES.md, peer-reviewed additions of the third sweep.}
\auditentry{106}{huang2017predictability}{peer reviewed, level~B}
\afield{Reference}{\fullcite{huang2017predictability}}
\afield{Entry}{DOI \href{https://doi.org/10.1145/3035918.3064016}{\texttt{10.1145/3035918.3064016}}.}
\afield{Cited in}{1 place.}
\begin{citedsentence}\textbf{\small Appendix B.}\enspace The fourth strand is the domain and the method: the tail-latency and microsecond-scale literature~\citemark{6, 7, 31, 106}, thread-per-core structuring moving the tail in its own right~\citemark{46}, O-RAN disaggregation~\citemark{3}, control-plane latency as an active target in the core network too~\citemark{107, 108}, adoption evidence pointing at the learning curve rather than run-time cost~\citemark{80, 81, 109}, and standard guidance on repeat-based evaluation and fault injection~\citemark{33, 60}.\end{citedsentence}
\afield{Claim it carries}{part of the microsecond-scale predictability literature}
\afield{Source says}{}
\begin{saidquote}(deposited abstract): \enquote{While much of the research on transaction processing has focused on improving overall performance in terms of throughput and mean latency, surprisingly less attention has been given to \textbf{performance predictability: how often individual transactions exhibit execution latency far from the mean}.}\end{saidquote}
\afield{Verdict}{OK. Latency far from the mean rather than the mean itself is exactly the distinction the thesis cites it for.}
\afield{Record}{Sources\_evidence.md, entry record.}
\auditentry{107}{liu2023l25gc}{peer reviewed, level~B}
\afield{Reference}{\fullcite{liu2023l25gc}}
\afield{Entry}{DOI \href{https://doi.org/10.1109/CloudNet59005.2023.10490024}{\texttt{10.1109/CloudNet59005.2023.10490024}}.}
\afield{Cited in}{1 place.}
\begin{citedsentence}\textbf{\small Appendix B.}\enspace The fourth strand is the domain and the method: the tail-latency and microsecond-scale literature~\citemark{6, 7, 31, 106}, thread-per-core structuring moving the tail in its own right~\citemark{46}, O-RAN disaggregation~\citemark{3}, control-plane latency as an active target in the core network too~\citemark{107, 108}, adoption evidence pointing at the learning curve rather than run-time cost~\citemark{80, 81, 109}, and standard guidance on repeat-based evaluation and fault injection~\citemark{33, 60}.\end{citedsentence}
\afield{Claim it carries}{control-plane latency is an active target in the core network}
\afield{Source says}{}
\begin{saidquote}(deposited abstract): \enquote{the \textbf{control plane in 5G core (5GC) still presents challenges} due to inefficiencies in handling control plane operations (including session establishment, handovers and idle-to-active state-transitions) ... introduces substantial overheads that \textbf{impact latency}.}\end{saidquote}
\afield{Verdict}{OK.}
\afield{Record}{Sources\_evidence.md, entry record.}
\auditentry{108}{park2025ctrlbench}{peer reviewed, level~B}
\afield{Reference}{\fullcite{park2025ctrlbench}}
\afield{Entry}{DOI \href{https://doi.org/10.1109/ICTC66702.2025.11388531}{\texttt{10.1109/ICTC66702.2025.11388531}}.}
\afield{Cited in}{1 place.}
\begin{citedsentence}\textbf{\small Appendix B.}\enspace The fourth strand is the domain and the method: the tail-latency and microsecond-scale literature~\citemark{6, 7, 31, 106}, thread-per-core structuring moving the tail in its own right~\citemark{46}, O-RAN disaggregation~\citemark{3}, control-plane latency as an active target in the core network too~\citemark{107, 108}, adoption evidence pointing at the learning curve rather than run-time cost~\citemark{80, 81, 109}, and standard guidance on repeat-based evaluation and fault injection~\citemark{33, 60}.\end{citedsentence}
\afield{Claim it carries}{control-plane latency is an active target, measured procedure by procedure}
\afield{Source says}{}
\begin{saidquote}(deposited abstract): \enquote{CtrlBench, a domain-specific \textbf{microbenchmarking tool for the 5G Core Control Plane} ... It can generate configurable load and chain-dependent APIs, giving granular, fine-grained, and reproducible latency measurements ... \textbf{Existing end-to-end metrics, such as Registration Time, compress entire procedures into a single number. This hides which Network Function or API actually causes delays.}}\end{saidquote}
\afield{Verdict}{OK, and the closing sentence is an independent statement of this thesis's own decomposition argument, that an end-to-end figure hides where the cost sits.}
\afield{Record}{Sources\_evidence.md, entry record.}
\auditentry{109}{ardito2021verbosity}{peer reviewed, level~B}
\afield{Reference}{\fullcite{ardito2021verbosity}}
\afield{Entry}{DOI \href{https://doi.org/10.7717/peerj-cs.406}{\texttt{10.7717/peerj-cs.406}}.}
\afield{Cited in}{1 place.}
\begin{citedsentence}\textbf{\small Appendix B.}\enspace The fourth strand is the domain and the method: the tail-latency and microsecond-scale literature~\citemark{6, 7, 31, 106}, thread-per-core structuring moving the tail in its own right~\citemark{46}, O-RAN disaggregation~\citemark{3}, control-plane latency as an active target in the core network too~\citemark{107, 108}, adoption evidence pointing at the learning curve rather than run-time cost~\citemark{80, 81, 109}, and standard guidance on repeat-based evaluation and fault injection~\citemark{33, 60}.\end{citedsentence}
\afield{Claim it carries}{adoption evidence points at the learning curve rather than run-time cost}
\afield{Source says}{}
\begin{saidquote}Crossref-deposited abstract: \enquote{The Rust code had an \textbf{average verbosity} in terms of the raw size of the code \ldots{} \textbf{Rust code had a better Cyclomatic Complexity, Halstead Metrics, and Maintainability Indexes than C and C++} but performed worse than the other considered object-oriented languages.}\end{saidquote}
\afield{Metadata}{[verified] Crossref matches every field.}
\afield{Verdict}{OK.}
\afield{Record}{Sources\_evidence.md, entry record.}
\auditentry{110}{qazi2017pepc}{peer reviewed, level~A}
\afield{Reference}{\fullcite{qazi2017pepc}}
\afield{Entry}{DOI \href{https://doi.org/10.1145/3098822.3098848}{\texttt{10.1145/3098822.3098848}}.}
\afield{Cited in}{1 place.}
\begin{citedsentence}\textbf{\small Appendix B.}\enspace Rust has mobile-network precedent in the datapath~\citemark{110} and, on the control-plane side, an open-source gNB-CU with a separated CU-CP and CU-UP over E1.\end{citedsentence}
\afield{Claim it carries}{Rust has mobile-network precedent in the datapath}
\afield{Source says}{}
\begin{saidquote}(author copy, \S{}4.1 Implementation): \enquote{NetBricks uses the Rust programming language [38], and \textbf{PEPC is also written in rust}, and runs inside NetBricks execution environment.}\end{saidquote}
\afield{Metadata}{[verified] Crossref, pages 348-361. Author copy at \texttt{web.lums.edu.pk/\textasciitilde{}zafar/sigcomm17.pdf}.}
\afield{Verdict}{OK. Worth recording how close this came to being a defect: the Crossref-deposited abstract is truncated and never mentions Rust, so on the abstract alone this citation looks unsupported. The implementation section settles it outright. \textbf{An abstract that omits the implementation language is not evidence that the language is absent.}}
\afield{Record}{Sources\_evidence.md, entry record.}
\auditentry{111}{alsoran}{tool record, level~A}
\afield{Reference}{\fullcite{alsoran}}
\afield{Entry}{\href{https://github.com/nplrkn/alsoran}{source page}. Proof of concept, development discontinued. Accessed August 2026.}
\afield{Cited in}{1 place.}
\begin{citedsentence}\textbf{\small Appendix B.}\enspace That last one is an existence proof and not a comparison, reporting no latency and carrying no C++ counterpart built to the same specification~\citemark{111}.\end{citedsentence}
\afield{Claim it carries}{an open-source proof-of-concept gNB-CU with separated CU-CP and CU-UP over E1, reporting no latency and with no C++ counterpart}
\afield{Source says}{}
\begin{saidquote}(repository README): \enquote{Alsoran is a \textbf{Rust implementation of the gNodeB Centralized Unit (gNB-CU)} of the 5G Radio Access Network}. \enquote{This project is currently a \textbf{proof of concept} and not yet a fully functional gNB-CU}, under Current support: \enquote{\textbf{Separate CU-CP and CU-UP communicating using E1}}, \enquote{Scale-out of GNB-CU-CP workers using multiple TNLAs}, \enquote{ASN.1 libraries for NGAP, E1AP, F1AP and RRC}, and the banner \enquote{This project is no longer active}.\end{saidquote}
\afield{Metadata}{[verified] Repository description: \enquote{5G gNodeB Centralized Unit written in Rust}. MIT licence, last pushed June 2025. The owner's GitHub account (\texttt{nplrkn}) shows the given name \enquote{Nic} only, the surname \enquote{Larkin} in the entry could not be confirmed from the repository itself, check the QCore project page if you want it airtight. Consider adding an access date, as the entry has no \texttt{year}.}
\afield{Verdict}{OK, the thesis's characterisation is accurate, including the absence of latency numbers.}
\afield{Record}{Sources\_evidence.md, entry record.}

\clearpage
\section*{Part 3. Findings and open items}
\addcontentsline{toc}{section}{Part 3. Findings and open items}

\subsection*{Citations that were attached to a claim the source does not make}

All six are closed as of
31 August 2026, and they are kept
here because the pattern recurs. In every case the work itself was sound and cited correctly
elsewhere, so each was closed by moving or dropping the citation rather than by replacing the
source. No source was added that the bibliography did not already hold.

Four of the six existed only in \texttt{Thesis\_v3} and the full draft had them right, which
locates the cause: compressing a chapter into five paragraphs detaches a citation from the clause
it was attached to. The two that reached the full draft are both in Chapter~1, which
\texttt{Thesis\_v3} carries as a copy. So the rule for the next compression pass is to re-read
what each surviving citation now sits next to.

\begin{enumerate}[leftmargin=1.6em,itemsep=4pt]
\item \textbf{Chapter 1, on Rust promising performance competitive with C++.}\par \texttt{klabnik2023rust}, thesis [10]. \texttt{jung2018rustbelt}, thesis [11].\par RustBelt is a machine-checked safety proof and makes no performance claim. The type-system half of the sentence was carried by the book. Nothing in the pair supported the performance half. The same sentence was in the full draft, because this chapter was copied from it, and in the plan in a different wording.\par\textit{Closed.} The pair now carries only the compile-time guarantee, phrased as enforced \enquote{on machine-checked foundations}, and the performance clause carries \texttt{balasubramanian2017rust}, which is the source the full draft already uses for that framing. The plan's sentence drops the performance claim, which its next paragraph states and cites on its own.
\item \textbf{Chapter 1, on thousands of procedures per second at a budget of a few milliseconds.}\par \texttt{3gpp38913}, thesis [5].\par TR 38.913 clause 7.4 sets a control-plane latency target. It states no procedure rate, so the citation covered the second half of a compound sentence only.\par\textit{Closed.} The sentence now reads that each procedure has to complete inside the millisecond-scale budget the 5G system sets for control-plane signalling, which is what clause 7.4 gives. The rate stands as this study's own workload assumption, in the wording already used later in the same chapter.
\item \textbf{Chapter 6, second use, on the seven boundary classes being chosen because safe Rust is expected not to prevent them.}\par \texttt{natella2016fault}, thesis [60].\par Natella's survey of software fault injection is the right source for the method sentence in Chapter 4. It says nothing about the boundary of Rust's guarantee, and Chapter 3 already states the catalogue is this study's own. The full draft never had this one.\par\textit{Closed.} The citation is dropped. The method sentence in Chapter 4 keeps it.
\item \textbf{Chapter 6, on the boundary defects being re-imported against the checker's own primitives and shrunk to model size.}\par \texttt{bornholt2021shardstore}, thesis [71].\par ShardStore is the peer-reviewed account of Loom and Shuttle and supports the division of labour between them, which is how the full draft uses it. The rewrite and the shrink are this study's own measured cost.\par\textit{Closed.} The citation moves onto the division of labour, naming Loom and Shuttle as the full draft does, and the rewrite cost is left as the study's own finding.
\item \textbf{Chapter 7, on confining unsafe to a seam crate making the reviewable surface a countable thing.}\par \texttt{evans2020rust}, thesis [44].\par Evans et al.\ report that fewer than 30 per cent of crates use unsafe directly and that more than half cannot be fully checked because of unsafe in the call chain. All five sites in the full draft cite that for dependency propagation. The countability of the surface is this thesis's own guideline.\par\textit{Closed.} The citation moves onto the dependency-propagation clause, matching the full draft, and the countable surface follows as the guideline it is.
\item \textbf{Chapter 7, on the parity suite bounding behaviour and not performance.}\par \texttt{emberg2023asn1}, thesis [57].\par That limitation is this study's own, and both places the full draft states it carry no citation. Emberg and Wadsten compared matched C++ and Rust ASN.1 decompilers on performance.\par\textit{Closed.} The citation is dropped. The entry's other use, on no earlier study isolating the language for a stateful concurrent subsystem, was sound and stands.
\end{enumerate}

\subsection*{Entries whose supporting passage was never located}

None. This list is closed.

The list was six entries until 31 August 2026, all of them student theses or a paywalled paper, and
three of the defects above sat on three of those keys. Where nobody had written down what the
source says, the prose drifted. Closing it took one afternoon, because five of the six were openly
deposited with full PDFs. The correlation is the reason to keep this list at zero.



\subsection*{Level C, metadata confirmed and passage not opened}

These 6 entries could not be closed without university access, and each is
used for a title-level claim or inside a list. Three have no abstract deposited anywhere, so there
is nothing short of the full text to read. Three have no DOI because the claim is a specification
clause or a book chapter rather than a finding, and in each of those the citation sits beside a
verified one that carries the same point.

\begin{itemize}[leftmargin=1.4em,itemsep=0pt,topsep=2pt]
\item \texttt{3gpp38331}\hfill [4]
\item \texttt{klabnik2023rust}\hfill [10]
\item \texttt{3gpp22261}\hfill [29]
\item \texttt{li2022crosslang}\hfill [50]
\item \texttt{mo2023rustshyper}\hfill [95]
\item \texttt{xiao2024rpkt}\hfill [96]
\end{itemize}

\subsection*{The one preprint still load-bearing}

\texttt{bugden2022rust} is the last arXiv-only entry carrying weight. It is cited twice and both
times beside a peer-reviewed source. Replacing it with a peer-reviewed equivalent is optional.

\subsection*{Differences that look like defects and are not}

Recorded so no later pass \enquote{corrects} them.

\begin{itemize}[leftmargin=1.4em,itemsep=2pt]
\item \textbf{Crossref year against issue year.} Crossref dates the POPL 2018 and POPL 2020 issues a year early because the deposit precedes the issue. The entries follow what ACM prints.
\item \textbf{Article-number page ranges.} \texttt{66:1--66:34} against Crossref's \texttt{1-34}, and eight more like it. ACM prints the article number and the entries follow ACM.
\item \textbf{Author names that differ from the publisher record.} D'Oro, O'Keeffe, Tao Wei, Garc\'ia S\'anchez-Torija and di Luzio all differ between the publisher record and the paper's own title page. The entries follow the papers.
\item \textbf{Semicolons in 3GPP titles.} \texttt{NG-RAN; Architecture description} and four like it keep their semicolons, because the punctuation is part of the official title. Changing them would misquote the standard.
\item \textbf{HTTP 403 and 202 on publisher links.} Bot blocking rather than dead links. Every one of those DOIs resolves through the Crossref API.
\end{itemize}


\clearpage
\section*{Part 4. Coverage}
\addcontentsline{toc}{section}{Part 4. Coverage}

\begin{tabular}{@{}lrr@{}}
\toprule
\textbf{Source kind} & \textbf{n} & \textbf{share} \\
\midrule
peer reviewed & 74 & 67\,\% \\
standard & 11 & 10\,\% \\
grey & 8 & 7\,\% \\
thesis & 7 & 6\,\% \\
tool record & 5 & 5\,\% \\
committee draft & 3 & 3\,\% \\
book & 2 & 2\,\% \\
regulation & 1 & 1\,\% \\
\midrule
\textbf{total} & \textbf{111} & \\
\bottomrule
\end{tabular}

\vspace{1em}

\begin{tabular}{@{}lr@{}}
\toprule
\textbf{Evidence} & \textbf{n} \\
\midrule
Level A, passage located & 73 \\
Level B, deposited record carries it & 32 \\
Level C, metadata only & 6 \\
No located passage & 0 \\
\midrule
\textbf{total} & \textbf{111} \\
\bottomrule
\end{tabular}

\vspace{1em}

Every entry in the bibliography is cited, so there is no orphan and no citation without an
entry. The thesis carries 176 citation commands naming 263 works between
them, since one command often cites several.
74 entries (67\,\%) are in peer-reviewed venues.

\end{document}